% Options for packages loaded elsewhere
% Options for packages loaded elsewhere
\PassOptionsToPackage{unicode}{hyperref}
\PassOptionsToPackage{hyphens}{url}
\PassOptionsToPackage{dvipsnames,svgnames,x11names}{xcolor}
%
\documentclass[
  letterpaper,
  DIV=11,
  numbers=noendperiod]{scrreprt}
\usepackage{xcolor}
\usepackage{amsmath,amssymb}
\setcounter{secnumdepth}{5}
\usepackage{iftex}
\ifPDFTeX
  \usepackage[T1]{fontenc}
  \usepackage[utf8]{inputenc}
  \usepackage{textcomp} % provide euro and other symbols
\else % if luatex or xetex
  \usepackage{unicode-math} % this also loads fontspec
  \defaultfontfeatures{Scale=MatchLowercase}
  \defaultfontfeatures[\rmfamily]{Ligatures=TeX,Scale=1}
\fi
\usepackage{lmodern}
\ifPDFTeX\else
  % xetex/luatex font selection
\fi
% Use upquote if available, for straight quotes in verbatim environments
\IfFileExists{upquote.sty}{\usepackage{upquote}}{}
\IfFileExists{microtype.sty}{% use microtype if available
  \usepackage[]{microtype}
  \UseMicrotypeSet[protrusion]{basicmath} % disable protrusion for tt fonts
}{}
\makeatletter
\@ifundefined{KOMAClassName}{% if non-KOMA class
  \IfFileExists{parskip.sty}{%
    \usepackage{parskip}
  }{% else
    \setlength{\parindent}{0pt}
    \setlength{\parskip}{6pt plus 2pt minus 1pt}}
}{% if KOMA class
  \KOMAoptions{parskip=half}}
\makeatother
% Make \paragraph and \subparagraph free-standing
\makeatletter
\ifx\paragraph\undefined\else
  \let\oldparagraph\paragraph
  \renewcommand{\paragraph}{
    \@ifstar
      \xxxParagraphStar
      \xxxParagraphNoStar
  }
  \newcommand{\xxxParagraphStar}[1]{\oldparagraph*{#1}\mbox{}}
  \newcommand{\xxxParagraphNoStar}[1]{\oldparagraph{#1}\mbox{}}
\fi
\ifx\subparagraph\undefined\else
  \let\oldsubparagraph\subparagraph
  \renewcommand{\subparagraph}{
    \@ifstar
      \xxxSubParagraphStar
      \xxxSubParagraphNoStar
  }
  \newcommand{\xxxSubParagraphStar}[1]{\oldsubparagraph*{#1}\mbox{}}
  \newcommand{\xxxSubParagraphNoStar}[1]{\oldsubparagraph{#1}\mbox{}}
\fi
\makeatother

\usepackage{color}
\usepackage{fancyvrb}
\newcommand{\VerbBar}{|}
\newcommand{\VERB}{\Verb[commandchars=\\\{\}]}
\DefineVerbatimEnvironment{Highlighting}{Verbatim}{commandchars=\\\{\}}
% Add ',fontsize=\small' for more characters per line
\usepackage{framed}
\definecolor{shadecolor}{RGB}{241,243,245}
\newenvironment{Shaded}{\begin{snugshade}}{\end{snugshade}}
\newcommand{\AlertTok}[1]{\textcolor[rgb]{0.68,0.00,0.00}{#1}}
\newcommand{\AnnotationTok}[1]{\textcolor[rgb]{0.37,0.37,0.37}{#1}}
\newcommand{\AttributeTok}[1]{\textcolor[rgb]{0.40,0.45,0.13}{#1}}
\newcommand{\BaseNTok}[1]{\textcolor[rgb]{0.68,0.00,0.00}{#1}}
\newcommand{\BuiltInTok}[1]{\textcolor[rgb]{0.00,0.23,0.31}{#1}}
\newcommand{\CharTok}[1]{\textcolor[rgb]{0.13,0.47,0.30}{#1}}
\newcommand{\CommentTok}[1]{\textcolor[rgb]{0.37,0.37,0.37}{#1}}
\newcommand{\CommentVarTok}[1]{\textcolor[rgb]{0.37,0.37,0.37}{\textit{#1}}}
\newcommand{\ConstantTok}[1]{\textcolor[rgb]{0.56,0.35,0.01}{#1}}
\newcommand{\ControlFlowTok}[1]{\textcolor[rgb]{0.00,0.23,0.31}{\textbf{#1}}}
\newcommand{\DataTypeTok}[1]{\textcolor[rgb]{0.68,0.00,0.00}{#1}}
\newcommand{\DecValTok}[1]{\textcolor[rgb]{0.68,0.00,0.00}{#1}}
\newcommand{\DocumentationTok}[1]{\textcolor[rgb]{0.37,0.37,0.37}{\textit{#1}}}
\newcommand{\ErrorTok}[1]{\textcolor[rgb]{0.68,0.00,0.00}{#1}}
\newcommand{\ExtensionTok}[1]{\textcolor[rgb]{0.00,0.23,0.31}{#1}}
\newcommand{\FloatTok}[1]{\textcolor[rgb]{0.68,0.00,0.00}{#1}}
\newcommand{\FunctionTok}[1]{\textcolor[rgb]{0.28,0.35,0.67}{#1}}
\newcommand{\ImportTok}[1]{\textcolor[rgb]{0.00,0.46,0.62}{#1}}
\newcommand{\InformationTok}[1]{\textcolor[rgb]{0.37,0.37,0.37}{#1}}
\newcommand{\KeywordTok}[1]{\textcolor[rgb]{0.00,0.23,0.31}{\textbf{#1}}}
\newcommand{\NormalTok}[1]{\textcolor[rgb]{0.00,0.23,0.31}{#1}}
\newcommand{\OperatorTok}[1]{\textcolor[rgb]{0.37,0.37,0.37}{#1}}
\newcommand{\OtherTok}[1]{\textcolor[rgb]{0.00,0.23,0.31}{#1}}
\newcommand{\PreprocessorTok}[1]{\textcolor[rgb]{0.68,0.00,0.00}{#1}}
\newcommand{\RegionMarkerTok}[1]{\textcolor[rgb]{0.00,0.23,0.31}{#1}}
\newcommand{\SpecialCharTok}[1]{\textcolor[rgb]{0.37,0.37,0.37}{#1}}
\newcommand{\SpecialStringTok}[1]{\textcolor[rgb]{0.13,0.47,0.30}{#1}}
\newcommand{\StringTok}[1]{\textcolor[rgb]{0.13,0.47,0.30}{#1}}
\newcommand{\VariableTok}[1]{\textcolor[rgb]{0.07,0.07,0.07}{#1}}
\newcommand{\VerbatimStringTok}[1]{\textcolor[rgb]{0.13,0.47,0.30}{#1}}
\newcommand{\WarningTok}[1]{\textcolor[rgb]{0.37,0.37,0.37}{\textit{#1}}}

\usepackage{longtable,booktabs,array}
\newcounter{none} % for unnumbered tables
\usepackage{calc} % for calculating minipage widths
% Correct order of tables after \paragraph or \subparagraph
\usepackage{etoolbox}
\makeatletter
\patchcmd\longtable{\par}{\if@noskipsec\mbox{}\fi\par}{}{}
\makeatother
% Allow footnotes in longtable head/foot
\IfFileExists{footnotehyper.sty}{\usepackage{footnotehyper}}{\usepackage{footnote}}
\makesavenoteenv{longtable}
\usepackage{graphicx}
\makeatletter
\newsavebox\pandoc@box
\newcommand*\pandocbounded[1]{% scales image to fit in text height/width
  \sbox\pandoc@box{#1}%
  \Gscale@div\@tempa{\textheight}{\dimexpr\ht\pandoc@box+\dp\pandoc@box\relax}%
  \Gscale@div\@tempb{\linewidth}{\wd\pandoc@box}%
  \ifdim\@tempb\p@<\@tempa\p@\let\@tempa\@tempb\fi% select the smaller of both
  \ifdim\@tempa\p@<\p@\scalebox{\@tempa}{\usebox\pandoc@box}%
  \else\usebox{\pandoc@box}%
  \fi%
}
% Set default figure placement to htbp
\def\fps@figure{htbp}
\makeatother


% definitions for citeproc citations
\NewDocumentCommand\citeproctext{}{}
\NewDocumentCommand\citeproc{mm}{%
  \begingroup\def\citeproctext{#2}\cite{#1}\endgroup}
\makeatletter
 % allow citations to break across lines
 \let\@cite@ofmt\@firstofone
 % avoid brackets around text for \cite:
 \def\@biblabel#1{}
 \def\@cite#1#2{{#1\if@tempswa , #2\fi}}
\makeatother
\newlength{\cslhangindent}
\setlength{\cslhangindent}{1.5em}
\newlength{\csllabelwidth}
\setlength{\csllabelwidth}{3em}
\newenvironment{CSLReferences}[2] % #1 hanging-indent, #2 entry-spacing
 {\begin{list}{}{%
  \setlength{\itemindent}{0pt}
  \setlength{\leftmargin}{0pt}
  \setlength{\parsep}{0pt}
  % turn on hanging indent if param 1 is 1
  \ifodd #1
   \setlength{\leftmargin}{\cslhangindent}
   \setlength{\itemindent}{-1\cslhangindent}
  \fi
  % set entry spacing
  \setlength{\itemsep}{#2\baselineskip}}}
 {\end{list}}
\usepackage{calc}
\newcommand{\CSLBlock}[1]{\hfill\break\parbox[t]{\linewidth}{\strut\ignorespaces#1\strut}}
\newcommand{\CSLLeftMargin}[1]{\parbox[t]{\csllabelwidth}{\strut#1\strut}}
\newcommand{\CSLRightInline}[1]{\parbox[t]{\linewidth - \csllabelwidth}{\strut#1\strut}}
\newcommand{\CSLIndent}[1]{\hspace{\cslhangindent}#1}



\setlength{\emergencystretch}{3em} % prevent overfull lines

\providecommand{\tightlist}{%
  \setlength{\itemsep}{0pt}\setlength{\parskip}{0pt}}



 


\KOMAoption{captions}{tableheading}
\makeatletter
\@ifpackageloaded{bookmark}{}{\usepackage{bookmark}}
\makeatother
\makeatletter
\@ifpackageloaded{caption}{}{\usepackage{caption}}
\AtBeginDocument{%
\ifdefined\contentsname
  \renewcommand*\contentsname{Table of contents}
\else
  \newcommand\contentsname{Table of contents}
\fi
\ifdefined\listfigurename
  \renewcommand*\listfigurename{List of Figures}
\else
  \newcommand\listfigurename{List of Figures}
\fi
\ifdefined\listtablename
  \renewcommand*\listtablename{List of Tables}
\else
  \newcommand\listtablename{List of Tables}
\fi
\ifdefined\figurename
  \renewcommand*\figurename{Figure}
\else
  \newcommand\figurename{Figure}
\fi
\ifdefined\tablename
  \renewcommand*\tablename{Table}
\else
  \newcommand\tablename{Table}
\fi
}
\@ifpackageloaded{float}{}{\usepackage{float}}
\floatstyle{ruled}
\@ifundefined{c@chapter}{\newfloat{codelisting}{h}{lop}}{\newfloat{codelisting}{h}{lop}[chapter]}
\floatname{codelisting}{Listing}
\newcommand*\listoflistings{\listof{codelisting}{List of Listings}}
\makeatother
\makeatletter
\makeatother
\makeatletter
\@ifpackageloaded{caption}{}{\usepackage{caption}}
\@ifpackageloaded{subcaption}{}{\usepackage{subcaption}}
\makeatother
\usepackage{bookmark}
\IfFileExists{xurl.sty}{\usepackage{xurl}}{} % add URL line breaks if available
\urlstyle{same}
\hypersetup{
  pdftitle={DS\_6021\_Predictive\_Modeling\_Maxwell\_Braseth\_First\_Quarto\_Project},
  pdfauthor={Maxwell Braseth},
  colorlinks=true,
  linkcolor={blue},
  filecolor={Maroon},
  citecolor={Blue},
  urlcolor={Blue},
  pdfcreator={LaTeX via pandoc}}


\title{DS\_6021\_Predictive\_Modeling\_Maxwell\_Braseth\_First\_Quarto\_Project}
\author{Maxwell Braseth}
\date{2026-05-26}
\begin{document}
\maketitle

\renewcommand*\contentsname{Table of contents}
{
\setcounter{tocdepth}{2}
\tableofcontents
}

\bookmarksetup{startatroot}

\chapter*{Preface}\label{preface}
\addcontentsline{toc}{chapter}{Preface}

\markboth{Preface}{Preface}

This is a Quarto book.

To learn more about Quarto books visit
\url{https://quarto.org/docs/books}.

\bookmarksetup{startatroot}

\chapter{Introduction}\label{introduction}

This is a book created from markdown and executable code.

See Knuth (1984) for additional discussion of literate programming.

\bookmarksetup{startatroot}

\chapter{Homework 1: Data Exploration and
Visualization}\label{homework-1-data-exploration-and-visualization}

This Chapter will focus on the topics surrounding Visualizing Data
Models

Topics includes:

\begin{itemize}
\tightlist
\item
  Partial effect plots
\item
  Partial dependent plots
\item
  Individual conditional expectation plots
\item
  Accumulated local effect plots
\item
  (Possible) Permutation feature importance
\end{itemize}

\bookmarksetup{startatroot}

\chapter{Module 1 homework}\label{module-1-homework}

\begin{enumerate}
\def\labelenumi{\arabic{enumi}.}
\item
  Create your course portfolio using the instructions.
\item
  In the HW1.qmd, create another partial effect plot of predictions with
  x1 ranging from -3 to 3 and x2 equal -0.5 and x3 set to 0. Repeat for
  all five models.
\end{enumerate}

\section{Partial Effect Plots}\label{partial-effect-plots}

\bookmarksetup{startatroot}

\chapter{```\{r\}}\label{r}

install.packages(``pak'') pak::pak(``thomasgstewart/ds6021b'')
require(ds6021b)

expanded\_grid\_for\_df\_mod\_01 \textless- expand.grid(x1 = seq(-3, 3,
0.5), x2 = -0.5, x3 = 0)

\bookmarksetup{startatroot}

\chapter{First prediction model}\label{first-prediction-model}

plot(expanded\_grid\_for\_df\_mod\_01\(x1, f1(expanded_grid_for_df_mod_01, x1 = expanded_grid_for_df_mod_01\)x1,
x2 = -0.5, x3 = 0), title = ``First Prediction Model'', xlab = ``x1'',
ylab = ``Black Box Model 1=F1 for
f1(expanded\_grid\_for\_df\_mod\_01)'', type = ``p'')

\bookmarksetup{startatroot}

\chapter{Second prediction model}\label{second-prediction-model}

plot(expanded\_grid\_for\_df\_mod\_01\(x1, f2(expanded_grid_for_df_mod_01, x1 = expanded_grid_for_df_mod_01\)x1,
x2 = -0.5, x3 = 0), title = ``Second Prediction Model'', xlab = ``x1'',
ylab = ``Black Box Model 2=F2 for
f2(expanded\_grid\_for\_df\_mod\_01)'', type = ``b'')

\bookmarksetup{startatroot}

\chapter{Third prediction model}\label{third-prediction-model}

plot(expanded\_grid\_for\_df\_mod\_01\(x1, f3(expanded_grid_for_df_mod_01, x1 = expanded_grid_for_df_mod_01\)x1,
x2 = -0.5, x3 = 0), title = ``Third Prediction Model'', xlab = ``x1'',
ylab = ``Black Box Model 3=F3 for
f3(expanded\_grid\_for\_df\_mod\_01)'', type = ``b'')

\bookmarksetup{startatroot}

\chapter{Fourth prediction model}\label{fourth-prediction-model}

plot(expanded\_grid\_for\_df\_mod\_01\(x1, f4(expanded_grid_for_df_mod_01, x1 = expanded_grid_for_df_mod_01\)x1,
x2 = -0.5, x3 = 0), title = ``Fourth Prediction Model'', xlab = ``x1'',
ylab = ``Black Box Model 4=F4 for
f4(expanded\_grid\_for\_df\_mod\_01)'', type = ``o'')

\bookmarksetup{startatroot}

\chapter{Fifth prediction model}\label{fifth-prediction-model}

plot(expanded\_grid\_for\_df\_mod\_01\(x1, f5(expanded_grid_for_df_mod_01, x1 = expanded_grid_for_df_mod_01\)x1,
x2 = -0.5, x3 = 0), title = ``Fifth Prediction Model'', xlab = ``x1'',
ylab = ``Black Box Model 5=F5 for
f5(expanded\_grid\_for\_df\_mod\_01)'', type = ``s'')```

\begin{enumerate}
\def\labelenumi{\arabic{enumi}.}
\setcounter{enumi}{2}
\tightlist
\item
  Create another plot that compares predictions for the following using
  all five models:
\end{enumerate}

\begin{itemize}
\tightlist
\item
  x1 ranging from -3 to 3 and x2 equal -0.5 and x3 set to 0
\item
  x1 ranging from -3 to 3 and x2 equal -0.5 and x3 set to 1
\end{itemize}

\# Generate predictions for ALL models

\bookmarksetup{startatroot}

\chapter{x3 = 0}\label{x3-0}

y1\_0 \textless- f1(x1 = seq(-3, 3, 0.5), x2 = -0.5, x3 = 0) y2\_0
\textless- f2(x1 = seq(-3, 3, 0.5), x2 = -0.5, x3 = 0) y3\_0 \textless-
f3(x1 = seq(-3, 3, 0.5), x2 = -0.5, x3 = 0) y4\_0 \textless- f4(x1 =
seq(-3, 3, 0.5), x2 = -0.5, x3 = 0) y5\_0 \textless- f5(x1 = seq(-3, 3,
0.5), x2 = -0.5, x3 = 0)

\bookmarksetup{startatroot}

\chapter{x3 = 1}\label{x3-1}

y1\_1 \textless- f1(x1 = seq(-3, 3, 0.5), x2 = -0.5, x3 = 1) y2\_1
\textless- f2(x1 = seq(-3, 3, 0.5), x2 = -0.5, x3 = 1) y3\_1 \textless-
f3(x1 = seq(-3, 3, 0.5), x2 = -0.5, x3 = 1) y4\_1 \textless- f4(x1 =
seq(-3, 3, 0.5), x2 = -0.5, x3 = 1) y5\_1 \textless- f5(x1 = seq(-3, 3,
0.5), x2 = -0.5, x3 = 1)

\bookmarksetup{startatroot}

\chapter{Create side-by-side plotting
area}\label{create-side-by-side-plotting-area}

par(mfrow = c(1,2))

plot( seq(-3, 3, 0.5), y1\_0, type = ``l'', col = ``black'', ylim =
range(c(y1\_0, y2\_0, y3\_0, y4\_0, y5\_0)), main = ``All Models
Combined (x3 = 0)'', xlab = ``x1'', ylab = ``Predictions for Black Box
Models (1-5)=(F1-F5)'' )

\bookmarksetup{startatroot}

\chapter{lines for all Black models f1-f5 for x3 = 0 and x3 =
1}\label{lines-for-all-black-models-f1-f5-for-x3-0-and-x3-1}

lines(seq(-3, 3, 0.5), y2\_0, col = ``red'', lwd = 2, lty = 2)
lines(seq(-3, 3, 0.5), y3\_0, col = ``blue'', lwd = 2, lty = 3)
lines(seq(-3, 3, 0.5), y4\_0, col = ``green'', lwd = 2, lty = 4)
lines(seq(-3, 3, 0.5), y5\_0, col = ``purple'', lwd = 2, lty = 5)

legend( ``topright'', legend = c(``Model 1 = F1'', ``Model 2 = F2'',
``Model 3 = F3'', ``Model 4 = F4'', ``Model 5 = F5''), col =
c(``black'', ``red'', ``blue'', ``green'', ``purple''), lty =
c(1,2,3,4,5) )

plot( seq(-3, 3, 0.5), y1\_1, type = ``l'', col = ``black'', ylim =
range(c(y1\_1, y2\_1, y3\_1, y4\_1, y5\_1)), main = ``All Models
Combined (x3 = 1)'', xlab = ``x1'', ylab = ``Predictions for Black Box
Models (1-5)=(F1-F5)'' )

lines(seq(-3, 3, 0.5), y2\_1, col = ``red'', lty = 2) lines(seq(-3, 3,
0.5), y3\_1, col = ``blue'', lwd = 2, lty = 3) lines(seq(-3, 3, 0.5),
y4\_1, col = ``green'', lwd = 2, lty = 4) lines(seq(-3, 3, 0.5), y5\_1,
col = ``purple'', lwd = 2, lty = 5)

legend( ``topright'', legend = c(``Model 1 = F1'', ``Model 2 = F2'',
``Model 3 = F3'', ``Model 4 = F4'', ``Model 5 = F5''), col =
c(``black'', ``red'', ``blue'', ``green'', ``purple''), lty =
c(1,2,3,4,5) ) \# ```

\begin{enumerate}
\def\labelenumi{\arabic{enumi}.}
\setcounter{enumi}{3}
\tightlist
\item
  What do we learn from the figures in question 3?
\end{enumerate}

From the figures in question 3, it becomes easy to learn that the model
predictions vary significantly across the five different models (F1 to
F5) when x1 ranges from -3 to 3, with x2 held constant at -0.5 and x3
set to either 0 or 1. These predictions varies uniquely for each black
model value from F1-F5 and x3 values changing and with these variations
we can gain an understanding how each model creates different outcomes
and graph/model behaviors regardless of utilizing the same input
variable x1 from -3 to 3 with incrementation at the value of 0.5.
Altogether, this figure helps us understand that there will always be
differences in predictions and we need a model by model comparison to
better understand these predictions and what is does to the
interpretations and decisions we make based on these models.

\bookmarksetup{startatroot}

\chapter{Homework 2: Data Plots on Measures of
Performance}\label{homework-2-data-plots-on-measures-of-performance}

This Chapter will focus on the topics surrounding Visualizing Data
Models

Topics includes:

\begin{itemize}
\tightlist
\item
  Discrimination (Somer's D/C-index)
\item
  Calibration plot
\item
  Residual plots
\item
  Proper scoring rules
\item
  Expected squared error
\item
  Expected absolute error
\item
  Error distribution Stratification across subgroups
\end{itemize}

\bookmarksetup{startatroot}

\chapter{Module 2 homework}\label{module-2-homework}

\begin{enumerate}
\def\labelenumi{\arabic{enumi}.}
\tightlist
\item
  In your repo, create the HW2.qmd and add it to your book
  (\_quarto.yml).
\item
  For each prediction rule, create two calibration curves: one for x3=0
  and another for x3=1.
\end{enumerate}

\bookmarksetup{startatroot}

\chapter{Installing the properties and packages if
needed:}\label{installing-the-properties-and-packages-if-needed}

\bookmarksetup{startatroot}

\chapter{```\{r\}}\label{r-1}

\#install.packages(``pak'') \#pak::pak(``thomasgstewart/ds6021b'')
\#pak::pak(``thomasgstewart/ds6021c'') library(ds6021b) library(ds6021c)

set.seed(23583) data\_for\_predictionModel1 \textless-
external\_data(n=100)

\bookmarksetup{startatroot}

\chapter{All prediction models for all prediction rules and create two
calibration curves: one for x3=0 and another for
x3=1.}\label{all-prediction-models-for-all-prediction-rules-and-create-two-calibration-curves-one-for-x30-and-another-for-x31.}

data \textless- external\_data(1000)

\bookmarksetup{startatroot}

\chapter{calibration curve for x3=0 and
x3=1}\label{calibration-curve-for-x30-and-x31}

x3\_0 \textless- data\(x3 == 0
x3_1 <- data\)x3 == 1 \# Create side-by-side plotting area par(mfrow =
c(1,2)) \# Loop through each prediction rule and create calibration
curves for x3=0 and x3=1 for (fi in c(``f1'', ``f2'', ``f3'', ``f4'',
``f5'')) \{ \# Create calibration curve for x3=0
plot(get(fi)(data{[}x3\_0, {]}), data{[}x3\_0, {]}\(y,
         xlab = paste("Predicted", fi, " values"),
         ylab = paste("Observed Y values from the Predicted ", fi, " values"),
         main = paste("Prediction Rule Model: ", fi, "at (x3=0)"),
         col = "#aa4848", type = "p")
    lines(lowess(get(fi)(data[x3_0, ]), data[x3_0, ]\)y), col =
``black'') abline(0, 1, lty = 2, lwd = 2, col = ``blue'') \# calibration
curve for x3=1 plot(get(fi)(data{[}x3\_1, {]}), data{[}x3\_1, {]}\(y,
         xlab = paste("Predicted", fi, " values"),
         ylab = paste("Observed Y values from the Predicted ", fi, " values"),
         main = paste("Prediction Rule Model: ", fi, " at (x3=1)"),
         col = "#aa4848", type = "p")
    lines(lowess(get(fi)(data[x3_1, ]), data[x3_1, ]\)y), col =
``black'') abline(0, 1, lty = 2, lwd = 2, col = ``darkgreen'') \}

\begin{enumerate}
\def\labelenumi{\arabic{enumi}.}
\setcounter{enumi}{2}
\item
  For each prediction rule, create a visualization of the residual
  distribution. for (fi in c(``f1'',``f2'', ``f3'', ``f4'', ``f5'')) \{
  resid\_value \textless- get(fi)(data) - data\$y hist(resid\_value,
  freq = FALSE, main = paste(``Residuals model values for the Prediction
  Rule Model:'', fi), xlab = paste(``Residual values (predicted value
  (f1, f2, f3, f4, f5) of y - observed value of y)''), ylab =
  paste(``Probabilities of the Y values from the Predicted'', fi, ''
  values''), col = ``\#aa4848'', border = ``black'', breaks = 30) \#
  Plot the density curve lines(density(resid\_value), col = ``blue'') \}
\item
  For each prediction rule, calculate expected squared error, expected
  absolute error, and the 90th quantile. Place the calculated quantities
  in a visualization for easy comparison. (You may want to create a
  different visualization for each summary measure.)
\end{enumerate}

\bookmarksetup{startatroot}

\chapter{Expected squared error}\label{expected-squared-error}

mean\_squared\_error\_vals \textless-
sapply(c(``f1'',``f2'',``f3'',``f4'',``f5''), function(fi) \{
mean((get(fi)(data) - data\$y)\^{}2) \})

plot(1:5, mean\_squared\_error\_vals, xlab = ``Prediction Rules (f1, f2,
f3, f4, f5)'', ylab = ``Mean Squared Error for each Prediction Rule
Model'', main = ``Expected Squared Error by each Prediction Rule
Model'', col = ``\#0c27bf'', ylim =
c(0,max(mean\_squared\_error\_vals)), pch = 16, cex = 1.5, xaxt = ``n'')
axis(1, 1:5, c(``f1'',``f2'',``f3'',``f4'',``f5''))

\bookmarksetup{startatroot}

\chapter{Expected absolute error}\label{expected-absolute-error}

mae\_vals \textless- sapply(c(``f1'',``f2'',``f3'',``f4'',``f5''),
function(fi) \{ mean(abs(get(fi)(data) - data\$y)) \})

plot(1:5, mae\_vals, xlab = ``Prediction Rules (f1, f2, f3, f4, f5)'',
ylab = ``Mean Absolute Error for each Prediction Rule Model'', main =
``Expected Absolute Error by each Prediction Rule Model'', col =
``\#098001'', ylim = c(0,max(mae\_vals)), pch = 16, cex = 1.5, xaxt =
``n'') axis(1, 1:5, c(``f1'',``f2'',``f3'',``f4'',``f5'')) \# 90th
quantile quantile\_90th\_vals \textless-
sapply(c(``f1'',``f2'',``f3'',``f4'',``f5''), function(fi) \{
quantile(abs(get(fi)(data) - data\$y), 0.9) \})

plot(1:5, quantile\_90th\_vals, xlab = ``Prediction Rules (f1, f2, f3,
f4, f5)'', ylab = ``90th Quantile for each Prediction Rule Model'', main
= ``90th Quantile by each Prediction Rule Model'', col = ``\#ff0000'',
ylim = c(0,max(quantile\_90th\_vals)), pch = 16, cex = 1.5, xaxt =
``n'') axis(1, 1:5, c(``f1'',``f2'',``f3'',``f4'',``f5'')) \# ```

\begin{enumerate}
\def\labelenumi{\arabic{enumi}.}
\setcounter{enumi}{4}
\tightlist
\item
  Describe a measure of model performance from the other section.
\end{enumerate}

After going through all of the sections of this homework assignment, I
found that each graph and predictual rule model revealed some unique
details and insights about the measure of performances and how these
models are used to help make accurate predictions and decions from these
models.

Starting with the calibration curves, these curves are used to evaluate
the performance of a Predictive model by comparing the Predicted (f1,
f2, f3, f4, f5) probabilities of all x related data whose data that
matches the x3 values of 0 and 1 to the observed Observed Y values from
the Predicted (f1, f2, f3, f4, f5) values. The first solid line of the
calibration curves for x3=0 along with x3=1 helped to showed how the
data points were correlated in a smoother linear way to be represent the
relationship between the predicted probabilities and the observed
frequencies. While the second dashed line in the calibration curves for
x3=0 along with x3=1 best represents the reference line of the
calibration curve that helps to represent the ideal straight linear
relationship between the predicted probabilities and the observed
frequencies. Altogether, these calibration curves for x3=0 along with
x3=1 helps us better understand how well the predicted probabilities
from the model align with the actual observed outcomes to see which
relationships in the data better represents the ideal calibration of how
data is predicted and observed in the models. From here, we can make
better decisions based on the accurate prediciton calibrations of the
models and we could also use these models to judge the performances of
data in models in a precise and calibrated way.

Going on into the residual distribution visualization, this histogram
model helps us understand the distribution of residuals, which are the
differences between the predicted values from the models and the
observed values specified in the y variable in the data. After graphing
this histogram, we can see how well the model fits the data into an
almost perfect bell curve distribution of the residuals which helps to
conclude that the model's predictions are well fitted to the observed
data since the highest frequency of residual values are clustered around
zero. This also means that this model's predictions are most perfectly
the same value as the observed y values of the data. If the model was
skewed or sparsely distributed into disjointed clusters, then the data
of the residuals would not be well fitted to the observed data and this
would lead to poor performance and inaccurate assessments of the data.
It is also important to note that the density curve of the residuals
that I added helps reinforce the accuracy and reliability of the model's
predictions by showing the overall shape of the residual distribution.
Altogether, this histogram of residuals for each prediction rule model
(f1, f2, f3, f4, f5) helps us understand how well the model fits the
data and whether there are any patterns or biases in the predictions
along with the data. On top of that, this histogram of residuals for
each prediction rule model (f1, f2, f3, f4, f5) helps us finalize that
data predictions could be relied upon in decision making surrounding the
precise calibration of the data in its later usage and how it does not
deviate away from any previous observed data values.

Finally for the measure of performances of calculate expected squared
error, expected absolute error, and the 90th quantile, anyone could see
by these visualizations that the reliability of the data predictions of
these models could be used to determine accuracy and validity of the
data. For example, within the Expected Squared Error visualization, we
can see that the model reveals that all of the squared error mean values
of each prediction rule (f1, f2, f3, f4, f5) are relatively close to
each other since the cover in the 1 probability range of the y axis.
However, the model with the lowest mean squared error is the most
accurate model in terms of its predictions being close to the observed
values on the y axis. It is important to know that there is one mean
squared error value in the prediction rule f1 that is the highest which
would indicate some potential biased and inaccurates prediction values
to the other prediction rules along with being slightly less accurate
than the other prediction rules. Regardless though, the mean squared
error values reveals that all of the predictions are close in value to
their observed values and thus are closely accurate predictions along
with having a linear relationship regardless if there is no relationship
line to connect them together. When viewing the expected absolute error
rate visualization, we can seem a similar linear pattern in the data
where all of the prediction values of each prediction rule (f1, f2, f3,
f4, f5) are relatively close to each other since the cover in the 0.8
probability range of the y axis. Thus, it would be easy for anyone to
conclude that all of the prediction results of the predicted value to
the observed data values are relatively close to each other and thus are
closely accurate predictions. Finally for the visualization describing
the 90th quantile, we can see that all of the prediction values of each
prediction rule (f1, f2, f3, f4, f5) are relatively close to each other
since the cover in the 1.5 probability range of the y axis. Thus, it
would be easy for anyone to conclude that all of the prediction results
of the predicted value to the observed data values are relatively close
to each other and thus are closely accurate predictions.

Altogether, these visualizations of expected squared error, expected
absolute error, and the 90th quantile help us understand which models
are more accurate and reliable in their predictions and which models may
be less accurate and reliable. This information can be used to make
informed decisions about which models to use for future decisions,
predictions and analysis in comparison to other models.

\bookmarksetup{startatroot}

\chapter{Module 3 homework: Internal measures of
performance}\label{module-3-homework-internal-measures-of-performance}

This Chapter will focus on the topics surrounding Internal measures of
performance

Topics includes:

\begin{itemize}
\tightlist
\item
  Apparent error
\item
  Estimating the error rate and its uncertainty
\item
  K-fold CV -- Bootstrap -- 0.632 \& 0.632+
\item
  Optimism
\end{itemize}

\bookmarksetup{startatroot}

\chapter{Module 3 homework}\label{module-3-homework}

Module 3 homework 1. In your repo, create HW3.qmd and add it to your
book (\_quarto.yml). 2. Calculate the apparent MAE and external MAE for
all 100 models in m using their corresponding training datasets in td.
Create a scatter plot with apparent MAE on the x-axis and external MAE
on the y-axis. Use a large external dataset to calculate MAE. Describe
what you observe.

\bookmarksetup{startatroot}

\chapter{Installing the properties and packages if
needed:}\label{installing-the-properties-and-packages-if-needed-1}

\bookmarksetup{startatroot}

\chapter{```\{r\}}\label{r-2}

install.packages(``pak'', quiet = TRUE)
pak::pak(``thomasgstewart/ds6021d'') library(ds6021d)

\bookmarksetup{startatroot}

\chapter{2.) Calculate the apparent MAE and external MAE for all 100
models in
m}\label{calculate-the-apparent-mae-and-external-mae-for-all-100-models-in-m}

\bookmarksetup{startatroot}

\chapter{using their corresponding training datasets in
td.}\label{using-their-corresponding-training-datasets-in-td.}

\bookmarksetup{startatroot}

\chapter{Create a scatter plot with apparent MAE on the x-axis and
external MAE
on}\label{create-a-scatter-plot-with-apparent-mae-on-the-x-axis-and-external-mae-on}

\bookmarksetup{startatroot}

\chapter{the y-axis. Use a large external dataset to calculate
MAE.}\label{the-y-axis.-use-a-large-external-dataset-to-calculate-mae.}

\bookmarksetup{startatroot}

\chapter{Describe what you observe.}\label{describe-what-you-observe.}

vector\_for\_apparent\_MAE\_100\_models \textless- rep(NA, 100)
vectors\_for\_external\_MAE\_100\_models \textless- rep(NA, 100)
set.seed(2343)

for(index in seq\_along(vector\_for\_apparent\_MAE\_100\_models))\{
m\_functions \textless- m{[}{[}index{]}{]} td\_internal\_training\_data
\textless- td{[}{[}index{]}{]}

\begin{verbatim}
td_internal_training_predition_data_funct_y_out <- m_functions(td_internal_training_data)
td_internal_training_y_data <- td_internal_training_data$y
apparent_rate_MAE <- mean(abs(td_internal_training_predition_data_funct_y_out - td_internal_training_y_data))
vector_for_apparent_MAE_100_models[index] <- apparent_rate_MAE

external_data_vals <- external_data(10000)  # larger external set
external_data_vals_funct_y_out <- m_functions(external_data_vals)
external_rate_MAE <- mean(abs(external_data_vals_funct_y_out - external_data_vals$y))
vectors_for_external_MAE_100_models[index] <- external_rate_MAE
\end{verbatim}

\}

plot(x = vector\_for\_apparent\_MAE\_100\_models, y =
vectors\_for\_external\_MAE\_100\_models, main = ``Apparent MAE vs
External MAE'', xlab = ``Apparent MAE (Training)'', ylab = ``External
MAE (Test)'', pch = 19, col = ``steelblue'', cex = 1.2)

\bookmarksetup{startatroot}

\chapter{3.) For each of the 100 models, compute the
optimism}\label{for-each-of-the-100-models-compute-the-optimism}

\bookmarksetup{startatroot}

\chapter{for MAE (apparent MAE minus external
MAE).}\label{for-mae-apparent-mae-minus-external-mae.}

\bookmarksetup{startatroot}

\chapter{Plot the distribution of
optimism.}\label{plot-the-distribution-of-optimism.}

\bookmarksetup{startatroot}

\chapter{Calculate the mean optimism for
MAE.}\label{calculate-the-mean-optimism-for-mae.}

\bookmarksetup{startatroot}

\chapter{Is the distribution roughly symmetric or
skewed?}\label{is-the-distribution-roughly-symmetric-or-skewed}

\bookmarksetup{startatroot}

\chapter{What does the figure tell you about using apparent
error}\label{what-does-the-figure-tell-you-about-using-apparent-error}

\bookmarksetup{startatroot}

\chapter{as a performance estimate?}\label{as-a-performance-estimate}

vector\_for\_optimism\_MAE\_100\_models \textless- rep(NA, 100)
for(index in seq\_along(vector\_for\_apparent\_MAE\_100\_models))\{
vector\_for\_optimism\_MAE\_100\_models{[}index{]} \textless-
vector\_for\_apparent\_MAE\_100\_models{[}index{]} -
vectors\_for\_external\_MAE\_100\_models{[}index{]} \}
mean\_optimism\_for\_MAE\_100\_models \textless-
mean(vector\_for\_optimism\_MAE\_100\_models)

hist(vector\_for\_optimism\_MAE\_100\_models, main = ``Distribution of
Optimism (Apparent MAE - External MAE)'', xlab = ``Optimism for MAE'',
col = ``lightblue'', border = ``darkblue'') abline(v =
mean\_optimism\_for\_MAE\_100\_models, col = ``red'', lwd = 2) text(x =
mean\_optimism\_for\_MAE\_100\_models, y = -0.5, labels = paste(``Mean
='', round(mean\_optimism\_for\_MAE\_100\_models, 5)), col = ``red'')

\bookmarksetup{startatroot}

\chapter{4.) Select td{[}{[}1{]}{]} as your working
dataset.}\label{select-td1-as-your-working-dataset.}

\bookmarksetup{startatroot}

\chapter{Implement 10-fold CV for the regression
model}\label{implement-10-fold-cv-for-the-regression-model}

\bookmarksetup{startatroot}

\chapter{in the section 3 examples and record the
resulting}\label{in-the-section-3-examples-and-record-the-resulting}

\bookmarksetup{startatroot}

\chapter{MAE estimate.}\label{mae-estimate.}

new\_train\_data \textless- td{[}{[}1{]}{]}
number\_of\_rows\_training\_data \textless- nrow(new\_train\_data)
k\_val \textless- 10 set.seed(123) folds\_randomized \textless-
sample(rep(1:k\_val, length.out = number\_of\_rows\_training\_data))

cv\_mae \textless- rep(NA, k\_val) for(k in 1:k\_val)\{ val\_idx\_of\_CV
\textless- which(folds\_randomized == k) train\_idx\_of\_CV \textless-
which(folds\_randomized != k)

\begin{verbatim}
train_fold_of_CV <- new_train_data[train_idx_of_CV, ]
val_fold_of_CV <- new_train_data[val_idx_of_CV, ]

linear_model_of_fold_of_CV <- lm(y ~ ., data = train_fold_of_CV)
predict_y_val_of_fold_of_CV <- predict(linear_model_of_fold_of_CV, newdata = val_fold_of_CV)

cv_mae[k] <- mean(abs(predict_y_val_of_fold_of_CV - val_fold_of_CV$y))

cat("10-fold CV for the regression model: Fold number: ", k, ", MAE =", round(cv_mae[k], 4), "\n")
\end{verbatim}

\}

\bookmarksetup{startatroot}

\chapter{mean of 10-fold CV for MAE}\label{mean-of-10-fold-cv-for-mae}

mean\_of\_10\_fold\_CV\_for\_MAE \textless- mean(cv\_mae)
print(mean\_of\_10\_fold\_CV\_for\_MAE)

\bookmarksetup{startatroot}

\chapter{5.) Using td{[}{[}1{]}{]}, implement the bootstrap optimism
estimator with
B=200}\label{using-td1-implement-the-bootstrap-optimism-estimator-with-b200}

\bookmarksetup{startatroot}

\chapter{bootstrap samples. Report the apparent MAE, the estimated
optimism,}\label{bootstrap-samples.-report-the-apparent-mae-the-estimated-optimism}

\bookmarksetup{startatroot}

\chapter{and the corrected MAE.}\label{and-the-corrected-mae.}

\bookmarksetup{startatroot}

\chapter{Compare the corrected MAE to the true external error (computed
with a large external
dataset).}\label{compare-the-corrected-mae-to-the-true-external-error-computed-with-a-large-external-dataset.}

b\_val \textless- 200 set.seed(42352) optimism\_b \textless- rep(NA,
b\_val)

m\_function \textless- m{[}{[}1{]}{]}
new\_training\_predition\_data\_funct\_y\_output \textless-
m\_function(new\_train\_data)

apparent\_error\_mae \textless-
mean(abs(new\_training\_predition\_data\_funct\_y\_output -
new\_train\_data\$y)) \# using original m{[}{[}1{]}{]}

for (b in 1:b\_val)\{ idx\_training\_data \textless-
sample(1:number\_of\_rows\_training\_data, replace = TRUE)
bootstrap\_training\_data \textless-
new\_train\_data{[}idx\_training\_data, {]}

\begin{verbatim}
linearly_fittted_model_b_value <- lm(y ~ ., data = bootstrap_training_data)

# In-sample on bootstrap
predicted_y_value_b_value_training_data_bootstrap_in_sample <- predict(linearly_fittted_model_b_value, newdata = bootstrap_training_data)
mae_optimism_in_sample_b_value <- mean(abs(predicted_y_value_b_value_training_data_bootstrap_in_sample - bootstrap_training_data$y))

# original training data
predicted_y_value_b_value_training_data_bootstrap_original <- predict(linearly_fittted_model_b_value, newdata = new_train_data)
mae_optimism_original_b_value <- mean(abs(predicted_y_value_b_value_training_data_bootstrap_original - new_train_data$y))

# optimism for bootstrap formula
optimism_b[b] <- mae_optimism_in_sample_b_value - mae_optimism_original_b_value
\end{verbatim}

\}

\bookmarksetup{startatroot}

\chapter{corrected optimism bootstrap
MAE}\label{corrected-optimism-bootstrap-mae}

bootstrap\_corrected\_mae \textless- apparent\_error\_mae -
mean(optimism\_b)

\bookmarksetup{startatroot}

\chapter{True external error}\label{true-external-error}

large\_external \textless- external\_data(20000)
true\_external\_error\_mae \textless- mean(abs(predict(lm(y
\textasciitilde{} ., data = new\_train\_data), newdata =
large\_external) - large\_external\$y))

cat(``Apparent MAE:'', round(apparent\_error\_mae, 5), ``\n'')
cat(``Bootstrap optimism:'', round(mean(optimism\_b), 5), ``\n'')
cat(``Corrected MAE:'', round(bootstrap\_corrected\_mae, 5), ``\n'')
cat(``True external MAE:'', round(true\_external\_error\_mae, 5),
``\n'')

hist(optimism\_b, main = ``Bootstrap estimates of optimism'', xlab =
``Optimism (in-sample MAE − original MAE)'', freq = FALSE) abline(v =
mean(optimism\_b), col = ``blue'', lwd = 2) text(x = mean(optimism\_b),
y = 0, labels = paste(``Mean ='', round(mean\_optimism, 5)), col =
``blue'') \# ```

\bookmarksetup{startatroot}

\chapter{Name: Maxwell Braseth}\label{name-maxwell-braseth}

\bookmarksetup{startatroot}

\chapter{Date: 06/14/2026}\label{date-06142026}

\bookmarksetup{startatroot}

\chapter{Module 4 Homework}\label{module-4-homework}

\bookmarksetup{startatroot}

\chapter{```\{r\}}\label{r-3}

install.packages(``plotly'') install.packages(``dplyr'')
install.packages(``ggplot2'') library(plotly) library(dplyr)
library(ggplot2) set.seed(42)

\bookmarksetup{startatroot}

\chapter{1.)}\label{section}

\bookmarksetup{startatroot}

\chapter{In your repo, create the HW4.qmd and add it to your
book}\label{in-your-repo-create-the-hw4.qmd-and-add-it-to-your-book}

\bookmarksetup{startatroot}

\chapter{(\_quarto.yml).}\label{quarto.yml.}

\bookmarksetup{startatroot}

\chapter{2.)}\label{section-1}

\bookmarksetup{startatroot}

\chapter{Suppose X1 has 4
categories,}\label{suppose-x1-has-4-categories}

\bookmarksetup{startatroot}

\chapter{X2 has 6 categories, and}\label{x2-has-6-categories-and}

\bookmarksetup{startatroot}

\chapter{X3 has 3 categories.}\label{x3-has-3-categories.}

\bookmarksetup{startatroot}

\chapter{X1 ∈ \{1, 2, 3, 4\} --\textgreater{} 4
categories}\label{x1-1-2-3-4-4-categories}

\bookmarksetup{startatroot}

\chapter{X2 ∈ \{1, 2, 3, 4, 5, 6\} --\textgreater{} 6
categories}\label{x2-1-2-3-4-5-6-6-categories}

\bookmarksetup{startatroot}

\chapter{X3 ∈ \{1, 2, 3\} --\textgreater{} 3
categories}\label{x3-1-2-3-3-categories}

\bookmarksetup{startatroot}

\chapter{X1 * X2 * X3 = 4 * 6 * 3 = 72
cells}\label{x1-x2-x3-4-6-3-72-cells}

x1 \textless- c(1, 2, 3, 4) x2 \textless- c(1, 2, 3, 4, 5, 6) x3
\textless- c(1, 2, 3)

\bookmarksetup{startatroot}

\chapter{How many cells does the largest possible model
contain?}\label{how-many-cells-does-the-largest-possible-model-contain}

total\_cells \textless- (length(x1) * length(x2) * length(x3)) \# (x1 *
x2 * x3) print(paste(``The maximum number of cells within the model
is:'', total\_cells)) print(paste(``This is because the largest possible
model (the saturated model) requires a unique cell for every possible
joint combination of the categories, calculated as 4 × 6 × 3 = 72
cells.'')) print(paste(``These 72 cells represent the complete
classification of all three covariates, spanning from (1, 1, 1) through
(4, 6, 3), allowing the model to capture all of the independently
predicted values of the combined covariate cell values.''))

\bookmarksetup{startatroot}

\chapter{How many parameters does the smallest possible model
require?}\label{how-many-parameters-does-the-smallest-possible-model-require}

parameters\_for\_smallest\_model \textless- 1 print(paste(``The minimum
number of parameters within the model is:'',
parameters\_for\_smallest\_model)) print(paste(``This is because the
smallest possible model (the intercept-only model) which ignores all
covariate information and assigns the same prediction value'', ``to
every observation regardless of the combinations of the values of X1,
X2, and X3.'')) print(paste(``At the same time, this model needs a
intercept constant value that provides the same prediction value
regardless of the combinations of the covariates values of x1, x2, and
x3. This also helps estimate the mean of this smallest model.''))

\bookmarksetup{startatroot}

\chapter{If a training dataset of n=30 observations is drawn uniformly
at}\label{if-a-training-dataset-of-n30-observations-is-drawn-uniformly-at}

\bookmarksetup{startatroot}

\chapter{random from this covariate
space,}\label{random-from-this-covariate-space}

\bookmarksetup{startatroot}

\chapter{roughly how many cells would you expect to be
unobserved?}\label{roughly-how-many-cells-would-you-expect-to-be-unobserved}

\bookmarksetup{startatroot}

\chapter{Set parameters}\label{set-parameters}

set.seed(42) n \textless- 30

\bookmarksetup{startatroot}

\chapter{Mathematical Formula for the Expected Number of Unobserved
Cells}\label{mathematical-formula-for-the-expected-number-of-unobserved-cells}

\bookmarksetup{startatroot}

\chapter{E(cells\_unobserved = UC) = total\_cells = M * (1 - 1 /
total\_cells +
M)\^{}n}\label{ecells_unobserved-uc-total_cells-m-1---1-total_cells-mn}

cells\_unobserved\_expected \textless- total\_cells * (1 - 1 /
total\_cells)\^{}n print(paste(``Theoretical expected unobserved
cells:'', round(cells\_unobserved\_expected, 0)))

\bookmarksetup{startatroot}

\chapter{graph explanation to the observed and unobserved
cells}\label{graph-explanation-to-the-observed-and-unobserved-cells}

\bookmarksetup{startatroot}

\chapter{Continuous Uniform
Simulation}\label{continuous-uniform-simulation}

\bookmarksetup{startatroot}

\chapter{uniformly between 0 and 1 for n=30 observations for x, y, and
z}\label{uniformly-between-0-and-1-for-n30-observations-for-x-y-and-z}

idx\_x1 \textless- sample(x1, n, replace = TRUE) idx\_x2 \textless-
sample(x2, n, replace = TRUE) idx\_x3 \textless- sample(x3, n, replace =
TRUE) \# creating the cell matrix/grid grid\_matrix \textless-
expand.grid( x1 = x1, x2 = x2, x3 = x3 ) \# Initialize observation count
grid\_matrix\(Obs_Count <- 0
# Looping throught the total n observations and filling the matrix/grid cells
for(i in 1:n){
  matched_row <- which(
    grid_matrix\)x1 == idx\_x1{[}i{]} \& grid\_matrix\(x2 == idx_x2[i] &
    grid_matrix\)x3 == idx\_x3{[}i{]} ) \# incrementing the observation
count of matched values the covariates of x1=4, x2=6, and x3=3
grid\_matrix\(Obs_Count[matched_row] <-
    grid_matrix\)Obs\_Count{[}matched\_row{]} + 1 \}

\bookmarksetup{startatroot}

\chapter{Number of unobserved cells and observed
cells}\label{number-of-unobserved-cells-and-observed-cells}

simulated\_unobserved \textless- sum(grid\_matrix\(Obs_Count == 0)
simulated_observed <- sum(grid_matrix\)Obs\_Count \textgreater{} 0)

print(paste(``Simulated unobserved cells:'', simulated\_unobserved))
print(paste(``Simulated observed cells:'', simulated\_observed))
print(paste(``Difference between simulated and expected unobserved
cells:'',round(simulated\_unobserved - cells\_unobserved\_expected, 0)))

\bookmarksetup{startatroot}

\chapter{3.)}\label{section-2}

\bookmarksetup{startatroot}

\chapter{Using the linear interpolation model, calculate the prediction
ŷ(X0 =
0.45)}\label{using-the-linear-interpolation-model-calculate-the-prediction-yx0-0.45}

\bookmarksetup{startatroot}

\chapter{Wa = (Xb - X0) / (Xb - Xa)}\label{wa-xb---x0-xb---xa}

\bookmarksetup{startatroot}

\chapter{Wb = (X0 - Xa) / (Xb - Xa)}\label{wb-x0---xa-xb---xa}

\bookmarksetup{startatroot}

\chapter{ŷ(X0) = Wa * ya + Wb * yb}\label{yx0-wa-ya-wb-yb}

\bookmarksetup{startatroot}

\chapter{given training points and the observed value of x0: x0=0.45, xa
= 0.38, ya = 1.92, xb = 0.57, yb =
2.04.}\label{given-training-points-and-the-observed-value-of-x0-x00.45-xa-0.38-ya-1.92-xb-0.57-yb-2.04.}

x0 \textless- 0.45 ya \textless- 1.92 xa \textless- 0.38 xb \textless-
0.57 yb \textless- 2.04 \# Show the weights Wa and Wb wa \textless- (xb
- x0) / (xb - xa) wb \textless- (x0 - xa) / (xb - xa) \# setting the
value of ŷ(X0) prediction\_y\_x0 \textless- wa * ya + wb * yb \#
verifying that the weights of Wa and Wb sum to 1 print(paste(``The sum
of the weights is:'', wa + wb)) print(paste(``Difference between the
prediction value of ŷ(X0) and the total weights of wa \& wb:'',
round(prediction\_y\_x0 - (wa + wb), 0)))

\bookmarksetup{startatroot}

\chapter{what is the relative magnitudes of the weights, tell you about
how
much}\label{what-is-the-relative-magnitudes-of-the-weights-tell-you-about-how-much}

\bookmarksetup{startatroot}

\chapter{each neighbor contributes to the
prediction.}\label{each-neighbor-contributes-to-the-prediction.}

\bookmarksetup{startatroot}

\chapter{given that the neighbors are considered to be the two training
observations that help show the prediction of
ŷ(X0)}\label{given-that-the-neighbors-are-considered-to-be-the-two-training-observations-that-help-show-the-prediction-of-yx0}

\bookmarksetup{startatroot}

\chapter{and the neighbors are the left neighbor: wa=(Xa = 0.38, ya =
1.92) and the right neighbor: wb=(Xb = 0.57, yb =
2.04).}\label{and-the-neighbors-are-the-left-neighbor-waxa-0.38-ya-1.92-and-the-right-neighbor-wbxb-0.57-yb-2.04.}

\bookmarksetup{startatroot}

\chapter{At the same time, these neighbords are closets values to the
observed value of x0 to the left and right of this observed
value.}\label{at-the-same-time-these-neighbords-are-closets-values-to-the-observed-value-of-x0-to-the-left-and-right-of-this-observed-value.}

\bookmarksetup{startatroot}

\chapter{relative magnitude(s) of the
weights}\label{relative-magnitudes-of-the-weights}

print(paste(``Relative magnitude ratio for wa/wb:'', round(wa/wb), 2))
print(paste(``Relative magnitude ratio for wb/wa:'', round(wb/wa), 2))
print(paste(``Relative magnitude of the left neighbor:'', round(wa / (wa
+ wb), 2))) print(paste(``Relative magnitude of the right neighbor:'',
round(wb / (wa + wb), 2))) print(paste(``Difference between the relative
magnitude of the left neighbor and the right neighbor:'', round(wa / (wa
+ wb) - wb / (wa + wb), 0))) print(paste(``Is the left or the right
neighbor larger or smaller?:'', ``The left neighbor is larger than the
right neighbor?:'', wa \textgreater{} wb, '' The right neighbor is
larger than the left neighbor?: ``, wb \textgreater{} wa)) if(wa
\textgreater{} wb)\{ print(paste(''The left neighbor is larger than the
right neighbor and this means that the relative magnitude of the left
neighbor contributes more to the prediction of ŷ(X0) at a wa value of :
``, round(wa, 2) * 100,''\%``)) print(paste(''This also means that the
left neighbor wa is closer to the observed value of x0 than the right
neighbor wb.'')) print(paste(``Which also means that this linear
interpolation model is better at predicting ŷ(X0) at a wa value of :'',
round(wa, 2) * 100, ``\%'')) \}else\{ print(paste(``The right neighbor
is larger than the left neighbor and this means that the relative
magnitude of the right neighbor contributes more to the prediction of
ŷ(X0) at a wb value of :'', round(wb, 2) * 100, ``\%''))
print(paste(``This also means that the right neighbor wb is closer to
the observed value of x0 than the left neighbor wa.''))
print(paste(``Which also means that this linear interpolation model is
better at predicting ŷ(X0) at a wb value of :'', round(wb, 2) * 100,
``\%'')) \}

\bookmarksetup{startatroot}

\chapter{4.) Using the following data, calculate the 3 nearest neigbhor
prediction}\label{using-the-following-data-calculate-the-3-nearest-neigbhor-prediction}

\bookmarksetup{startatroot}

\chapter{for X0 = 0.2, 0.5, 0.7 by
hand.}\label{for-x0-0.2-0.5-0.7-by-hand.}

\bookmarksetup{startatroot}

\chapter{x y}\label{x-y}

\bookmarksetup{startatroot}

\chapter{0.22 2.14}\label{section-3}

\bookmarksetup{startatroot}

\chapter{0.37 1.89}\label{section-4}

\bookmarksetup{startatroot}

\chapter{0.57 1.99}\label{section-5}

\bookmarksetup{startatroot}

\chapter{0.66 1.88}\label{section-6}

\bookmarksetup{startatroot}

\chapter{0.03 2.26}\label{section-7}

\bookmarksetup{startatroot}

\chapter{0.92 2.11}\label{section-8}

\bookmarksetup{startatroot}

\chapter{0.77 1.92}\label{section-9}

\bookmarksetup{startatroot}

\chapter{0.73 1.98}\label{section-10}

\bookmarksetup{startatroot}

\chapter{0.11 2.47}\label{section-11}

\bookmarksetup{startatroot}

\chapter{0.40 2.10}\label{section-12}

\bookmarksetup{startatroot}

\chapter{Before selecting the 3 closest k=3 neighbors, measuring the
distance between the new observed point (x) and a known training point
(y) across (n) different
features.}\label{before-selecting-the-3-closest-k3-neighbors-measuring-the-distance-between-the-new-observed-point-x-and-a-known-training-point-y-across-n-different-features.}

\bookmarksetup{startatroot}

\chapter{Distance Metrics Formulas:}\label{distance-metrics-formulas}

\bookmarksetup{startatroot}

\chapter{Euclidean Distance: Σ(√(xi-yi)\^{}2) - yi : sqrt((x2 - x1)\^{}2
+ (y2 -
y1)\^{}2)}\label{euclidean-distance-ux3c3xi-yi2---yi-sqrtx2---x12-y2---y12}

\bookmarksetup{startatroot}

\chapter{Euclidean Distance in this context: Σ(x0-xi) - yi : \textbar x2
-
x1\textbar{}}\label{euclidean-distance-in-this-context-ux3c3x0-xi---yi-x2---x1}

\bookmarksetup{startatroot}

\chapter{Manhatten Distance: Σ\textbar xi - yi\textbar{} - yi :
\textbar x2 - x1\textbar{} + \textbar y2 -
y1\textbar{}}\label{manhatten-distance-ux3c3xi---yi---yi-x2---x1-y2---y1}

\bookmarksetup{startatroot}

\chapter{3 nearest neighbors for x0 =
0.2}\label{nearest-neighbors-for-x0-0.2}

\bookmarksetup{startatroot}

\chapter{3 nearest neighbors for x0 =
0.5}\label{nearest-neighbors-for-x0-0.5}

\bookmarksetup{startatroot}

\chapter{3 nearest neighbors for x0 =
0.7}\label{nearest-neighbors-for-x0-0.7}

\bookmarksetup{startatroot}

\chapter{where x0 = 0.2}\label{where-x0-0.2}

\bookmarksetup{startatroot}

\chapter{f(xi) = Σ(x0-xi)}\label{fxi-ux3c3x0-xi}

\bookmarksetup{startatroot}

\chapter{f(x=0.22) = 0.20 - 0.22 = \textbar-0.02\textbar{} = 0.02
--\textgreater{} y = 2.14}\label{fx0.22-0.20---0.22--0.02-0.02-y-2.14}

\bookmarksetup{startatroot}

\chapter{f(x=0.37) = 0.20 - 0.37 = \textbar-0.17\textbar{} = 0.17
--\textgreater{} y = 1.89}\label{fx0.37-0.20---0.37--0.17-0.17-y-1.89}

\bookmarksetup{startatroot}

\chapter{f(x=0.57) = 0.20 - 0.57 = \textbar-0.37\textbar{} = 0.37
--\textgreater{} y = 1.99}\label{fx0.57-0.20---0.57--0.37-0.37-y-1.99}

\bookmarksetup{startatroot}

\chapter{f(x=0.66) = 0.20 - 0.66 = \textbar-0.54\textbar{} = 0.54
--\textgreater{} y = 1.88}\label{fx0.66-0.20---0.66--0.54-0.54-y-1.88}

\bookmarksetup{startatroot}

\chapter{f(x=0.03) = 0.20 - 0.03 = \textbar0.17\textbar{} = 0.17
--\textgreater{} y = 2.26}\label{fx0.03-0.20---0.03-0.17-0.17-y-2.26}

\bookmarksetup{startatroot}

\chapter{f(x=0.92) = 0.20 - 0.92 = \textbar-0.72\textbar{} = 0.72
--\textgreater{} y = 2.11}\label{fx0.92-0.20---0.92--0.72-0.72-y-2.11}

\bookmarksetup{startatroot}

\chapter{f(x=0.77) = 0.20 - 0.77 = \textbar-0.43\textbar{} = 0.43
--\textgreater{} y = 1.92}\label{fx0.77-0.20---0.77--0.43-0.43-y-1.92}

\bookmarksetup{startatroot}

\chapter{f(x=0.73) = 0.20 - 0.73 = \textbar-0.46\textbar{} = 0.46
--\textgreater{} y = 1.98}\label{fx0.73-0.20---0.73--0.46-0.46-y-1.98}

\bookmarksetup{startatroot}

\chapter{f(x=0.11) = 0.20 - 0.11 = \textbar0.09\textbar{} = 0.09
--\textgreater{} y = 2.47}\label{fx0.11-0.20---0.11-0.09-0.09-y-2.47}

\bookmarksetup{startatroot}

\chapter{f(x=0.40) = 0.20 - 0.40 = \textbar-0.20\textbar{} = 0.20
--\textgreater{} y = 2.10}\label{fx0.40-0.20---0.40--0.20-0.20-y-2.10}

\bookmarksetup{startatroot}

\chapter{where x0 = 0.5}\label{where-x0-0.5}

\bookmarksetup{startatroot}

\chapter{f(xi) = Σ(x0-xi)}\label{fxi-ux3c3x0-xi-1}

\bookmarksetup{startatroot}

\chapter{f(x=0.22) = 0.50 - 0.22 = \textbar0.28\textbar{} = 0.28
--\textgreater{} y = 2.14}\label{fx0.22-0.50---0.22-0.28-0.28-y-2.14}

\bookmarksetup{startatroot}

\chapter{f(x=0.37) = 0.50 - 0.37 = \textbar0.13\textbar{} = 0.13
--\textgreater{} y = 1.89}\label{fx0.37-0.50---0.37-0.13-0.13-y-1.89}

\bookmarksetup{startatroot}

\chapter{f(x=0.57) = 0.50 - 0.57 = \textbar-0.07\textbar{} = 0.07
--\textgreater{} y = 1.99}\label{fx0.57-0.50---0.57--0.07-0.07-y-1.99}

\bookmarksetup{startatroot}

\chapter{f(x=0.66) = 0.50 - 0.66 = \textbar-0.19\textbar{} = 0.19
--\textgreater{} y = 1.88}\label{fx0.66-0.50---0.66--0.19-0.19-y-1.88}

\bookmarksetup{startatroot}

\chapter{f(x=0.03) = 0.50 - 0.03 = \textbar0.47\textbar{} = 0.47
--\textgreater{} y = 2.26}\label{fx0.03-0.50---0.03-0.47-0.47-y-2.26}

\bookmarksetup{startatroot}

\chapter{f(x=0.92) = 0.50 - 0.92 = \textbar-0.58\textbar{} = 0.58
--\textgreater{} y = 2.11}\label{fx0.92-0.50---0.92--0.58-0.58-y-2.11}

\bookmarksetup{startatroot}

\chapter{f(x=0.77) = 0.50 - 0.77 = \textbar-0.23\textbar{} = 0.23
--\textgreater{} y = 1.92}\label{fx0.77-0.50---0.77--0.23-0.23-y-1.92}

\bookmarksetup{startatroot}

\chapter{f(x=0.73) = 0.50 - 0.73 = \textbar-0.27\textbar{} = 0.27
--\textgreater{} y = 1.98}\label{fx0.73-0.50---0.73--0.27-0.27-y-1.98}

\bookmarksetup{startatroot}

\chapter{f(x=0.11) = 0.50 - 0.11 = \textbar0.39\textbar{} = 0.39
--\textgreater{} y = 2.47}\label{fx0.11-0.50---0.11-0.39-0.39-y-2.47}

\bookmarksetup{startatroot}

\chapter{f(x=0.40) = 0.50 - 0.40 = \textbar-0.10\textbar{} = 0.10
--\textgreater{} y = 2.10}\label{fx0.40-0.50---0.40--0.10-0.10-y-2.10}

\bookmarksetup{startatroot}

\chapter{where x0 = 0.7}\label{where-x0-0.7}

\bookmarksetup{startatroot}

\chapter{f(xi) = Σ(x0-xi)}\label{fxi-ux3c3x0-xi-2}

\bookmarksetup{startatroot}

\chapter{f(x=0.22) = 0.70 - 0.22 = \textbar0.48\textbar{} = 0.48
--\textgreater{} y = 2.14}\label{fx0.22-0.70---0.22-0.48-0.48-y-2.14}

\bookmarksetup{startatroot}

\chapter{f(x=0.37) = 0.70 - 0.37 = \textbar0.33\textbar{} = 0.33
--\textgreater{} y = 1.89}\label{fx0.37-0.70---0.37-0.33-0.33-y-1.89}

\bookmarksetup{startatroot}

\chapter{f(x=0.57) = 0.70 - 0.57 = \textbar0.13\textbar{} = 0.13
--\textgreater{} y = 1.99}\label{fx0.57-0.70---0.57-0.13-0.13-y-1.99}

\bookmarksetup{startatroot}

\chapter{f(x=0.66) = 0.70 - 0.66 = \textbar-0.02\textbar{} = 0.02
--\textgreater{} y = 1.88}\label{fx0.66-0.70---0.66--0.02-0.02-y-1.88}

\bookmarksetup{startatroot}

\chapter{f(x=0.03) = 0.70 - 0.03 = \textbar0.67\textbar{} = 0.67
--\textgreater{} y = 2.26}\label{fx0.03-0.70---0.03-0.67-0.67-y-2.26}

\bookmarksetup{startatroot}

\chapter{f(x=0.92) = 0.70 - 0.92 = \textbar-0.08\textbar{} = 0.08
--\textgreater{} y = 2.11}\label{fx0.92-0.70---0.92--0.08-0.08-y-2.11}

\bookmarksetup{startatroot}

\chapter{f(x=0.77) = 0.70 - 0.77 = \textbar-0.03\textbar{} = 0.03
--\textgreater{} y = 1.92}\label{fx0.77-0.70---0.77--0.03-0.03-y-1.92}

\bookmarksetup{startatroot}

\chapter{f(x=0.73) = 0.70 - 0.73 = \textbar-0.07\textbar{} = 0.07
--\textgreater{} y = 1.98}\label{fx0.73-0.70---0.73--0.07-0.07-y-1.98}

\bookmarksetup{startatroot}

\chapter{f(x=0.11) = 0.70 - 0.11 = \textbar0.59\textbar{} = 0.59
--\textgreater{} y = 2.47}\label{fx0.11-0.70---0.11-0.59-0.59-y-2.47}

\bookmarksetup{startatroot}

\chapter{f(x=0.40) = 0.70 - 0.40 = \textbar0.30\textbar{} = 0.30
--\textgreater{} y = 2.10}\label{fx0.40-0.70---0.40-0.30-0.30-y-2.10}

\bookmarksetup{startatroot}

\chapter{The top three closest neighbors for x0 = 0.2
are:}\label{the-top-three-closest-neighbors-for-x0-0.2-are}

\bookmarksetup{startatroot}

\chapter{y = 2.14 due to ŷ(x = 0.22)=0.02, y = 2.47 due to ŷ(x =
0.11)=0.09, y = 1.89 due to
ŷ(x=0.37)=0.17}\label{y-2.14-due-to-yx-0.220.02-y-2.47-due-to-yx-0.110.09-y-1.89-due-to-yx0.370.17}

\bookmarksetup{startatroot}

\chapter{3 nearest neigbhor prediction for x0 = 0.2 for ŷ(x0 = 0.2) =
(2.14 + 2.47 + 1.89) / 3 = 2.1667 =
2.17}\label{nearest-neigbhor-prediction-for-x0-0.2-for-yx0-0.2-2.14-2.47-1.89-3-2.1667-2.17}

\bookmarksetup{startatroot}

\chapter{The top three closest neighbors for x0 = 0.5
are:}\label{the-top-three-closest-neighbors-for-x0-0.5-are}

\bookmarksetup{startatroot}

\chapter{ŷ(x=0.57)= 0.07 --\textgreater{} y = 1.99, ŷ(x = 0.37)=0.13
--\textgreater{} y = 1.89, ŷ(x = 0.40)=0.10 --\textgreater{} y =
2.10}\label{yx0.57-0.07-y-1.99-yx-0.370.13-y-1.89-yx-0.400.10-y-2.10}

\bookmarksetup{startatroot}

\chapter{3 nearest neigbhor prediction for x0 = 0.5 for ŷ(x0 = 0.5) =
(1.99 + 1.89 + 2.10) / 3 = 1.99333 =
1.99}\label{nearest-neigbhor-prediction-for-x0-0.5-for-yx0-0.5-1.99-1.89-2.10-3-1.99333-1.99}

\bookmarksetup{startatroot}

\chapter{The top three closest neighbors for x0 = 0.7
are:}\label{the-top-three-closest-neighbors-for-x0-0.7-are}

\bookmarksetup{startatroot}

\chapter{ŷ(x=0.66)= 0.02 --\textgreater{} y = 1.88, ŷ(x=0.77)=0.03
--\textgreater{} y = 1.92, ŷ(x=0.73)=0.07 --\textgreater{} y =
1.98}\label{yx0.66-0.02-y-1.88-yx0.770.03-y-1.92-yx0.730.07-y-1.98}

\bookmarksetup{startatroot}

\chapter{3 nearest neigbhor prediction for x0 = 0.7 for ŷ(x0 = 0.7) =
(1.88 + 1.92 + 1.98) / 3 = 1.926667 =
1.93}\label{nearest-neigbhor-prediction-for-x0-0.7-for-yx0-0.7-1.88-1.92-1.98-3-1.926667-1.93}

\bookmarksetup{startatroot}

\chapter{5.) Generate a plot of the 3 nearest neigbhor
prediction}\label{generate-a-plot-of-the-3-nearest-neigbhor-prediction}

\bookmarksetup{startatroot}

\chapter{function overlayed on top of the scatterplot of x and
y.}\label{function-overlayed-on-top-of-the-scatterplot-of-x-and-y.}

\bookmarksetup{startatroot}

\chapter{You may do this by hand or by
code.}\label{you-may-do-this-by-hand-or-by-code.}

\bookmarksetup{startatroot}

\chapter{x y}\label{x-y-1}

\bookmarksetup{startatroot}

\chapter{0.22 2.14}\label{section-13}

\bookmarksetup{startatroot}

\chapter{0.37 1.89}\label{section-14}

\bookmarksetup{startatroot}

\chapter{0.57 1.99}\label{section-15}

\bookmarksetup{startatroot}

\chapter{0.66 1.88}\label{section-16}

\bookmarksetup{startatroot}

\chapter{0.03 2.26}\label{section-17}

\bookmarksetup{startatroot}

\chapter{0.92 2.11}\label{section-18}

\bookmarksetup{startatroot}

\chapter{0.77 1.92}\label{section-19}

\bookmarksetup{startatroot}

\chapter{0.73 1.98}\label{section-20}

\bookmarksetup{startatroot}

\chapter{0.11 2.47}\label{section-21}

\bookmarksetup{startatroot}

\chapter{0.40 2.10}\label{section-22}

data\_original \textless- data.frame( x = c(0.22, 0.37, 0.57, 0.66,
0.03, 0.92, 0.77, 0.73, 0.11, 0.40), y = c(2.14, 1.89, 1.99, 1.88, 2.26,
2.11, 1.92, 1.98, 2.47, 2.10) )

x0\_values \textless- c(0.2, 0.5, 0.7)

for (x0 in x0\_values) \{ smallest\_dist \textless- data\_original
\%\textgreater\% mutate(dist = abs(x - x0)) \%\textgreater\%
arrange(dist) \%\textgreater\% slice(1:3)

predict\_y\_vals\_from\_x0 \textless- mean(smallest\_dist\$y)

cat(``For X0 ='', x0, ``\n'') print(smallest\_dist{[}, c(``x'', ``y'',
``dist''){]}) cat(``Prediction ŷ ='', round(predict\_y\_vals\_from\_x0,
3), ``\n\n'') \}

\bookmarksetup{startatroot}

\chapter{Data from the table above}\label{data-from-the-table-above}

x\_data\_new \textless- seq(0, 1, length.out = n)

\bookmarksetup{startatroot}

\chapter{3-NN prediction function using the ORIGINAL
data}\label{nn-prediction-function-using-the-original-data}

knn3\_pred \textless- function(x\_new, original\_data\_1, k=3) \{
sapply(x\_new, function(xi) \{ data \textless- original\_data\_1
\%\textgreater\% mutate(dist = abs(original\_data\_1\(x - xi))
    mean((data %>% arrange(dist) %>% slice(1:k))
\)y) \}) \}

preds \textless- knn3\_pred(x\_data\_new, data\_original)

\bookmarksetup{startatroot}

\chapter{R scatter plot}\label{r-scatter-plot}

plot(data\_original\(x, data_original\)y, pch = 16, col = ``steelblue'',
cex = 1.5, main = ``3-Nearest Neighbor Prediction (k=3)'', xlab = ``X'',
ylab = ``Y'', ylim = range(c(data\_original\$y, preds)))

\bookmarksetup{startatroot}

\chapter{Add the 3-NN prediction
line}\label{add-the-3-nn-prediction-line}

lines(x\_data\_new, preds, col = ``tomato'', lwd = 2)

\bookmarksetup{startatroot}

\chapter{Adding grid and legend}\label{adding-grid-and-legend}

grid() legend(``topright'', legend = c(``New X Data: Training Data'',
``3-NN Prediction (k=3) value from the original data''), col =
c(``steelblue'', ``tomato''), pch = c(16, NA), lty = c(NA, 1), lwd =
c(NA, 2), bty = ``n'')

\#6.) In your own words, explain why the ``largest model'' \# --- which
assigns an independent estimate to every covariate \# combination ---
can be a poor choice in practice even though it \# fits the training
data perfectly. \# Connect your answer to the concept of information
sharing introduced \# in the chapter, and describe what characteristic
of the data-generating \# process would make the largest model a
reasonable choice.

\bookmarksetup{startatroot}

\chapter{When considering the topic of why the ``Largest Model''
independently predicts and estimates the outcome for each cell in the
combinations of covariates, it is important to recognize that the model
assigns a separate prediction to every possible combination of covariate
values (for example, x1 = 4 categories, x2 = 6 categories, and x3 = 3
categories, resulting in 72 possible
cells).}\label{when-considering-the-topic-of-why-the-largest-model-independently-predicts-and-estimates-the-outcome-for-each-cell-in-the-combinations-of-covariates-it-is-important-to-recognize-that-the-model-assigns-a-separate-prediction-to-every-possible-combination-of-covariate-values-for-example-x1-4-categories-x2-6-categories-and-x3-3-categories-resulting-in-72-possible-cells.}

\bookmarksetup{startatroot}

\chapter{The largest model can be a poor choice in practice because it
simply tracks the mean outcome within each observed cell and does not
share information across cells. While this allows the model to fit the
training data perfectly, it does not provide a reliable way to make
predictions for cells that are unobserved or contain very few
observations. If some covariate combinations are not represented in the
training data, then no cell-specific mean can be calculated for those
combinations.}\label{the-largest-model-can-be-a-poor-choice-in-practice-because-it-simply-tracks-the-mean-outcome-within-each-observed-cell-and-does-not-share-information-across-cells.-while-this-allows-the-model-to-fit-the-training-data-perfectly-it-does-not-provide-a-reliable-way-to-make-predictions-for-cells-that-are-unobserved-or-contain-very-few-observations.-if-some-covariate-combinations-are-not-represented-in-the-training-data-then-no-cell-specific-mean-can-be-calculated-for-those-combinations.}

\bookmarksetup{startatroot}

\chapter{This lack of information sharing can result in overfitting. The
model may appear to perform extremely well on the observed training data
because each cell is estimated independently, but its predictions may
not generalize well to new data. In other words, the apparent error on
the training data may be very small, while the prediction error on
future observations may be much
larger.}\label{this-lack-of-information-sharing-can-result-in-overfitting.-the-model-may-appear-to-perform-extremely-well-on-the-observed-training-data-because-each-cell-is-estimated-independently-but-its-predictions-may-not-generalize-well-to-new-data.-in-other-words-the-apparent-error-on-the-training-data-may-be-very-small-while-the-prediction-error-on-future-observations-may-be-much-larger.}

\bookmarksetup{startatroot}

\chapter{When considering the concept of Information Sharing in this
context, information sharing describes how observations contribute to
predictions beyond their own cell. The best models often lie between the
largest model (no information sharing) and the smallest model (maximum
information sharing). These models allow observations to influence
predictions in neighboring cells or groups, either locally (such as
nearest-neighbor methods) or globally (such as regression
models).}\label{when-considering-the-concept-of-information-sharing-in-this-context-information-sharing-describes-how-observations-contribute-to-predictions-beyond-their-own-cell.-the-best-models-often-lie-between-the-largest-model-no-information-sharing-and-the-smallest-model-maximum-information-sharing.-these-models-allow-observations-to-influence-predictions-in-neighboring-cells-or-groups-either-locally-such-as-nearest-neighbor-methods-or-globally-such-as-regression-models.}

\bookmarksetup{startatroot}

\chapter{Keeping this in mind, the largest model would be considered
reasonable if the data-generating process truly produced independent
outcomes for every covariate combination and there were sufficient
observations within every cell to estimate the cell-specific mean
accurately. In that situation, little or no information sharing would be
needed because each cell could be reliably estimated using its own
observed
data.}\label{keeping-this-in-mind-the-largest-model-would-be-considered-reasonable-if-the-data-generating-process-truly-produced-independent-outcomes-for-every-covariate-combination-and-there-were-sufficient-observations-within-every-cell-to-estimate-the-cell-specific-mean-accurately.-in-that-situation-little-or-no-information-sharing-would-be-needed-because-each-cell-could-be-reliably-estimated-using-its-own-observed-data.}

\bookmarksetup{startatroot}

\chapter{Finally, the largest model would be acceptable when all
covariate combinations are well represented in the dataset and each cell
contains enough observations to produce stable predictions. Under those
conditions, the model could accurately describe the data-generating
process without needing to borrow information from neighboring cells or
groups.}\label{finally-the-largest-model-would-be-acceptable-when-all-covariate-combinations-are-well-represented-in-the-dataset-and-each-cell-contains-enough-observations-to-produce-stable-predictions.-under-those-conditions-the-model-could-accurately-describe-the-data-generating-process-without-needing-to-borrow-information-from-neighboring-cells-or-groups.}

\bookmarksetup{startatroot}

\chapter{```}\label{section-23}

\bookmarksetup{startatroot}

\chapter{Name: Maxwell Braseth}\label{name-maxwell-braseth-1}

\bookmarksetup{startatroot}

\chapter{Date: 06/22/2026}\label{date-06222026}

\bookmarksetup{startatroot}

\chapter{Module 5 Homework}\label{module-5-homework}

\bookmarksetup{startatroot}

\chapter{```\{r\}}\label{r-4}

require(FNN) require(dplyr) require(ggplot2) require(ranger)
require(rms) require(nnet) library(FNN) library(dplyr) library(ggplot2)
library(ranger) library(rms) library(nnet)

abalone\_data \textless- read.csv(
``https://archive.ics.uci.edu/ml/machine-learning-databases/abalone/abalone.data'',
header = FALSE, col.names =
c(``sex'',``length'',``diameter'',``height'',
``whole\_wt'',``shucked\_wt'',``viscera\_wt'',``shell\_wt'',``rings''))

abalone\_data\(sex <- factor(
  abalone_data\)sex, levels = c(``F'', ``I'', ``M''), labels =
c(``Female'', ``Infant'', ``Male'') )

abalone\_data\(sex <- factor(abalone_data\)sex)

predictors \textless- c(``length'', ``diameter'', ``height'',
``whole\_wt'', ``shucked\_wt'', ``viscera\_wt'', ``shell\_wt'')

\bookmarksetup{startatroot}

\chapter{1.) Fit all five model types to the Abalone data using a 70/30
training-test
split:}\label{fit-all-five-model-types-to-the-abalone-data-using-a-7030-training-test-split}

\bookmarksetup{startatroot}

\chapter{Training-Test Split (70/30)}\label{training-test-split-7030}

set.seed(123)

total\_rows\_in\_abalone\_data \textless- nrow(abalone\_data)

training\_test\_split\_70\_30 \textless- sample(
seq\_len(total\_rows\_in\_abalone\_data), size = floor(0.70 *
total\_rows\_in\_abalone\_data) )

training\_data\_70 \textless-
abalone\_data{[}training\_test\_split\_70\_30, {]} testing\_data\_30
\textless- abalone\_data{[}-training\_test\_split\_70\_30, {]}

\bookmarksetup{startatroot}

\chapter{* KNN (FNN::knn.reg(), choose k using cross-validation or a
reasonable
default)}\label{knn-fnnknn.reg-choose-k-using-cross-validation-or-a-reasonable-default}

y\_training\_data\_rings \textless- training\_data\_70\(rings
y_testing_data_rings  <- testing_data_30\)rings

x\_training\_data \textless- training\_data\_70{[}, predictors{]}
x\_testing\_data \textless- testing\_data\_30{[}, predictors{]}

\bookmarksetup{startatroot}

\chapter{Choosing k values using
cross-validation}\label{choosing-k-values-using-cross-validation}

k\_values\_to\_test \textless- seq(1, 20, by = 1)

cross\_validation\_errors \textless- sapply(k\_values\_to\_test,
function(k) \{ preds \textless- knn.reg(train = x\_training\_data, y =
y\_training\_data\_rings, test = x\_training\_data, k = k)\$pred
mean((preds - y\_training\_data\_rings)\^{}2) \})

\bookmarksetup{startatroot}

\chapter{determines the best k value from how the cross-validation error
values show which k has the smallest distance from the predicted values
to the observed
values}\label{determines-the-best-k-value-from-how-the-cross-validation-error-values-show-which-k-has-the-smallest-distance-from-the-predicted-values-to-the-observed-values}

best\_k \textless-
k\_values\_to\_test{[}min(which(cross\_validation\_errors ==
min(cross\_validation\_errors))){]}

\bookmarksetup{startatroot}

\chapter{Final KNN with best k}\label{final-knn-with-best-k}

final\_knn\_predictions \textless- knn.reg( train = x\_training\_data,
test = x\_testing\_data, y = y\_training\_data\_rings, k = best\_k
)\$pred

\bookmarksetup{startatroot}

\chapter{* Random forest (ranger::ranger(), default
hyperparameters)}\label{random-forest-rangerranger-default-hyperparameters}

random\_forest\_model\_training \textless- ranger( formula = rings
\textasciitilde{} ., data = training\_data\_70, num.trees = 500,
respect.unordered.factors = ``order'' )

random\_forest\_model\_testing \textless- ranger( formula = rings
\textasciitilde{} ., data = testing\_data\_30, num.trees = 500,
respect.unordered.factors = ``order'' )

random\_forest\_training\_predictions \textless-
predict(random\_forest\_model\_training, data =
training\_data\_70)\(predictions
random_forest_testing_predictions <- predict(random_forest_model_testing, data = testing_data_30)\)predictions

\bookmarksetup{startatroot}

\chapter{* Ordinal regression with Restricted Cubic Splines (RCS)
(rms::orm() with rcs()
terms)}\label{ordinal-regression-with-restricted-cubic-splines-rcs-rmsorm-with-rcs-terms}

orm\_fit \textless- orm(rings \textasciitilde{} sex + rcs(length) +
rcs(diameter) + rcs(height) + rcs(whole\_wt) + rcs(shucked\_wt) +
rcs(viscera\_wt) + rcs(shell\_wt), data = training\_data\_70 )

orm\_predictions\_training\_data \textless- predict(orm\_fit, newdata =
training\_data\_70, type = ``mean'') orm\_predictions\_testing\_data
\textless- predict(orm\_fit, newdata = testing\_data\_30, type =
``mean'')

\bookmarksetup{startatroot}

\chapter{* Linear regression with RCS (rms::ols() with rcs()
terms)}\label{linear-regression-with-rcs-rmsols-with-rcs-terms}

dd \textless- datadist(training\_data\_70) options(datadist = `dd')

lm\_fit \textless- ols(rings \textasciitilde{} sex + rcs(length) +
rcs(diameter) + rcs(height) + rcs(whole\_wt) + rcs(shucked\_wt) +
rcs(viscera\_wt) + rcs(shell\_wt), data = training\_data\_70, x = TRUE,
y = TRUE)

lm\_testing\_predictions \textless- predict(lm\_fit, newdata =
testing\_data\_30)

\bookmarksetup{startatroot}

\chapter{* Neural network (nnet::nnet(), linout =
TRUE)}\label{neural-network-nnetnnet-linout-true}

train\_nn \textless- training\_data\_70 test\_nn \textless-
testing\_data\_30

means \textless- sapply(train\_nn{[}, predictors{]}, mean) sds
\textless- sapply(train\_nn{[}, predictors{]}, sd)

train\_nn{[}, predictors{]} \textless- scale(train\_nn{[},
predictors{]}, center = means, scale = sds) test\_nn{[}, predictors{]}
\textless- scale(test\_nn{[}, predictors{]}, center = means, scale =
sds)

x\_train\_nn \textless- model.matrix(rings \textasciitilde{} ., data =
train\_nn){[}, -1{]} x\_test\_nn \textless- model.matrix(rings
\textasciitilde{} ., data = test\_nn){[}, -1{]}

set.seed(123)

nn\_fit \textless- nnet(x = x\_train\_nn, y = train\_nn\$rings, size =
5, linout = TRUE, maxit = 500, trace = FALSE)

nn\_pred \textless- as.vector(predict(nn\_fit, x\_test\_nn))

\bookmarksetup{startatroot}

\chapter{2. For each fitted model, compute on the test
set:}\label{for-each-fitted-model-compute-on-the-test-set}

\bookmarksetup{startatroot}

\chapter{* A calibration plot (predicted vs.~observed, with a loess
smooth and the line of
identity)}\label{a-calibration-plot-predicted-vs.-observed-with-a-loess-smooth-and-the-line-of-identity}

\bookmarksetup{startatroot}

\chapter{* Mean squared error (MSE)}\label{mean-squared-error-mse}

\bookmarksetup{startatroot}

\chapter{KNN}\label{knn}

\bookmarksetup{startatroot}

\chapter{Root Mean Squared Error}\label{root-mean-squared-error}

final\_knn\_prediction\_RMSE\_testing \textless-
sqrt(mean((final\_knn\_predictions - y\_testing\_data\_rings)\^{}2))
final\_knn\_prediction\_RMSE\_training \textless-
sqrt(mean((final\_knn\_predictions - y\_training\_data\_rings)\^{}2))
cat(``Final KNN Root Mean Squared Error (RMSE) for Testing Data ='',
round(final\_knn\_prediction\_RMSE\_testing, 4), ``\n'') cat(``Final KNN
Root Mean Squared Error (RMSE) for Training Data ='',
round(final\_knn\_prediction\_RMSE\_training, 4), ``\n'') \# Calibration
Plot knn\_results \textless- data.frame( Actual =
as.numeric(y\_testing\_data\_rings), Predicted =
as.numeric(final\_knn\_predictions) ) ggplot(knn\_results, aes(x =
Actual, y = Predicted)) + geom\_point(alpha = 0.5, color =
``steelblue'') + geom\_smooth(method = ``loess'', color = ``red'', se =
FALSE) + geom\_abline(slope = 1, intercept = 0, linetype = ``dashed'',
color = ``black'') + labs(title = paste(``KNN Calibration Plot (best k
='', best\_k, ``)''), x = ``Actual Rings or Age of the Abalone Data
Subjects/Individuals'', y = ``Predicted Rings or Age of the Abalone Data
Subjects/Individuals'') + theme\_minimal()

\bookmarksetup{startatroot}

\chapter{Random Forest}\label{random-forest}

\bookmarksetup{startatroot}

\chapter{Root Mean Squared Error
value}\label{root-mean-squared-error-value}

random\_forest\_rmse\_testing \textless-
sqrt(mean((random\_forest\_testing\_predictions -
y\_testing\_data\_rings)\^{}2)) random\_forest\_rmse\_training
\textless- sqrt(mean((random\_forest\_training\_predictions -
y\_training\_data\_rings)\^{}2)) cat(``\nRandom Forest Training RMSE
='', round(random\_forest\_rmse\_training, 4), ``\n'')
cat(``\nRandom Forest Testing RMSE ='',
round(random\_forest\_rmse\_testing, 4), ``\n'') \# calibration plot
rf\_results\_testing \textless- data.frame( Actual\_Rings\_testing =
y\_testing\_data\_rings, Predicted\_testing =
random\_forest\_testing\_predictions ) ggplot(rf\_results\_testing,
aes(x = Actual\_Rings\_testing, y = Predicted\_testing)) +
geom\_point(alpha = 0.5, color = ``steelblue'') + geom\_smooth(method =
``loess'', color = ``red'', se = FALSE, linewidth = 1.2) +
geom\_abline(slope = 1, intercept = 0, linetype = ``dashed'', color =
``black'', linewidth = 1) + labs(title = ``Random Forest Calibration
Plot (Test Set)'', subtitle = paste(``Test RMSE ='',
round(random\_forest\_rmse\_testing, 4)), x = ``Actual Rings'', y =
``Predicted Rings'') + theme\_minimal()

\bookmarksetup{startatroot}

\chapter{Ordinal Regression (ORM +
RCS)}\label{ordinal-regression-orm-rcs}

\bookmarksetup{startatroot}

\chapter{Root Mean Squared Error
value}\label{root-mean-squared-error-value-1}

orm\_rmse\_testing \textless- sqrt(mean((orm\_predictions\_testing\_data
- y\_testing\_data\_rings)\^{}2)) cat(``\nOrdinal Regression Testing
RMSE ='', round(orm\_rmse\_testing, 4), ``\n'') orm\_rmse\_training
\textless- sqrt(mean((orm\_predictions\_training\_data -
y\_training\_data\_rings)\^{}2)) cat(``\nOrdinal Regression Training
RMSE ='', round(orm\_rmse\_training, 4), ``\n'') \# calibration plot
ordinal\_results\_testing \textless- data.frame( Actual =
y\_testing\_data\_rings, Predicted = orm\_predictions\_testing\_data )
ggplot(ordinal\_results\_testing, aes(x = Actual, y = Predicted)) +
geom\_point(alpha = 0.5, color = ``steelblue'') + geom\_smooth(method =
``loess'', color = ``red'', se = FALSE) + geom\_abline(slope = 1,
intercept = 0, linetype = ``dashed'', color = ``black'') + labs(title =
``Ordinal Regression (RCS) Calibration Plot - Testing Data'', x =
``Observed Rings'', y = ``Predicted Rings'') + theme\_minimal()

\bookmarksetup{startatroot}

\chapter{Linear Regression (OLS + RCS)}\label{linear-regression-ols-rcs}

\bookmarksetup{startatroot}

\chapter{Root Mean Squared Error}\label{root-mean-squared-error-1}

lm\_rmse\_testing \textless- sqrt(mean((lm\_testing\_predictions -
y\_testing\_data\_rings)\^{}2)) cat(``Linear Regression (OLS + RCS) Test
RMSE ='', round(lm\_rmse\_testing, 4), ``\n'')

lm\_rmse\_training \textless- sqrt(mean((predict(lm\_fit) -
y\_training\_data\_rings)\^{}2)) cat(``Linear Regression (OLS + RCS)
Training RMSE ='', round(lm\_rmse\_training, 4), ``\n'')

\bookmarksetup{startatroot}

\chapter{calibration plot}\label{calibration-plot}

lm\_results\_testing \textless- data.frame( Actual\_Rings\_testing =
y\_testing\_data\_rings, Predicted\_testing = lm\_testing\_predictions )
ggplot(lm\_results\_testing, aes(x = Actual\_Rings\_testing, y =
Predicted\_testing)) + geom\_point(alpha = 0.5, color = ``steelblue'') +
geom\_smooth(method = ``loess'', color = ``red'', se = FALSE, linewidth
= 1.2) + geom\_abline(slope = 1, intercept = 0, linetype = ``dashed'',
color = ``black'', linewidth = 1) + labs(title = ``Linear Regression
(OLS + RCS) Calibration Plot - Test Set'', subtitle = paste(``Test RMSE
='', round(lm\_rmse\_testing, 4)), x = ``Actual Rings'', y = ``Predicted
Rings'') + theme\_minimal()

\bookmarksetup{startatroot}

\chapter{Neural Network}\label{neural-network}

\bookmarksetup{startatroot}

\chapter{Root Mean Squared Error}\label{root-mean-squared-error-2}

nn\_rmse \textless- sqrt(mean((nn\_pred -
y\_testing\_data\_rings)\^{}2)) cat(``Neural Network RMSE ='',
round(nn\_rmse, 4), ``\n'')

\bookmarksetup{startatroot}

\chapter{calibration plot}\label{calibration-plot-1}

nn\_results\_testing \textless- data.frame( Actual =
y\_testing\_data\_rings, Predicted = nn\_pred )
ggplot(nn\_results\_testing, aes(x = Actual, y = Predicted)) +
geom\_point(alpha = 0.5, color = ``steelblue'') + geom\_smooth(method =
``loess'', color = ``red'', se = FALSE) + geom\_abline(slope = 1,
intercept = 0, linetype = ``dashed'', color = ``black'') + labs(title =
``Neural Network Calibration Plot (Testing Data)'', x = ``Observed
Rings'', y = ``Predicted Rings'') + theme\_minimal()

\bookmarksetup{startatroot}

\chapter{3. Written reflection: Which model performs best on the Abalone
data,}\label{written-reflection-which-model-performs-best-on-the-abalone-data}

\bookmarksetup{startatroot}

\chapter{and which definition of best are you using to decide? What are
some of
the}\label{and-which-definition-of-best-are-you-using-to-decide-what-are-some-of-the}

\bookmarksetup{startatroot}

\chapter{different considerations that might lead to different choices
about which prediction model to
use?}\label{different-considerations-that-might-lead-to-different-choices-about-which-prediction-model-to-use}

\bookmarksetup{startatroot}

\chapter{Given the MSE and Calibration plots values and conclusions, it
was clear to mean and to
others}\label{given-the-mse-and-calibration-plots-values-and-conclusions-it-was-clear-to-mean-and-to-others}

\bookmarksetup{startatroot}

\chapter{that the Random Forest model performed the best in terms of
predicting the age of the Abalone data in cross comparison to the
actual}\label{that-the-random-forest-model-performed-the-best-in-terms-of-predicting-the-age-of-the-abalone-data-in-cross-comparison-to-the-actual}

\bookmarksetup{startatroot}

\chapter{observations of the age of the Abalone witihin the data. This
is evidenced by the fact that the Random Forest model had the
lowest}\label{observations-of-the-age-of-the-abalone-witihin-the-data.-this-is-evidenced-by-the-fact-that-the-random-forest-model-had-the-lowest}

\bookmarksetup{startatroot}

\chapter{RMSE values for both the training and testing data which means
that the predictions were very close to the actual
observations}\label{rmse-values-for-both-the-training-and-testing-data-which-means-that-the-predictions-were-very-close-to-the-actual-observations}

\bookmarksetup{startatroot}

\chapter{in terms of all of the non-linear relationships of the Abalone
rings(age) values with all of the variables values of the data (both
numeric and
non-numeric).}\label{in-terms-of-all-of-the-non-linear-relationships-of-the-abalone-ringsage-values-with-all-of-the-variables-values-of-the-data-both-numeric-and-non-numeric.}

\bookmarksetup{startatroot}

\chapter{This conclusions that the Random Forest model performing the
best in term of
predictions}\label{this-conclusions-that-the-random-forest-model-performing-the-best-in-term-of-predictions}

\bookmarksetup{startatroot}

\chapter{also helped enforce my definition that it was the best model to
use due to how the Random Forest model was able to efficiently predict
the age of the Abalone data and is
very}\label{also-helped-enforce-my-definition-that-it-was-the-best-model-to-use-due-to-how-the-random-forest-model-was-able-to-efficiently-predict-the-age-of-the-abalone-data-and-is-very}

\bookmarksetup{startatroot}

\chapter{scalable and adaptable to new testing data and even with
predicting the outcome values given the actual training data
observations.}\label{scalable-and-adaptable-to-new-testing-data-and-even-with-predicting-the-outcome-values-given-the-actual-training-data-observations.}

\bookmarksetup{startatroot}

\chapter{It is even important to note that the calibration plot of the
Random Forest model showed the predictions were close in distance from
the actual observations of the
data.}\label{it-is-even-important-to-note-that-the-calibration-plot-of-the-random-forest-model-showed-the-predictions-were-close-in-distance-from-the-actual-observations-of-the-data.}

\bookmarksetup{startatroot}

\chapter{Which this meant that the Random Forest model prediction
ring(age) values/counts were very close to the actual observations of
the data and the model's predictions
matched}\label{which-this-meant-that-the-random-forest-model-prediction-ringage-valuescounts-were-very-close-to-the-actual-observations-of-the-data-and-the-models-predictions-matched}

\bookmarksetup{startatroot}

\chapter{the identity line of the calibration plot that means that seems
to fit the model perfectly along with avoiding over and under fitting
when dealing with new
data.}\label{the-identity-line-of-the-calibration-plot-that-means-that-seems-to-fit-the-model-perfectly-along-with-avoiding-over-and-under-fitting-when-dealing-with-new-data.}

\bookmarksetup{startatroot}

\chapter{However, it is important to note that the best model in terms
of accurate predictions would need to be a model that is robust and able
to keep accurate predictions in the face of new testing
data.}\label{however-it-is-important-to-note-that-the-best-model-in-terms-of-accurate-predictions-would-need-to-be-a-model-that-is-robust-and-able-to-keep-accurate-predictions-in-the-face-of-new-testing-data.}

\bookmarksetup{startatroot}

\chapter{In this case of the 70-30 split of the data and considering the
robustness of the models, Random Forest Models would actually not be the
best model since when I altered the data or entered new data the
predictions dramatically
shifted.}\label{in-this-case-of-the-70-30-split-of-the-data-and-considering-the-robustness-of-the-models-random-forest-models-would-actually-not-be-the-best-model-since-when-i-altered-the-data-or-entered-new-data-the-predictions-dramatically-shifted.}

\bookmarksetup{startatroot}

\chapter{As a result of this reality, the next models that I would
consider to be a good model for prediction would be the Oridinal
Regression and the Linear Regression (OLS + RCS)
models.}\label{as-a-result-of-this-reality-the-next-models-that-i-would-consider-to-be-a-good-model-for-prediction-would-be-the-oridinal-regression-and-the-linear-regression-ols-rcs-models.}

\bookmarksetup{startatroot}

\chapter{Within those models, the Oridinal Regression model performed
the best in terms of accurate predictions through its RMSE calculations
and the Linear Regression (OLS + RCS) model performed the best in terms
of interpretability since its predictors
had}\label{within-those-models-the-oridinal-regression-model-performed-the-best-in-terms-of-accurate-predictions-through-its-rmse-calculations-and-the-linear-regression-ols-rcs-model-performed-the-best-in-terms-of-interpretability-since-its-predictors-had}

\bookmarksetup{startatroot}

\chapter{a deep influence on its calibration plot outputs (in other
words its predictor variables gave the expected accurate predictions in
term of the number of rings of the Abalone data and they were accurate
against the training data in the
model).}\label{a-deep-influence-on-its-calibration-plot-outputs-in-other-words-its-predictor-variables-gave-the-expected-accurate-predictions-in-term-of-the-number-of-rings-of-the-abalone-data-and-they-were-accurate-against-the-training-data-in-the-model.}

\bookmarksetup{startatroot}

\chapter{It is also important to note that when I tested these models on
the training data, the Oridinal Regression model performed the best and
showed adaptability when confronted when I ran the test data against
this model which showed a semi-good performance in dealing with over and
underfitting.}\label{it-is-also-important-to-note-that-when-i-tested-these-models-on-the-training-data-the-oridinal-regression-model-performed-the-best-and-showed-adaptability-when-confronted-when-i-ran-the-test-data-against-this-model-which-showed-a-semi-good-performance-in-dealing-with-over-and-underfitting.}

\bookmarksetup{startatroot}

\chapter{At the same time, the Linear Regression (OLS + RCS) model
performed better than the Ordinal Regression Model in terms of its
calibration models which means that the predictions of the regression
actually closely matched the observed ring(age) values/counts of the
Abalone
data.}\label{at-the-same-time-the-linear-regression-ols-rcs-model-performed-better-than-the-ordinal-regression-model-in-terms-of-its-calibration-models-which-means-that-the-predictions-of-the-regression-actually-closely-matched-the-observed-ringage-valuescounts-of-the-abalone-data.}

\bookmarksetup{startatroot}

\chapter{In conclusion, the best ways that I would describe a good
prediction model is that it does need a lower MSE and RMSE amount and
needs a plot such as a calibration plot that reinforces the conclusions
that the predictions ring(age) values/counts are close to the actual
observations of the
data,}\label{in-conclusion-the-best-ways-that-i-would-describe-a-good-prediction-model-is-that-it-does-need-a-lower-mse-and-rmse-amount-and-needs-a-plot-such-as-a-calibration-plot-that-reinforces-the-conclusions-that-the-predictions-ringage-valuescounts-are-close-to-the-actual-observations-of-the-data}

\bookmarksetup{startatroot}

\chapter{while still being adaptable to new testing data and out of
sample data to be considered an efficient and robust prediction
model.}\label{while-still-being-adaptable-to-new-testing-data-and-out-of-sample-data-to-be-considered-an-efficient-and-robust-prediction-model.}

\bookmarksetup{startatroot}

\chapter{```}\label{section-24}

\bookmarksetup{startatroot}

\chapter{Name: Maxwell Braseth}\label{name-maxwell-braseth-2}

\bookmarksetup{startatroot}

\chapter{Date: 06/30/2026}\label{date-06302026}

\bookmarksetup{startatroot}

\chapter{Module 6 Homework}\label{module-6-homework}

\begin{Shaded}
\begin{Highlighting}[]
\CommentTok{\# Loading the insurance.csv data}
\NormalTok{insurance\_data }\OtherTok{\textless{}{-}} \FunctionTok{read.csv}\NormalTok{(}\StringTok{"data/insurance.csv"}\NormalTok{, }\AttributeTok{encoding =} \StringTok{"UTF{-}8"}\NormalTok{)}
\FunctionTok{library}\NormalTok{(MASS)}
\FunctionTok{library}\NormalTok{(rms)}
\end{Highlighting}
\end{Shaded}

\begin{verbatim}
Loading required package: Hmisc
\end{verbatim}

\begin{verbatim}

Attaching package: 'Hmisc'
\end{verbatim}

\begin{verbatim}
The following objects are masked from 'package:base':

    format.pval, units
\end{verbatim}

\begin{Shaded}
\begin{Highlighting}[]
\CommentTok{\# 1.) Write down the linear regression model }
\CommentTok{\# mathematically. Use matrix notation Y = Xβ + ε.}
\CommentTok{\# Identify each component: what is Y, what }
\CommentTok{\# columns are in X, what does each element of β}
\CommentTok{\# represent?}

\CommentTok{\# Y = Xβ + ε}
\CommentTok{\# Y = dependent variable {-}{-}\textgreater{} response or obsered dependent response variable}
\CommentTok{\# X = independent variable {-}{-}\textgreater{} Design matrix}
\CommentTok{\# β = coefficient {-}{-}\textgreater{} (X\^{}T X)\^{}{-}1 X\^{}T Y}
\CommentTok{\# X\^{}T = transpose of the matrix}
\CommentTok{\# (X\^{}T X)\^{}{-}1 {-}{-}\textgreater{} inverse of the matrix product X\^{}T * X}
\CommentTok{\#   * example of an inverse of a matrix: a\^{}{-}1 = 1/a {-}{-}\textgreater{} a = [a, b, c, d] and a\^{}{-}1 = [1/a, 1/b, 1/c, 1/d] or A\^{}{-}1 = 1/|A| * Adj(A) = 1/(ad{-}bc) * [d, {-}b, {-}c, a]}
\CommentTok{\# ε = error {-}{-} \textgreater{} \# ε = Y {-} Xβ = ŷ {-} y}
\CommentTok{\# ŷ = predicted value of y = Xβ}

\CommentTok{\# Set factor levels}
\NormalTok{insurance\_data}\SpecialCharTok{$}\NormalTok{sex    }\OtherTok{\textless{}{-}} \FunctionTok{factor}\NormalTok{(insurance\_data}\SpecialCharTok{$}\NormalTok{sex,    }\AttributeTok{levels =} \FunctionTok{c}\NormalTok{(}\StringTok{"female"}\NormalTok{, }\StringTok{"male"}\NormalTok{))}
\NormalTok{insurance\_data}\SpecialCharTok{$}\NormalTok{region }\OtherTok{\textless{}{-}} \FunctionTok{factor}\NormalTok{(insurance\_data}\SpecialCharTok{$}\NormalTok{region, }\AttributeTok{levels =} \FunctionTok{c}\NormalTok{(}\StringTok{"northeast"}\NormalTok{, }\StringTok{"northwest"}\NormalTok{, }\StringTok{"southeast"}\NormalTok{, }\StringTok{"southwest"}\NormalTok{))}
\NormalTok{insurance\_data}\SpecialCharTok{$}\NormalTok{smoker }\OtherTok{\textless{}{-}} \FunctionTok{factor}\NormalTok{(insurance\_data}\SpecialCharTok{$}\NormalTok{smoker, }\AttributeTok{levels =} \FunctionTok{c}\NormalTok{(}\StringTok{"no"}\NormalTok{, }\StringTok{"yes"}\NormalTok{))}

\CommentTok{\# creating design matrices to combine together all of the independent predictor variables of x}
\NormalTok{sex\_matrix    }\OtherTok{\textless{}{-}} \FunctionTok{model.matrix}\NormalTok{(}\SpecialCharTok{\textasciitilde{}} \DecValTok{0} \SpecialCharTok{+}\NormalTok{ sex,    }\AttributeTok{data =}\NormalTok{ insurance\_data)}
\NormalTok{region\_matrix }\OtherTok{\textless{}{-}} \FunctionTok{model.matrix}\NormalTok{(}\SpecialCharTok{\textasciitilde{}} \DecValTok{0} \SpecialCharTok{+}\NormalTok{ region, }\AttributeTok{data =}\NormalTok{ insurance\_data)}
\NormalTok{smoker\_matrix }\OtherTok{\textless{}{-}} \FunctionTok{model.matrix}\NormalTok{(}\SpecialCharTok{\textasciitilde{}} \DecValTok{0} \SpecialCharTok{+}\NormalTok{ smoker, }\AttributeTok{data =}\NormalTok{ insurance\_data)}
\NormalTok{age\_matrix     }\OtherTok{\textless{}{-}} \FunctionTok{as.matrix}\NormalTok{(insurance\_data}\SpecialCharTok{$}\NormalTok{age)}
\NormalTok{bmi\_matrix     }\OtherTok{\textless{}{-}} \FunctionTok{as.matrix}\NormalTok{(insurance\_data}\SpecialCharTok{$}\NormalTok{bmi)}
\NormalTok{children\_matrix }\OtherTok{\textless{}{-}} \FunctionTok{as.matrix}\NormalTok{(insurance\_data}\SpecialCharTok{$}\NormalTok{children)}

\CommentTok{\# Combine all of the x variables within matrices format into full design matrix}
\CommentTok{\# what are columns are in X and what does they represent?}
\CommentTok{\# These x variables are the independent predictor variables that influences the linear regression relationship of showing how }
\CommentTok{\# the features of individuals (such as age, sex, bmi, number of children, etc...) affect the cost of health insurance charges in a positive or negative way.}
\NormalTok{x }\OtherTok{\textless{}{-}} \FunctionTok{cbind}\NormalTok{(}
\NormalTok{  sex\_matrix, }
\NormalTok{  region\_matrix, }
\NormalTok{  smoker\_matrix, }
\NormalTok{  age\_matrix, }
\NormalTok{  bmi\_matrix, }
\NormalTok{  children\_matrix}
\NormalTok{)}

\CommentTok{\# Give clear column names}
\FunctionTok{colnames}\NormalTok{(x) }\OtherTok{\textless{}{-}} \FunctionTok{c}\NormalTok{(}
  \StringTok{"sexfemale"}\NormalTok{, }\StringTok{"sexmale"}\NormalTok{,}
  \StringTok{"regionnortheast"}\NormalTok{, }\StringTok{"regionnorthwest"}\NormalTok{, }\StringTok{"regionsoutheast"}\NormalTok{, }\StringTok{"regionsouthwest"}\NormalTok{,}
  \StringTok{"smokerno"}\NormalTok{, }\StringTok{"smokeryes"}\NormalTok{,}
  \StringTok{"age"}\NormalTok{, }\StringTok{"bmi"}\NormalTok{, }\StringTok{"children"}
\NormalTok{)}
\CommentTok{\# Y dependent variable for all x independent variables}
\NormalTok{y }\OtherTok{\textless{}{-}}\NormalTok{ insurance\_data}\SpecialCharTok{$}\NormalTok{charges}

\CommentTok{\# Use generalized inverse for the Coefficient Estimate equation (X\^{}T * X)\^{}{-}1 * X\^{}T * Y}
\CommentTok{\# (X\^{}T X)\^{}{-}1 X\^{}T Y}
\NormalTok{b }\OtherTok{\textless{}{-}} \FunctionTok{ginv}\NormalTok{(}\FunctionTok{t}\NormalTok{(x) }\SpecialCharTok{\%*\%}\NormalTok{ x) }\SpecialCharTok{\%*\%} \FunctionTok{t}\NormalTok{(x) }\SpecialCharTok{\%*\%}\NormalTok{ y}

\CommentTok{\# assigning the column and row names}
\FunctionTok{rownames}\NormalTok{(b) }\OtherTok{\textless{}{-}} \FunctionTok{colnames}\NormalTok{(x)}
\FunctionTok{colnames}\NormalTok{(b) }\OtherTok{\textless{}{-}} \StringTok{"β Coefficient Estimate"}


\CommentTok{\# For all of the columns are in X, what does each element of β represent?}

\CommentTok{\# For each element of β in combination with the columns of x (age, bmi, children, etc...),}
\CommentTok{\# the element of β tells us how much that particular column in X contributes and influences the calculation of the predicted value of Y.}
\CommentTok{\# In this case though, we are using the linear combinations of the columns describing the features of the individuals (age, bmi, children, etc...)}
\CommentTok{\# along with the elements of β (to increase or decrease weight or influence) to calculate the predicted values of y or in this case the charges of health insurance.}
\CommentTok{\# In otherwords, the linear combinations of the columns of x and the elements of β helps determine which features of the individuals has more or less influence on the predicted insurance charges for these individuals.}


\CommentTok{\# ŷ = Xβ}
\NormalTok{predicted\_y }\OtherTok{\textless{}{-}}\NormalTok{ x }\SpecialCharTok{\%*\%}\NormalTok{ b}

\CommentTok{\# ε = Y {-} Xβ = ŷ {-} y}
\NormalTok{e1 }\OtherTok{\textless{}{-}}\NormalTok{ y }\SpecialCharTok{{-}}\NormalTok{ predicted\_y}
\NormalTok{e2 }\OtherTok{\textless{}{-}}\NormalTok{ y }\SpecialCharTok{{-}}\NormalTok{ x }\SpecialCharTok{\%*\%}\NormalTok{ b}
\FunctionTok{print}\NormalTok{(}\StringTok{"error 1 vs error 2 equal ?: "}\NormalTok{)}
\end{Highlighting}
\end{Shaded}

\begin{verbatim}
[1] "error 1 vs error 2 equal ?: "
\end{verbatim}

\begin{Shaded}
\begin{Highlighting}[]
\FunctionTok{print}\NormalTok{(}\FunctionTok{head}\NormalTok{(}\FunctionTok{mean}\NormalTok{(e1)) }\SpecialCharTok{==}\NormalTok{ (}\FunctionTok{head}\NormalTok{(}\FunctionTok{mean}\NormalTok{(e2))))}
\end{Highlighting}
\end{Shaded}

\begin{verbatim}
[1] TRUE
\end{verbatim}

\begin{Shaded}
\begin{Highlighting}[]
\CommentTok{\# What is Y in this scenario?}
\CommentTok{\# the Y dependent variable is the observed dependent response variable that we are trying to predict or estimate.}
\CommentTok{\# In this case, the Y dependent variable is the cost of health insurance charges that is dependent on the variables/features of the individuals (age, sex, bmi, number of children, etc...)}
\CommentTok{\# along with determining how the linear combinations of the coefficient elements (β) helps indicate and determine the impact/weight from which these individuals features of x (age, sex, bmi, number of children, etc...)}
\CommentTok{\# actually has on specific predicted calculation costs of health insurance charges of these individuals.}

\CommentTok{\# linear regression model : Y = Xβ + ε}
\NormalTok{linear\_regression\_model\_y }\OtherTok{\textless{}{-}}\NormalTok{ (x }\SpecialCharTok{\%*\%}\NormalTok{ b) }\SpecialCharTok{+}\NormalTok{ e1}

\CommentTok{\# Show breakdown for first 5 Linear Regression Model values along with the predicted y and epison value ε}
\NormalTok{breakdown\_of\_linear\_regression\_model }\OtherTok{\textless{}{-}} \FunctionTok{data.frame}\NormalTok{(}
  \AttributeTok{Observed\_Y\_Linear\_Regression\_value =} \FunctionTok{round}\NormalTok{(linear\_regression\_model\_y, }\DecValTok{2}\NormalTok{),}
  \AttributeTok{Predicted\_Y\_value =} \FunctionTok{round}\NormalTok{(predicted\_y, }\DecValTok{2}\NormalTok{),}
\NormalTok{  Error\_ε}\AttributeTok{\_value =} \FunctionTok{round}\NormalTok{(e1, }\DecValTok{2}\NormalTok{)}
\NormalTok{)}
\NormalTok{coefficient\_breakdown }\OtherTok{\textless{}{-}} \FunctionTok{data.frame}\NormalTok{(}
  \AttributeTok{Predictor\_Feature\_of\_Individuals\_X =} \FunctionTok{rownames}\NormalTok{(b),}
\NormalTok{  Coefficient\_β}\AttributeTok{\_value =} \FunctionTok{round}\NormalTok{(b[, }\DecValTok{1}\NormalTok{], }\DecValTok{2}\NormalTok{)}
\NormalTok{)}


\FunctionTok{colnames}\NormalTok{(breakdown\_of\_linear\_regression\_model) }\OtherTok{\textless{}{-}} \FunctionTok{c}\NormalTok{(}
  \StringTok{"Observed\_Y\_Linear\_Regression\_For\_Insurance\_Charges"}\NormalTok{,}
  \StringTok{"Predicted\_Y\_value"}\NormalTok{,}
  \StringTok{"Error\_ε\_value"}
\NormalTok{)}

\FunctionTok{print}\NormalTok{(}\StringTok{"Linear Regression Model for Insurance Charges for Individuals:"}\NormalTok{)}
\end{Highlighting}
\end{Shaded}

\begin{verbatim}
[1] "Linear Regression Model for Insurance Charges for Individuals:"
\end{verbatim}

\begin{Shaded}
\begin{Highlighting}[]
\FunctionTok{print}\NormalTok{(}\FunctionTok{head}\NormalTok{(breakdown\_of\_linear\_regression\_model, }\DecValTok{10}\NormalTok{))}
\end{Highlighting}
\end{Shaded}

\begin{verbatim}
   Observed_Y_Linear_Regression_For_Insurance_Charges Predicted_Y_value
1                                            16884.92          25293.71
2                                             1725.55           3448.60
3                                             4449.46           6706.99
4                                            21984.47           3754.83
5                                             3866.86           5592.49
6                                             3756.62           3719.83
7                                             8240.59          10659.96
8                                             7281.51           8047.91
9                                             6406.41           8502.97
10                                           28923.14          11884.64
   Error_ε_value
1       -8408.79
2       -1723.05
3       -2257.53
4       18229.64
5       -1725.64
6          36.80
7       -2419.37
8        -766.41
9       -2096.56
10      17038.50
\end{verbatim}

\begin{Shaded}
\begin{Highlighting}[]
\FunctionTok{print}\NormalTok{(}\StringTok{"β Coefficient Estimates by Feature of Individuals (X):"}\NormalTok{)}
\end{Highlighting}
\end{Shaded}

\begin{verbatim}
[1] "β Coefficient Estimates by Feature of Individuals (X):"
\end{verbatim}

\begin{Shaded}
\begin{Highlighting}[]
\FunctionTok{print}\NormalTok{(coefficient\_breakdown)}
\end{Highlighting}
\end{Shaded}

\begin{verbatim}
                Predictor_Feature_of_Individuals_X Coefficient_β_value
sexfemale                                sexfemale             -201.12
sexmale                                    sexmale             -332.43
regionnortheast                    regionnortheast              453.62
regionnorthwest                    regionnorthwest              100.66
regionsoutheast                    regionsoutheast             -581.40
regionsouthwest                    regionsouthwest             -506.43
smokerno                                  smokerno           -12191.04
smokeryes                                smokeryes            11657.49
age                                            age              256.86
bmi                                            bmi              339.19
children                                  children              475.50
\end{verbatim}

\begin{Shaded}
\begin{Highlighting}[]
\CommentTok{\# 2. Explain in words how the model answers the research question.}
\CommentTok{\# Which coefficient directly addresses whether smoking is associated with }
\CommentTok{\# higher charges, and what does its magnitude represent?}

\NormalTok{smokeryes\_coef }\OtherTok{\textless{}{-}}\NormalTok{ b[}\StringTok{"smokeryes"}\NormalTok{, ]}
\NormalTok{bmi\_coef }\OtherTok{\textless{}{-}}\NormalTok{ b[}\StringTok{"bmi"}\NormalTok{, ]}
\NormalTok{age\_coef }\OtherTok{\textless{}{-}}\NormalTok{ b[}\StringTok{"age"}\NormalTok{, ]}
\NormalTok{children\_coef }\OtherTok{\textless{}{-}}\NormalTok{ b[}\StringTok{"children"}\NormalTok{, ]}

\NormalTok{max\_values\_df\_smokeryes }\OtherTok{\textless{}{-}} \FunctionTok{data.frame}\NormalTok{(}
  \AttributeTok{Item =} \FunctionTok{c}\NormalTok{(}
    \StringTok{"Max\_SmokerYes\_Coefficient\_Overall"}\NormalTok{,}
    \StringTok{"Max\_Observed\_Y\_Linear\_Regression\_value"}\NormalTok{,}
    \StringTok{"Max\_Predicted\_Y\_value"}
\NormalTok{  ),}
  \AttributeTok{Value =} \FunctionTok{c}\NormalTok{(}
    \FunctionTok{max}\NormalTok{(smokeryes\_coef),}
    \FunctionTok{max}\NormalTok{(y),}
    \FunctionTok{max}\NormalTok{(predicted\_y)}
\NormalTok{  )}
\NormalTok{)}
\NormalTok{max\_values\_df\_bmi }\OtherTok{\textless{}{-}} \FunctionTok{data.frame}\NormalTok{(}
  \AttributeTok{Item =} \FunctionTok{c}\NormalTok{(}
    \StringTok{"Max\_BMI\_Coefficient\_Overall"}\NormalTok{,}
    \StringTok{"Max\_Observed\_Y\_Linear\_Regression\_value"}\NormalTok{,}
    \StringTok{"Max\_Predicted\_Y\_value"}
\NormalTok{  ),}
  \AttributeTok{Value =} \FunctionTok{c}\NormalTok{(}
    \FunctionTok{max}\NormalTok{(bmi\_coef),}
    \FunctionTok{max}\NormalTok{(y),}
    \FunctionTok{max}\NormalTok{(predicted\_y)}
\NormalTok{  )}
\NormalTok{)}
\NormalTok{max\_values\_df\_age }\OtherTok{\textless{}{-}} \FunctionTok{data.frame}\NormalTok{(}
  \AttributeTok{Item =} \FunctionTok{c}\NormalTok{(}
    \StringTok{"Max\_Age\_Coefficient\_Overall"}\NormalTok{,}
    \StringTok{"Max\_Observed\_Y\_Linear\_Regression\_value"}\NormalTok{,}
    \StringTok{"Max\_Predicted\_Y\_value"}
\NormalTok{  ),}
  \AttributeTok{Value =} \FunctionTok{c}\NormalTok{(}
    \FunctionTok{max}\NormalTok{(age\_coef),}
    \FunctionTok{max}\NormalTok{(y),}
    \FunctionTok{max}\NormalTok{(predicted\_y)}
\NormalTok{  )}
\NormalTok{)}
\NormalTok{max\_values\_df\_children }\OtherTok{\textless{}{-}} \FunctionTok{data.frame}\NormalTok{(}
  \AttributeTok{Item =} \FunctionTok{c}\NormalTok{(}
    \StringTok{"Max\_children\_Coefficient\_Overall"}\NormalTok{,}
    \StringTok{"Max\_Observed\_Y\_Linear\_Regression\_value"}\NormalTok{,}
    \StringTok{"Max\_Predicted\_Y\_value"}
\NormalTok{  ),}
  \AttributeTok{Value =} \FunctionTok{c}\NormalTok{(}
    \FunctionTok{max}\NormalTok{(children\_coef),}
    \FunctionTok{max}\NormalTok{(y),}
    \FunctionTok{max}\NormalTok{(predicted\_y)}
\NormalTok{  )}
\NormalTok{)}

\FunctionTok{print}\NormalTok{(max\_values\_df\_smokeryes)}
\end{Highlighting}
\end{Shaded}

\begin{verbatim}
                                    Item    Value
1      Max_SmokerYes_Coefficient_Overall 11657.49
2 Max_Observed_Y_Linear_Regression_value 63770.43
3                  Max_Predicted_Y_value 40834.81
\end{verbatim}

\begin{Shaded}
\begin{Highlighting}[]
\FunctionTok{print}\NormalTok{(max\_values\_df\_bmi)}
\end{Highlighting}
\end{Shaded}

\begin{verbatim}
                                    Item      Value
1            Max_BMI_Coefficient_Overall   339.1935
2 Max_Observed_Y_Linear_Regression_value 63770.4280
3                  Max_Predicted_Y_value 40834.8087
\end{verbatim}

\begin{Shaded}
\begin{Highlighting}[]
\FunctionTok{print}\NormalTok{(max\_values\_df\_age)}
\end{Highlighting}
\end{Shaded}

\begin{verbatim}
                                    Item      Value
1            Max_Age_Coefficient_Overall   256.8564
2 Max_Observed_Y_Linear_Regression_value 63770.4280
3                  Max_Predicted_Y_value 40834.8087
\end{verbatim}

\begin{Shaded}
\begin{Highlighting}[]
\FunctionTok{print}\NormalTok{(max\_values\_df\_children)}
\end{Highlighting}
\end{Shaded}

\begin{verbatim}
                                    Item      Value
1       Max_children_Coefficient_Overall   475.5005
2 Max_Observed_Y_Linear_Regression_value 63770.4280
3                  Max_Predicted_Y_value 40834.8087
\end{verbatim}

\begin{Shaded}
\begin{Highlighting}[]
\CommentTok{\# After some evaluations of the model along with the summary code above, this model definitely addresses the research question,}
\CommentTok{\# which is whether smoking is associated with higher charges, and what does its magnitude represent.}
\CommentTok{\# The highest coefficient β estimate for Smoker Yes is 11657.67, which when compared with the other coefficients β estimates for the remaining x features/variables}
\CommentTok{\# such as BMI, Age, and Number of Children (339.1935, 256.8564, 475.5005), which indicates that smoking is associated with higher coefficient value which in turn shows that}
\CommentTok{\# smoking is associated with the prediction of the highest or greater cost of health insurance charges of $63770.4280.}

\CommentTok{\# Finally in terms of its magnitude representation, the highest coefficient β estimate for Smoker Yes }
\CommentTok{\# which is 11657.67 basically indicates that the feature of smoking is one of the most influential/impactful}
\CommentTok{\# feature that contributes to predicting the insurance charges of the individuals in this data set. Altogether,}
\CommentTok{\# any of the highest coefficient β estimates for any feature of the individuals is heavily influential in the predicting of the health insurance charges of the individuals of this data set.}


\CommentTok{\# 3. Fit the model in R using lm(). Report the estimated coefficients, }
\CommentTok{\# their standard errors, and a confidence interval for the smoking coefficient. }
\CommentTok{\# Write one sentence interpreting the smoking coefficient in the context of the research question.}

\NormalTok{lm\_fit }\OtherTok{\textless{}{-}} \FunctionTok{lm}\NormalTok{(charges }\SpecialCharTok{\textasciitilde{}}\NormalTok{ age }\SpecialCharTok{+}\NormalTok{ bmi }\SpecialCharTok{+}\NormalTok{ children }\SpecialCharTok{+}\NormalTok{ sex }\SpecialCharTok{+}\NormalTok{ smoker }\SpecialCharTok{+}\NormalTok{ region,}
             \AttributeTok{data =}\NormalTok{ insurance\_data)}

\NormalTok{coefficient\_breakdown\_lm }\OtherTok{\textless{}{-}} \FunctionTok{data.frame}\NormalTok{(}
  \AttributeTok{Predictor =} \FunctionTok{rownames}\NormalTok{(}\FunctionTok{coef}\NormalTok{(}\FunctionTok{summary}\NormalTok{(lm\_fit))),}
  \AttributeTok{Estimate =} \FunctionTok{coef}\NormalTok{(}\FunctionTok{summary}\NormalTok{(lm\_fit))[,}\DecValTok{1}\NormalTok{],}
  \AttributeTok{Std\_Error =} \FunctionTok{coef}\NormalTok{(}\FunctionTok{summary}\NormalTok{(lm\_fit))[,}\DecValTok{2}\NormalTok{],}
  \AttributeTok{stringsAsFactors =} \ConstantTok{FALSE}
\NormalTok{)}

\FunctionTok{cat}\NormalTok{(}\StringTok{"====β coefficient estimates and standard errors for Smoking (lm version)====}\SpecialCharTok{\textbackslash{}n}\StringTok{"}\NormalTok{)}
\end{Highlighting}
\end{Shaded}

\begin{verbatim}
====β coefficient estimates and standard errors for Smoking (lm version)====
\end{verbatim}

\begin{Shaded}
\begin{Highlighting}[]
\FunctionTok{print}\NormalTok{(coefficient\_breakdown\_lm[}\StringTok{"smokeryes"}\NormalTok{, }\DecValTok{2}\SpecialCharTok{:}\DecValTok{3}\NormalTok{])}
\end{Highlighting}
\end{Shaded}

\begin{verbatim}
          Estimate Std_Error
smokeryes 23848.53  413.1534
\end{verbatim}

\begin{Shaded}
\begin{Highlighting}[]
\FunctionTok{cat}\NormalTok{(}\StringTok{"====95\% CI for Smoking====}\SpecialCharTok{\textbackslash{}n}\StringTok{"}\NormalTok{)}
\end{Highlighting}
\end{Shaded}

\begin{verbatim}
====95% CI for Smoking====
\end{verbatim}

\begin{Shaded}
\begin{Highlighting}[]
\NormalTok{smoking\_ci }\OtherTok{\textless{}{-}} \FunctionTok{confint}\NormalTok{(lm\_fit, }\StringTok{"smokeryes"}\NormalTok{, }\AttributeTok{level =} \FloatTok{0.95}\NormalTok{)}
\FunctionTok{print}\NormalTok{(}\FunctionTok{round}\NormalTok{(smoking\_ci,}\DecValTok{4}\NormalTok{))}
\end{Highlighting}
\end{Shaded}

\begin{verbatim}
             2.5 %   97.5 %
smokeryes 23038.03 24659.04
\end{verbatim}

\begin{Shaded}
\begin{Highlighting}[]
\CommentTok{\# The best one sentence interpretation of the smoking coefficient in the context of the research question of "which is whether smoking is associated with higher charges and what does its magnitude represent" is the following.}
\CommentTok{\# From the lm() fitted model, the smoker\textquotesingle{}s coefficient β estimate, the smoker\textquotesingle{}s standard error along with the Confidence Interval of 95\% all having high values and these high values all signals that smoking is highly associated }
\CommentTok{\# with how these values helps predict the cost of health insurance charges of the individuals in this data set just like the answer given in question 2.}


\CommentTok{\# 4. Refit the model using rms::ols() instead of lm(). Compare the coefficient table from both fits. }
\CommentTok{\# Then compute a 95\% confidence interval for the smoking coefficient from lm() (using confint()) }
\CommentTok{\# and from ols() (using summary()). Are the intervals identical? }
\CommentTok{\#What additional output does ols() provide that lm() does not?}

\CommentTok{\# datadist setup}
\NormalTok{dd }\OtherTok{\textless{}{-}} \FunctionTok{datadist}\NormalTok{(insurance\_data)}
\FunctionTok{options}\NormalTok{(}\AttributeTok{datadist =} \StringTok{"dd"}\NormalTok{)}

\CommentTok{\# Fitting ols model}
\NormalTok{ols\_fit }\OtherTok{\textless{}{-}} \FunctionTok{ols}\NormalTok{(charges }\SpecialCharTok{\textasciitilde{}}\NormalTok{ age }\SpecialCharTok{+}\NormalTok{ bmi }\SpecialCharTok{+}\NormalTok{ children }\SpecialCharTok{+}\NormalTok{ sex }\SpecialCharTok{+}\NormalTok{ smoker }\SpecialCharTok{+}\NormalTok{ region, }
               \AttributeTok{data =}\NormalTok{ insurance\_data, }\AttributeTok{x =} \ConstantTok{TRUE}\NormalTok{, }\AttributeTok{y =} \ConstantTok{TRUE}\NormalTok{)}

\CommentTok{\# Coefficient tables for both lm() and ols()}
\FunctionTok{cat}\NormalTok{(}\StringTok{"}\SpecialCharTok{\textbackslash{}n}\StringTok{========== lm() Coefficients ==========}\SpecialCharTok{\textbackslash{}n}\StringTok{"}\NormalTok{)}
\end{Highlighting}
\end{Shaded}

\begin{verbatim}

========== lm() Coefficients ==========
\end{verbatim}

\begin{Shaded}
\begin{Highlighting}[]
\NormalTok{lm\_coef }\OtherTok{\textless{}{-}} \FunctionTok{summary}\NormalTok{(lm\_fit)}\SpecialCharTok{$}\NormalTok{coefficients}
\FunctionTok{print}\NormalTok{(}\FunctionTok{round}\NormalTok{(lm\_coef, }\DecValTok{4}\NormalTok{))}
\end{Highlighting}
\end{Shaded}

\begin{verbatim}
                   Estimate Std. Error  t value Pr(>|t|)
(Intercept)     -11938.5386   987.8192 -12.0858   0.0000
age                256.8564    11.8988  21.5867   0.0000
bmi                339.1935    28.5995  11.8601   0.0000
children           475.5005   137.8041   3.4506   0.0006
sexmale           -131.3144   332.9454  -0.3944   0.6933
smokeryes        23848.5345   413.1534  57.7232   0.0000
regionnorthwest   -352.9639   476.2758  -0.7411   0.4588
regionsoutheast  -1035.0220   478.6922  -2.1622   0.0308
regionsouthwest   -960.0510   477.9330  -2.0088   0.0448
\end{verbatim}

\begin{Shaded}
\begin{Highlighting}[]
\FunctionTok{cat}\NormalTok{(}\StringTok{"}\SpecialCharTok{\textbackslash{}n}\StringTok{========== rms::ols() Summary ==========}\SpecialCharTok{\textbackslash{}n}\StringTok{"}\NormalTok{)}
\end{Highlighting}
\end{Shaded}

\begin{verbatim}

========== rms::ols() Summary ==========
\end{verbatim}

\begin{Shaded}
\begin{Highlighting}[]
\NormalTok{ols\_summary }\OtherTok{\textless{}{-}} \FunctionTok{summary}\NormalTok{(ols\_fit)}
\NormalTok{ols\_df }\OtherTok{\textless{}{-}} \FunctionTok{as.data.frame}\NormalTok{(ols\_summary)}
\FunctionTok{print}\NormalTok{(ols\_df)}
\end{Highlighting}
\end{Shaded}

\begin{verbatim}
                                  Low     High   Diff.      Effect     S.E.
age                          27.00000 51.00000 24.0000  6164.55246 285.5724
bmi                          26.29625 34.69375  8.3975  2848.37703 240.1641
children                      0.00000  2.00000  2.0000   951.00109 275.6082
sex - female:male             2.00000  1.00000      NA   131.31436 332.9454
smoker - yes:no               1.00000  2.00000      NA 23848.53454 413.1534
region - northeast:southeast  3.00000  1.00000      NA  1035.02205 478.6922
region - northwest:southeast  3.00000  2.00000      NA   682.05815 478.9592
region - southwest:southeast  3.00000  4.00000      NA    74.97106 470.6386
                              Lower 0.95 Upper 0.95 Type
age                           5604.33068  6724.7742    1
bmi                           2377.23505  3319.5190    1
children                       410.32657  1491.6756    1
sex - female:male             -521.84155   784.4703    1
smoker - yes:no              23038.03071 24659.0384    1
region - northeast:southeast    95.94733  1974.0968    1
region - northwest:southeast  -257.54026  1621.6566    1
region - southwest:southeast  -848.30457   998.2467    1
\end{verbatim}

\begin{Shaded}
\begin{Highlighting}[]
\FunctionTok{cat}\NormalTok{(}\StringTok{"}\SpecialCharTok{\textbackslash{}n}\StringTok{========== Equality Testing for smoker coefficient and 95\% CI ==========}\SpecialCharTok{\textbackslash{}n}\StringTok{"}\NormalTok{)}
\end{Highlighting}
\end{Shaded}

\begin{verbatim}

========== Equality Testing for smoker coefficient and 95% CI ==========
\end{verbatim}

\begin{Shaded}
\begin{Highlighting}[]
\CommentTok{\# Find smoker for both lm() and ols()}
\NormalTok{lm\_smoker\_name  }\OtherTok{\textless{}{-}} \FunctionTok{grep}\NormalTok{(}\StringTok{"smoker"}\NormalTok{, }\FunctionTok{rownames}\NormalTok{(lm\_coef), }\AttributeTok{ignore.case =} \ConstantTok{TRUE}\NormalTok{, }\AttributeTok{value =} \ConstantTok{TRUE}\NormalTok{)}
\NormalTok{ols\_smoker\_name }\OtherTok{\textless{}{-}} \FunctionTok{grep}\NormalTok{(}\StringTok{"smoker"}\NormalTok{, }\FunctionTok{rownames}\NormalTok{(ols\_df),   }\AttributeTok{ignore.case =} \ConstantTok{TRUE}\NormalTok{, }\AttributeTok{value =} \ConstantTok{TRUE}\NormalTok{)}

\CommentTok{\# Getting the smoker coefficient β}
\NormalTok{lm\_est\_coefficient   }\OtherTok{\textless{}{-}}\NormalTok{ lm\_coef[lm\_smoker\_name, }\StringTok{"Estimate"}\NormalTok{]}
\NormalTok{ols\_eff\_coefficient  }\OtherTok{\textless{}{-}}\NormalTok{ ols\_df[ols\_smoker\_name, }\StringTok{"Effect"}\NormalTok{]}

\NormalTok{lm\_ci    }\OtherTok{\textless{}{-}} \FunctionTok{confint}\NormalTok{(lm\_fit, lm\_smoker\_name, }\AttributeTok{level =} \FloatTok{0.95}\NormalTok{)}
\NormalTok{ols\_ci\_lower }\OtherTok{\textless{}{-}}\NormalTok{ ols\_df[ols\_smoker\_name, }\StringTok{"Lower 0.95"}\NormalTok{]}
\NormalTok{ols\_ci\_upper }\OtherTok{\textless{}{-}}\NormalTok{ ols\_df[ols\_smoker\_name, }\StringTok{"Upper 0.95"}\NormalTok{]}

\CommentTok{\# Tests}
\FunctionTok{cat}\NormalTok{(}\StringTok{"Smoker coefficient comparison:}\SpecialCharTok{\textbackslash{}n}\StringTok{"}\NormalTok{)}
\end{Highlighting}
\end{Shaded}

\begin{verbatim}
Smoker coefficient comparison:
\end{verbatim}

\begin{Shaded}
\begin{Highlighting}[]
\FunctionTok{cat}\NormalTok{(}\FunctionTok{sprintf}\NormalTok{(}\StringTok{"lm()  Estimate : \%.6f}\SpecialCharTok{\textbackslash{}n}\StringTok{"}\NormalTok{, lm\_est\_coefficient))}
\end{Highlighting}
\end{Shaded}

\begin{verbatim}
lm()  Estimate : 23848.534542
\end{verbatim}

\begin{Shaded}
\begin{Highlighting}[]
\FunctionTok{cat}\NormalTok{(}\FunctionTok{sprintf}\NormalTok{(}\StringTok{"ols() Effect   : \%.6f}\SpecialCharTok{\textbackslash{}n}\StringTok{"}\NormalTok{, ols\_eff\_coefficient))}
\end{Highlighting}
\end{Shaded}

\begin{verbatim}
ols() Effect   : 23848.534542
\end{verbatim}

\begin{Shaded}
\begin{Highlighting}[]
\FunctionTok{cat}\NormalTok{(}\StringTok{"Are they equal?   "}\NormalTok{, }\FunctionTok{all.equal}\NormalTok{(lm\_est\_coefficient, ols\_eff\_coefficient, }\AttributeTok{tolerance =} \FloatTok{1e{-}8}\NormalTok{), }\StringTok{"}\SpecialCharTok{\textbackslash{}n\textbackslash{}n}\StringTok{"}\NormalTok{)}
\end{Highlighting}
\end{Shaded}

\begin{verbatim}
Are they equal?    TRUE 
\end{verbatim}

\begin{Shaded}
\begin{Highlighting}[]
\FunctionTok{cat}\NormalTok{(}\StringTok{"95\% Confidence Intervals:}\SpecialCharTok{\textbackslash{}n}\StringTok{"}\NormalTok{)}
\end{Highlighting}
\end{Shaded}

\begin{verbatim}
95% Confidence Intervals:
\end{verbatim}

\begin{Shaded}
\begin{Highlighting}[]
\FunctionTok{cat}\NormalTok{(}\FunctionTok{sprintf}\NormalTok{(}\StringTok{"lm()  CI: [\%.6f, \%.6f]}\SpecialCharTok{\textbackslash{}n}\StringTok{"}\NormalTok{, lm\_ci[}\DecValTok{1}\NormalTok{], lm\_ci[}\DecValTok{2}\NormalTok{]))}
\end{Highlighting}
\end{Shaded}

\begin{verbatim}
lm()  CI: [23038.030705, 24659.038379]
\end{verbatim}

\begin{Shaded}
\begin{Highlighting}[]
\FunctionTok{cat}\NormalTok{(}\FunctionTok{sprintf}\NormalTok{(}\StringTok{"ols() CI: [\%.6f, \%.6f]}\SpecialCharTok{\textbackslash{}n}\StringTok{"}\NormalTok{, ols\_ci\_lower, ols\_ci\_upper))}
\end{Highlighting}
\end{Shaded}

\begin{verbatim}
ols() CI: [23038.030705, 24659.038379]
\end{verbatim}

\begin{Shaded}
\begin{Highlighting}[]
\FunctionTok{cat}\NormalTok{(}\StringTok{"Are CIs equal?    "}\NormalTok{, }
    \FunctionTok{all.equal}\NormalTok{(lm\_ci[}\DecValTok{1}\NormalTok{], ols\_ci\_lower, }\AttributeTok{tolerance =} \FloatTok{1e{-}8}\NormalTok{) }\SpecialCharTok{\&\&} 
    \FunctionTok{all.equal}\NormalTok{(lm\_ci[}\DecValTok{2}\NormalTok{], ols\_ci\_upper, }\AttributeTok{tolerance =} \FloatTok{1e{-}8}\NormalTok{), }\StringTok{"}\SpecialCharTok{\textbackslash{}n}\StringTok{"}\NormalTok{)}
\end{Highlighting}
\end{Shaded}

\begin{verbatim}
Are CIs equal?     TRUE 
\end{verbatim}

\begin{Shaded}
\begin{Highlighting}[]
\CommentTok{\# What additional output does ols() provide that lm() does not?}

\CommentTok{\# When cross comparing the outputs of ols() and lm(), we can see that ols() provides a number of different}
\CommentTok{\# features/statistical measures that give different statistical aggregated measures about the data that lm() does not have.}
\CommentTok{\# For example, ols() provides features/variables such as Lower 0.95 and Upper 0.95 that shows the lower and upper bonds of a 95\% confidence interval}
\CommentTok{\# along with Low and High features/variables that highlights the maximum and minimum values of the data along with a Diff = Difference feature that outlines}
\CommentTok{\# the difference between the maximum and minimum values in the ols() outputs and finally it has a Type feature that highlights the type of the data.}
\CommentTok{\# While the output of lm() only uniquely posseses the t value and Pr(\textgreater{}|t|) features/variables that highlights the p{-}value of a observation that is higher than the t{-}test statistic value and}
\CommentTok{\# the actual t{-}statistic value of the data.}
\CommentTok{\# The only features/variables that these outputs actually shares is the Estimate and Std. Error features/variables along with the Effect and S.E. features/variables that is partically the same features/variables}
\CommentTok{\# which is that they are the Coefficient β and Standard Error features/variables for each model.}


\CommentTok{\# 5. The dataset includes region as a categorical predictor with four levels. }
\CommentTok{\# How many indicator variables does R create? Which level is the reference? }
\CommentTok{\# Refit the model using cell means coding for region ({-}1 in the formula) and }
\CommentTok{\# verify that the predicted values are identical to the reference cell model.}

\CommentTok{\# Cell means linear model for the reference coding ONLY for the region feature/variable}
\NormalTok{lm\_cellmeans\_for\_only\_region }\OtherTok{\textless{}{-}} \FunctionTok{lm}\NormalTok{(}
\NormalTok{  charges }\SpecialCharTok{\textasciitilde{}}\NormalTok{ age }\SpecialCharTok{+}\NormalTok{ bmi }\SpecialCharTok{+}\NormalTok{ children }\SpecialCharTok{+}\NormalTok{ sex }\SpecialCharTok{+}\NormalTok{ smoker }\SpecialCharTok{+}\NormalTok{ region }\SpecialCharTok{{-}} \DecValTok{1}\NormalTok{,}
  \AttributeTok{data =}\NormalTok{ insurance\_data}
\NormalTok{)}

\FunctionTok{cat}\NormalTok{(}\StringTok{"}\SpecialCharTok{\textbackslash{}n}\StringTok{========== Coefficients (Cell Means Model) ==========}\SpecialCharTok{\textbackslash{}n}\StringTok{"}\NormalTok{)}
\end{Highlighting}
\end{Shaded}

\begin{verbatim}

========== Coefficients (Cell Means Model) ==========
\end{verbatim}

\begin{Shaded}
\begin{Highlighting}[]
\FunctionTok{print}\NormalTok{(}\FunctionTok{round}\NormalTok{(}\FunctionTok{coef}\NormalTok{(lm\_cellmeans\_for\_only\_region), }\DecValTok{4}\NormalTok{))}
\end{Highlighting}
\end{Shaded}

\begin{verbatim}
            age             bmi        children       sexfemale         sexmale 
       256.8564        339.1935        475.5005     -11938.5386     -12069.8529 
      smokeryes regionnorthwest regionsoutheast regionsouthwest 
     23848.5345       -352.9639      -1035.0220       -960.0510 
\end{verbatim}

\begin{Shaded}
\begin{Highlighting}[]
\CommentTok{\# Predictions from both models}
\NormalTok{pred\_reference }\OtherTok{\textless{}{-}} \FunctionTok{predict}\NormalTok{(lm\_fit)}
\NormalTok{pred\_cellmeans }\OtherTok{\textless{}{-}} \FunctionTok{predict}\NormalTok{(lm\_cellmeans\_for\_only\_region)}

\FunctionTok{cat}\NormalTok{(}\StringTok{"}\SpecialCharTok{\textbackslash{}n}\StringTok{========== Prediction Comparison ==========}\SpecialCharTok{\textbackslash{}n}\StringTok{"}\NormalTok{)}
\end{Highlighting}
\end{Shaded}

\begin{verbatim}

========== Prediction Comparison ==========
\end{verbatim}

\begin{Shaded}
\begin{Highlighting}[]
\FunctionTok{cat}\NormalTok{(}\StringTok{"Predictions identical? "}\NormalTok{,}
  \FunctionTok{isTRUE}\NormalTok{(}\FunctionTok{all.equal}\NormalTok{(pred\_reference,}
\NormalTok{                   pred\_cellmeans,}
                   \AttributeTok{tolerance =} \FloatTok{1e{-}8}\NormalTok{)),}\StringTok{"}\SpecialCharTok{\textbackslash{}n}\StringTok{"}\NormalTok{)}
\end{Highlighting}
\end{Shaded}

\begin{verbatim}
Predictions identical?  TRUE 
\end{verbatim}

\begin{Shaded}
\begin{Highlighting}[]
\CommentTok{\# How many indicator variables does R create?}

\CommentTok{\# Since this code references the region feature/variable as a categorical feature/variable for the model fit, then }
\CommentTok{\# the predictor indicator variables that is generated for the region feature/variable is 3 since there is one indicator variable }
\CommentTok{\# generated to serve as a reference level/feature/variable for the region feature/variable. Since this is a region feature/variable generation focused model,}
\CommentTok{\# The three indicator variables generated from the levels of the region feature/variable are regionnorthwest, regionsoutheast and regionsouthwest along with the }
\CommentTok{\# regionnortheast indicator variable that serves as the reference level/feature/variable for the code.}
\end{Highlighting}
\end{Shaded}

\bookmarksetup{startatroot}

\chapter{Name: Maxwell Braseth}\label{name-maxwell-braseth-3}

\bookmarksetup{startatroot}

\chapter{Date: 07/11/2026}\label{date-07112026}

\bookmarksetup{startatroot}

\chapter{Course : DS 6021 - Machine Learning I - Introduction to
Predictive
Modeling}\label{course-ds-6021---machine-learning-i---introduction-to-predictive-modeling}

\bookmarksetup{startatroot}

\chapter{Module 7 Homework}\label{module-7-homework}

\begin{Shaded}
\begin{Highlighting}[]
\CommentTok{\# installing packages}
\CommentTok{\#install.packages("ISLR")}
\CommentTok{\#install.packages("car")}
\CommentTok{\#install.packages("splines")}
\CommentTok{\#install.packages("rms")}
\CommentTok{\#install.packages("ggplot2")}
\CommentTok{\#install.packages("FNN")}

\CommentTok{\# loading packages into the program}
\FunctionTok{require}\NormalTok{(ISLR)}
\end{Highlighting}
\end{Shaded}

\begin{verbatim}
Loading required package: ISLR
\end{verbatim}

\begin{Shaded}
\begin{Highlighting}[]
\FunctionTok{require}\NormalTok{(car)}
\end{Highlighting}
\end{Shaded}

\begin{verbatim}
Loading required package: car
\end{verbatim}

\begin{verbatim}
Loading required package: carData
\end{verbatim}

\begin{Shaded}
\begin{Highlighting}[]
\FunctionTok{require}\NormalTok{(splines)}
\end{Highlighting}
\end{Shaded}

\begin{verbatim}
Loading required package: splines
\end{verbatim}

\begin{Shaded}
\begin{Highlighting}[]
\FunctionTok{require}\NormalTok{(rms)}
\end{Highlighting}
\end{Shaded}

\begin{verbatim}
Loading required package: rms
\end{verbatim}

\begin{verbatim}
Loading required package: Hmisc
\end{verbatim}

\begin{verbatim}

Attaching package: 'Hmisc'
\end{verbatim}

\begin{verbatim}
The following objects are masked from 'package:base':

    format.pval, units
\end{verbatim}

\begin{verbatim}

Attaching package: 'rms'
\end{verbatim}

\begin{verbatim}
The following objects are masked from 'package:car':

    Predict, vif
\end{verbatim}

\begin{Shaded}
\begin{Highlighting}[]
\FunctionTok{require}\NormalTok{(ggplot2)}
\end{Highlighting}
\end{Shaded}

\begin{verbatim}
Loading required package: ggplot2
\end{verbatim}

\begin{Shaded}
\begin{Highlighting}[]
\FunctionTok{require}\NormalTok{(FNN)}
\end{Highlighting}
\end{Shaded}

\begin{verbatim}
Loading required package: FNN
\end{verbatim}

\begin{Shaded}
\begin{Highlighting}[]
\FunctionTok{library}\NormalTok{(ISLR)}
\FunctionTok{library}\NormalTok{(car)}
\FunctionTok{library}\NormalTok{(splines)}
\FunctionTok{library}\NormalTok{(rms)}
\FunctionTok{library}\NormalTok{(ggplot2)}
\FunctionTok{library}\NormalTok{(FNN)}

\CommentTok{\# load Wage dataset}
\FunctionTok{data}\NormalTok{(Wage)}

\CommentTok{\# Task 1. Fit a linear model with age (and any relevant covariates,}
\CommentTok{\# such as education or race) as predictors. }
\CommentTok{\# Create a partial effect plot of age and a component + residual plot (crPlots from the car package). }
\CommentTok{\# Does the relationship between age and wage look linear?}

\CommentTok{\# Rename all the columns of the Wage dataset}
\FunctionTok{colnames}\NormalTok{(Wage) }\OtherTok{\textless{}{-}} \FunctionTok{c}\NormalTok{(}\StringTok{"year"}\NormalTok{, }\StringTok{"age"}\NormalTok{, }\StringTok{"marital\_status"}\NormalTok{, }\StringTok{"race"}\NormalTok{, }\StringTok{"education"}\NormalTok{, }\StringTok{"region"}\NormalTok{,}
 \StringTok{"job\_class"}\NormalTok{, }\StringTok{"health\_status"}\NormalTok{, }\StringTok{"health\_insurance"}\NormalTok{, }\StringTok{"log\_wage"}\NormalTok{, }\StringTok{"wages"}\NormalTok{)}

\CommentTok{\# fitting a linear model with age and wage as predictors along with other relevant covariates such as education, etc...}
\NormalTok{lm\_age }\OtherTok{\textless{}{-}} \FunctionTok{lm}\NormalTok{(wages }\SpecialCharTok{\textasciitilde{}}\NormalTok{ age }\SpecialCharTok{+}\NormalTok{ education }\SpecialCharTok{+}\NormalTok{ race }\SpecialCharTok{+}\NormalTok{ job\_class }\SpecialCharTok{+}\NormalTok{ marital\_status, }\AttributeTok{x =} \ConstantTok{TRUE}\NormalTok{, }\AttributeTok{y =} \ConstantTok{TRUE}\NormalTok{, }\AttributeTok{data =}\NormalTok{ Wage)}

\CommentTok{\# creating a partial effect plot and component + residual plots to visualize this relationship}
\FunctionTok{crPlots}\NormalTok{(lm\_age,}
       \AttributeTok{main =} \StringTok{"The Partial Effect Plot and a Component + Residual Plot: Wages by Age"}\NormalTok{,}
       \AttributeTok{ylab =} \StringTok{"Wages ($)"}\NormalTok{,}
       \AttributeTok{col =} \StringTok{"steelblue"}\NormalTok{,}
       \AttributeTok{pch =} \DecValTok{16}\NormalTok{,}
       \AttributeTok{lwd =} \FloatTok{2.5}\NormalTok{,}
       \AttributeTok{smooth =}\NormalTok{ loessLine,}
       \AttributeTok{span =} \FloatTok{0.7}\NormalTok{)}
\end{Highlighting}
\end{Shaded}

\begin{verbatim}
Warning in applyDefaults(smooth, defaults = list(smoother = loessLine, var =
FALSE), : unnamed smooth arguments, will be ignored
\end{verbatim}

\begin{verbatim}
Warning in plot.window(...): "span" is not a graphical parameter
\end{verbatim}

\begin{verbatim}
Warning in plot.xy(xy, type, ...): "span" is not a graphical parameter
\end{verbatim}

\begin{verbatim}
Warning in axis(side = side, at = at, labels = labels, ...): "span" is not a
graphical parameter
Warning in axis(side = side, at = at, labels = labels, ...): "span" is not a
graphical parameter
\end{verbatim}

\begin{verbatim}
Warning in box(...): "span" is not a graphical parameter
\end{verbatim}

\begin{verbatim}
Warning in title(...): "span" is not a graphical parameter
\end{verbatim}

\begin{verbatim}
Warning in applyDefaults(smooth, defaults = list(smoother = loessLine, var =
FALSE), : unnamed smooth arguments, will be ignored
Warning in applyDefaults(smooth, defaults = list(smoother = loessLine, var =
FALSE), : unnamed smooth arguments, will be ignored
Warning in applyDefaults(smooth, defaults = list(smoother = loessLine, var =
FALSE), : unnamed smooth arguments, will be ignored
Warning in applyDefaults(smooth, defaults = list(smoother = loessLine, var =
FALSE), : unnamed smooth arguments, will be ignored
\end{verbatim}

\pandocbounded{\includegraphics[keepaspectratio]{HW7_files/figure-pdf/unnamed-chunk-1-1.pdf}}

\begin{Shaded}
\begin{Highlighting}[]
\CommentTok{\# Does the relationship between age and wage look linear?}

\CommentTok{\# From the partial effect plot and component + residual plots listed above, }
\CommentTok{\# it appears that the relationship between age and wage does not appear to be strictly linear and mostly non{-}linear. }
\CommentTok{\# In fact, the partial effect plot that shows the relationship between age and wages shows a large collection of data points that clusters together around the ages from 10{-}70 years old, but the clusters of the data}
\CommentTok{\# are not smooth and linear. Instead, the shape of the data points distribution is more scattered and does not appear to follow a straight line, but }
\CommentTok{\# instead appears to show the data points being blotched together and randomly scattered in a way that does not follow a linear pattern. }
\CommentTok{\# The only thing to note is that the Purple (solid) line that represents the LOESS (locally weighted smoothing) in this plot shows that}
\CommentTok{\# the data points do not appear to follow a linear pattern due to how the LOESS curve bends slightly downward and even deviated away from the }
\CommentTok{\# linear dottted Blue (dotted) line that represents the linear regression line implied by the fitted model to show the relationship of wages and ages along with the other }
\CommentTok{\# covariates (education, race, job class, marital status, etc..). From these details, it can be concluded that the relationship between age and wages is not linear along}
\CommentTok{\# with the other covariates with being non{-}linear, not smooth and does not have good relationships between the wages and ages along with the covariates.}


\CommentTok{\# Task 2. Fit a model replacing the linear age term with a 4{-}knot restricted cubic spline (rcs(age, 4)}
\CommentTok{\# from the rms package). Create a partial effect plot of the RCS term. }
\CommentTok{\# How does the implied curve differ from the linear fit?}

\CommentTok{\# Wage data being restructured to ensure that the column names are consistent with the previous code and to avoid any potential issues with column name mismatches.}
\NormalTok{dd }\OtherTok{\textless{}{-}} \FunctionTok{datadist}\NormalTok{(Wage)}
\FunctionTok{options}\NormalTok{(}\AttributeTok{datadist =} \StringTok{"dd"}\NormalTok{)}

\CommentTok{\# fitting a model replacing the linear age term with a 4{-}knot restricted cubic spline}
\NormalTok{rcs\_4\_cubic\_spline\_model }\OtherTok{\textless{}{-}} \FunctionTok{ols}\NormalTok{(}
\NormalTok{  wages }\SpecialCharTok{\textasciitilde{}} \FunctionTok{rcs}\NormalTok{(age, }\DecValTok{4}\NormalTok{) }\SpecialCharTok{+}\NormalTok{ education }\SpecialCharTok{+}\NormalTok{ race }\SpecialCharTok{+}\NormalTok{ job\_class }\SpecialCharTok{+}\NormalTok{ marital\_status,}
  \AttributeTok{data =}\NormalTok{ Wage, }\AttributeTok{x =} \ConstantTok{TRUE}\NormalTok{, }\AttributeTok{y =} \ConstantTok{TRUE}
\NormalTok{)}

\CommentTok{\# Also fit linear version of the model from task 1 for comparison}
\NormalTok{linear\_rms }\OtherTok{\textless{}{-}} \FunctionTok{ols}\NormalTok{(}
\NormalTok{  wages }\SpecialCharTok{\textasciitilde{}}\NormalTok{ age }\SpecialCharTok{+}\NormalTok{ education }\SpecialCharTok{+}\NormalTok{ race }\SpecialCharTok{+}\NormalTok{ job\_class }\SpecialCharTok{+}\NormalTok{ marital\_status,}
  \AttributeTok{data =}\NormalTok{ Wage, }\AttributeTok{x =} \ConstantTok{TRUE}\NormalTok{, }\AttributeTok{y =} \ConstantTok{TRUE}
\NormalTok{)}

\CommentTok{\# Partial effects (both models hold other covariates at datadist reference values)}
\NormalTok{p\_spline }\OtherTok{\textless{}{-}} \FunctionTok{Predict}\NormalTok{(rcs\_4\_cubic\_spline\_model, age, }\AttributeTok{conf.int =} \ConstantTok{FALSE}\NormalTok{)}
\NormalTok{p\_linear }\OtherTok{\textless{}{-}} \FunctionTok{Predict}\NormalTok{(linear\_rms, age, }\AttributeTok{conf.int =} \ConstantTok{FALSE}\NormalTok{)}

\CommentTok{\# Knot locations (for annotation)}
\NormalTok{knots }\OtherTok{\textless{}{-}} \FunctionTok{rcspline.eval}\NormalTok{(Wage}\SpecialCharTok{$}\NormalTok{age, }\AttributeTok{nk =} \DecValTok{4}\NormalTok{, }\AttributeTok{knots.only =} \ConstantTok{TRUE}\NormalTok{)}
\CommentTok{\# Partial effects for the knot locations to at knots locations}
\NormalTok{p\_knots }\OtherTok{\textless{}{-}} \FunctionTok{Predict}\NormalTok{(rcs\_4\_cubic\_spline\_model, }\AttributeTok{age =}\NormalTok{ knots)}

\CommentTok{\# plotting the partial effect plot}
\FunctionTok{plot}\NormalTok{(p\_spline}\SpecialCharTok{$}\NormalTok{age, p\_spline}\SpecialCharTok{$}\NormalTok{yhat,}
     \AttributeTok{type =} \StringTok{"l"}\NormalTok{, }\AttributeTok{lwd =} \FloatTok{2.8}\NormalTok{, }\AttributeTok{col =} \StringTok{"darkorange"}\NormalTok{,}
     \AttributeTok{xlab =} \StringTok{"Age (years)"}\NormalTok{, }\AttributeTok{ylab =} \StringTok{"Partial Effect on Wages ($)"}\NormalTok{,}
     \AttributeTok{main =} \StringTok{"Partial Effect of Age on Wages}\SpecialCharTok{\textbackslash{}n}\StringTok{Linear vs 4{-}Knot Restricted Cubic Spline Curve"}\NormalTok{)}

\CommentTok{\# Linear partial effect placed into the plot for comparison}
\FunctionTok{lines}\NormalTok{(p\_linear}\SpecialCharTok{$}\NormalTok{age, p\_linear}\SpecialCharTok{$}\NormalTok{yhat, }\AttributeTok{col =} \StringTok{"blue"}\NormalTok{, }\AttributeTok{lwd =} \FloatTok{2.2}\NormalTok{, }\AttributeTok{lty =} \DecValTok{2}\NormalTok{)}

\CommentTok{\# Setting Knots linear lines}
\FunctionTok{abline}\NormalTok{(}\AttributeTok{v =}\NormalTok{ knots, }\AttributeTok{col =} \StringTok{"red"}\NormalTok{, }\AttributeTok{lty =} \DecValTok{3}\NormalTok{, }\AttributeTok{lwd =} \FloatTok{1.6}\NormalTok{)}
\CommentTok{\# Add knot points}
\FunctionTok{points}\NormalTok{(p\_knots}\SpecialCharTok{$}\NormalTok{age, p\_knots}\SpecialCharTok{$}\NormalTok{yhat, }\AttributeTok{col =} \StringTok{"red"}\NormalTok{, }\AttributeTok{pch =} \DecValTok{19}\NormalTok{, }\AttributeTok{cex =} \FloatTok{1.3}\NormalTok{)}
\CommentTok{\# Add knot labels}
\FunctionTok{text}\NormalTok{(p\_knots}\SpecialCharTok{$}\NormalTok{age, p\_knots}\SpecialCharTok{$}\NormalTok{yhat, }\AttributeTok{labels =} \FunctionTok{round}\NormalTok{(knots, }\DecValTok{1}\NormalTok{), }\AttributeTok{col =} \StringTok{"red"}\NormalTok{, }\AttributeTok{pos =} \DecValTok{3}\NormalTok{, }\AttributeTok{cex =} \FloatTok{0.85}\NormalTok{)}

\CommentTok{\# setting the legend to determine the curved and linear partial effect values along with the knot values in the plot}
\FunctionTok{legend}\NormalTok{(}\StringTok{"topright"}\NormalTok{,}
       \AttributeTok{legend =} \FunctionTok{c}\NormalTok{(}\StringTok{"RCS Partial Effect (4 knots)"}\NormalTok{, }
                  \StringTok{"Linear Partial Effect"}\NormalTok{, }\StringTok{"Knot locations"}\NormalTok{),}
       \AttributeTok{col =} \FunctionTok{c}\NormalTok{(}\StringTok{"darkorange"}\NormalTok{, }\StringTok{"blue"}\NormalTok{, }\StringTok{"red"}\NormalTok{),}
       \AttributeTok{lwd =} \FunctionTok{c}\NormalTok{(}\FloatTok{2.8}\NormalTok{, }\FloatTok{2.2}\NormalTok{, }\FloatTok{1.6}\NormalTok{),}
       \AttributeTok{lty =} \FunctionTok{c}\NormalTok{(}\DecValTok{1}\NormalTok{, }\DecValTok{2}\NormalTok{, }\DecValTok{3}\NormalTok{),}
       \AttributeTok{pch =} \FunctionTok{c}\NormalTok{(}\ConstantTok{NA}\NormalTok{, }\ConstantTok{NA}\NormalTok{, }\DecValTok{19}\NormalTok{),}
       \AttributeTok{bty =} \StringTok{"n"}\NormalTok{, }\AttributeTok{cex =} \FloatTok{0.9}\NormalTok{)}
\end{Highlighting}
\end{Shaded}

\pandocbounded{\includegraphics[keepaspectratio]{HW7_files/figure-pdf/unnamed-chunk-1-2.pdf}}

\begin{Shaded}
\begin{Highlighting}[]
\CommentTok{\# How does the implied curve differ from the linear fit?}
\CommentTok{\# The implied curve from the 4{-}knot restricted cubic spline (RCS) model (orange line) differs from the linear fit model (blue line) in several ways.}
\CommentTok{\# The Restricted Cubic Spline (RCS) model allows for more flexibility in capturing the relationship between age and wages, }
\CommentTok{\# resulting in a curve that can bend and adapt to the data. On top of that, the RCS model can capture non{-}linear patterns and }
\CommentTok{\# the knots at the specified locations in the model allows a slight rate of change and/or slope change to occur in the relationship }
\CommentTok{\# to where the data values in the relationship could experience a sudden exponential increase vertically, a slight increase and then decrease to create a small curve that signals}
\CommentTok{\# a regression back to the previous linear relationship between the wages and ages, a sudden plateau where the data values in the relationship could experience a sudden decrease in a curved shape}
\CommentTok{\# and finally another sudden decrease in values that will show the relationship of the ages and wages to be weaker than usual in the relationship.}
\CommentTok{\# In comparision to the original linear fit relationship, which shows a constant linear and positive rate of change across all ages to wages. The only thing to note is that}
\CommentTok{\# the linear fitted model tends to base its postive linear relationship on the data values that actually have a positive linear relationship between the ages and wages}
\CommentTok{\# which is reflected in how all of the data values of the blue line starts around wages of 92.5 and goes up vertically as the line goes through all of the ages in comparison to the}
\CommentTok{\# orange line which starts its wages measures around 72.5 and its line experience many changes (both vertically and horizontally) as the data line relationship goes through all of the ages values.}
\CommentTok{\# Altogether, the RCS model provides a more nuanced and accurate representation of the relationship between age and wages along with the other covariates in the model in comparison to the linear fit model}
\CommentTok{\# of task 1 which only reflects a somewhat artificial relationship between age and wages along with the covariates in the model.}


\CommentTok{\# Task 3. Compute the cross{-}validated MSE for both models using k{-}fold cross{-}validation (following the approach from earlier modules)}
\CommentTok{\#. Does the spline model offer a meaningful improvement in predictive accuracy?}

\NormalTok{total\_rows\_in\_wages\_data }\OtherTok{\textless{}{-}} \FunctionTok{nrow}\NormalTok{(Wage)}

\CommentTok{\# new restructuring of Wage data}
\NormalTok{dd }\OtherTok{\textless{}{-}} \FunctionTok{datadist}\NormalTok{(Wage)}
\FunctionTok{options}\NormalTok{(}\AttributeTok{datadist =} \StringTok{"dd"}\NormalTok{)}

\CommentTok{\# setting k{-}fold cross{-}validation k values}
\NormalTok{k\_vals }\OtherTok{\textless{}{-}} \DecValTok{20}                            
\NormalTok{folds }\OtherTok{\textless{}{-}}\NormalTok{ folds }\OtherTok{\textless{}{-}} \FunctionTok{sample}\NormalTok{(}\FunctionTok{rep}\NormalTok{(}\DecValTok{1}\SpecialCharTok{:}\NormalTok{k\_vals, }\AttributeTok{length.out =}\NormalTok{ total\_rows\_in\_wages\_data))}

\CommentTok{\# Function that computes MSE for one fold}
\NormalTok{compute\_k\_cross\_validation\_mse\_value }\OtherTok{\textless{}{-}} \ControlFlowTok{function}\NormalTok{(fold\_num) \{}
  \CommentTok{\# setting the training and testing data}
\NormalTok{  train\_data }\OtherTok{\textless{}{-}}\NormalTok{ Wage[folds }\SpecialCharTok{!=}\NormalTok{ fold\_num, ]}
\NormalTok{  test\_data  }\OtherTok{\textless{}{-}}\NormalTok{ Wage[folds }\SpecialCharTok{==}\NormalTok{ fold\_num, ]}
  
  \CommentTok{\# Fit BOTH models on training data}
\NormalTok{  fit\_lin }\OtherTok{\textless{}{-}} \FunctionTok{ols}\NormalTok{(}
\NormalTok{    wages }\SpecialCharTok{\textasciitilde{}}\NormalTok{ age }\SpecialCharTok{+}\NormalTok{ education }\SpecialCharTok{+}\NormalTok{ race }\SpecialCharTok{+}\NormalTok{ job\_class }\SpecialCharTok{+}\NormalTok{ marital\_status,}
    \AttributeTok{data =}\NormalTok{ train\_data, }\AttributeTok{x =} \ConstantTok{TRUE}\NormalTok{, }\AttributeTok{y =} \ConstantTok{TRUE}
\NormalTok{  )}
  
\NormalTok{  fit\_rcs }\OtherTok{\textless{}{-}} \FunctionTok{ols}\NormalTok{(}
\NormalTok{    wages }\SpecialCharTok{\textasciitilde{}} \FunctionTok{rcs}\NormalTok{(age, }\DecValTok{4}\NormalTok{) }\SpecialCharTok{+}\NormalTok{ education }\SpecialCharTok{+}\NormalTok{ race }\SpecialCharTok{+}\NormalTok{ job\_class }\SpecialCharTok{+}\NormalTok{ marital\_status,}
    \AttributeTok{data =}\NormalTok{ train\_data, }\AttributeTok{x =} \ConstantTok{TRUE}\NormalTok{, }\AttributeTok{y =} \ConstantTok{TRUE}     \CommentTok{\# ← FIXED: now using train\_data}
\NormalTok{  )}
  
  \CommentTok{\# Predict on the }\AlertTok{TEST}\CommentTok{ data}
\NormalTok{  pred\_lin }\OtherTok{\textless{}{-}} \FunctionTok{predict}\NormalTok{(fit\_lin, }\AttributeTok{newdata =}\NormalTok{ test\_data)}
\NormalTok{  pred\_rcs }\OtherTok{\textless{}{-}} \FunctionTok{predict}\NormalTok{(fit\_rcs, }\AttributeTok{newdata =}\NormalTok{ test\_data)}
  \CommentTok{\# Return MSEs}
  \FunctionTok{c}\NormalTok{(}
    \AttributeTok{linear\_mse =} \FunctionTok{mean}\NormalTok{((pred\_lin }\SpecialCharTok{{-}}\NormalTok{ test\_data}\SpecialCharTok{$}\NormalTok{wages)}\SpecialCharTok{\^{}}\DecValTok{2}\NormalTok{),}
    \AttributeTok{rcs\_mse    =} \FunctionTok{mean}\NormalTok{((pred\_rcs }\SpecialCharTok{{-}}\NormalTok{ test\_data}\SpecialCharTok{$}\NormalTok{wages)}\SpecialCharTok{\^{}}\DecValTok{2}\NormalTok{)}
\NormalTok{  )}
\NormalTok{\}}

\CommentTok{\# Apply the function to each fold}
\NormalTok{cv\_results }\OtherTok{\textless{}{-}} \FunctionTok{lapply}\NormalTok{(}\DecValTok{1}\SpecialCharTok{:}\NormalTok{k\_vals, compute\_k\_cross\_validation\_mse\_value)}

\CommentTok{\# Convert results to data frame and calculate averages}
\NormalTok{cv\_results\_df }\OtherTok{\textless{}{-}} \FunctionTok{do.call}\NormalTok{(rbind, cv\_results)}

\NormalTok{mean\_cv\_mse\_linear }\OtherTok{\textless{}{-}} \FunctionTok{mean}\NormalTok{(cv\_results\_df[, }\StringTok{"linear\_mse"}\NormalTok{])}
\NormalTok{mean\_cv\_mse\_rcs    }\OtherTok{\textless{}{-}} \FunctionTok{mean}\NormalTok{(cv\_results\_df[, }\StringTok{"rcs\_mse"}\NormalTok{])}

\FunctionTok{cat}\NormalTok{(}\StringTok{"Linear Model CV MSE:   "}\NormalTok{, }\FunctionTok{round}\NormalTok{(mean\_cv\_mse\_linear, }\DecValTok{2}\NormalTok{), }\StringTok{"}\SpecialCharTok{\textbackslash{}n}\StringTok{"}\NormalTok{)}
\end{Highlighting}
\end{Shaded}

\begin{verbatim}
Linear Model CV MSE:    1236.25 
\end{verbatim}

\begin{Shaded}
\begin{Highlighting}[]
\FunctionTok{cat}\NormalTok{(}\StringTok{"RCS Spline Model CV MSE:"}\NormalTok{, }\FunctionTok{round}\NormalTok{(mean\_cv\_mse\_rcs, }\DecValTok{2}\NormalTok{), }\StringTok{"}\SpecialCharTok{\textbackslash{}n}\StringTok{"}\NormalTok{)}
\end{Highlighting}
\end{Shaded}

\begin{verbatim}
RCS Spline Model CV MSE: 1210.75 
\end{verbatim}

\begin{Shaded}
\begin{Highlighting}[]
\FunctionTok{cat}\NormalTok{(}\StringTok{"Difference between models (RCS {-} Linear):"}\NormalTok{, }\FunctionTok{round}\NormalTok{(mean\_cv\_mse\_rcs }\SpecialCharTok{{-}}\NormalTok{ mean\_cv\_mse\_linear, }\DecValTok{2}\NormalTok{), }\StringTok{"}\SpecialCharTok{\textbackslash{}n}\StringTok{"}\NormalTok{)}
\end{Highlighting}
\end{Shaded}

\begin{verbatim}
Difference between models (RCS - Linear): -25.5 
\end{verbatim}

\begin{Shaded}
\begin{Highlighting}[]
\FunctionTok{cat}\NormalTok{(}\StringTok{"Percentage improvement:"}\NormalTok{, }\FunctionTok{round}\NormalTok{((mean\_cv\_mse\_rcs }\SpecialCharTok{{-}}\NormalTok{ mean\_cv\_mse\_linear) }\SpecialCharTok{/}\NormalTok{ mean\_cv\_mse\_linear }\SpecialCharTok{*} \DecValTok{100}\NormalTok{, }\DecValTok{2}\NormalTok{), }\StringTok{"\%}\SpecialCharTok{\textbackslash{}n}\StringTok{"}\NormalTok{)}
\end{Highlighting}
\end{Shaded}

\begin{verbatim}
Percentage improvement: -2.06 %
\end{verbatim}

\begin{Shaded}
\begin{Highlighting}[]
\CommentTok{\# Does the spline model offer a meaningful improvement in predictive accuracy in comparison with the linear model?}

\CommentTok{\# From the calculations that we see from this code, it would be easy to interpret that the Spline or RCS model}
\CommentTok{\# does not offer a meaningful improvement in predictive accuracy because the difference between the two models is }
\CommentTok{\# noticable even if it is noticably small. When the age and wage data from the Linear (training) model and RCS (testing) spline model is loaded into}
\CommentTok{\# the compute\_k\_cross\_validation\_mse\_value function along with the specified k value (1{-}20 in total), we can start to see that}
\CommentTok{\# training data and testing data are very different and as they attempt to predict each other\textquotesingle{}s y values with each iteration of k,}
\CommentTok{\# they usually produce wage y values that are different from each other. In particular, the predicted y values from the RCS (testing) spline model and}
\CommentTok{\# the Linear (training) model are always differing from each other regardless of the k value used in this cross valdidation process which in term }
\CommentTok{\# means that there is always a difference in the calculated MSE values between the two models. In this context though, the MSE Cross Validation value for}
\CommentTok{\# the RCS (testing) spline model is much smaller than the MSE Cross Validation value for the linear model.}
\CommentTok{\# As a result of these results, we can conclude the that the Spline or RCS model does not offer a meaningful improvement in predictive accuracy in comparison with the linear model,}
\CommentTok{\# at least in the terms of believing that predictive accuracy is determined by the pureness of the linear line of data that these models are suppose to follow.}
\CommentTok{\# From another context though, the Spline or RCS model does offer a more nuanced and complex predictive accuracy in terms of calculating a K{-}Fold Cross Validation MSE value}
\CommentTok{\# that takes under the consideration all of the values and changes that comes with the covariates still left in the models when considering that a RCS model is able to capture non{-}linear behavior}
\CommentTok{\# in its model in comparison to the linear model that only seems to be seeking out}
\CommentTok{\# a pure linear relationship between the ages and wages along with the other covariates in the models.}


\CommentTok{\#Task 4. Create a calibration plot for both models. Does either model show systematic }
\CommentTok{\# miscalibration across the predicted range?}

\CommentTok{\# New Wage data being restructured to avoid column and row mismatches}
\NormalTok{dd }\OtherTok{\textless{}{-}} \FunctionTok{datadist}\NormalTok{(Wage)}
\FunctionTok{options}\NormalTok{(}\AttributeTok{datadist =} \StringTok{"dd"}\NormalTok{)}

\CommentTok{\# recreating Linear and RCS spline models for task 4}

\CommentTok{\# fitting a model for 4{-}knot restricted cubic spline}
\NormalTok{rcs\_spline\_model }\OtherTok{\textless{}{-}} \FunctionTok{ols}\NormalTok{(}
\NormalTok{  wages }\SpecialCharTok{\textasciitilde{}} \FunctionTok{rcs}\NormalTok{(age, }\DecValTok{4}\NormalTok{) }\SpecialCharTok{+}\NormalTok{ education }\SpecialCharTok{+}\NormalTok{ race }\SpecialCharTok{+}\NormalTok{ job\_class }\SpecialCharTok{+}\NormalTok{ marital\_status,}
  \AttributeTok{data =}\NormalTok{ Wage, }\AttributeTok{x =} \ConstantTok{TRUE}\NormalTok{, }\AttributeTok{y =} \ConstantTok{TRUE}
\NormalTok{)}

\CommentTok{\# fitting a model for being linear from task 1}
\NormalTok{linear\_rms }\OtherTok{\textless{}{-}} \FunctionTok{ols}\NormalTok{(}
\NormalTok{  wages }\SpecialCharTok{\textasciitilde{}}\NormalTok{ age }\SpecialCharTok{+}\NormalTok{ education }\SpecialCharTok{+}\NormalTok{ race }\SpecialCharTok{+}\NormalTok{ job\_class }\SpecialCharTok{+}\NormalTok{ marital\_status,}
  \AttributeTok{data =}\NormalTok{ Wage, }\AttributeTok{x =} \ConstantTok{TRUE}\NormalTok{, }\AttributeTok{y =} \ConstantTok{TRUE}
\NormalTok{)}

\CommentTok{\# Predictions for partial effect plots for age}
\NormalTok{rcs\_spline\_prediction }\OtherTok{\textless{}{-}} \FunctionTok{Predict}\NormalTok{(rcs\_4\_cubic\_spline\_model, age, }\AttributeTok{conf.int =} \ConstantTok{FALSE}\NormalTok{)}
\NormalTok{linear\_prediction }\OtherTok{\textless{}{-}} \FunctionTok{Predict}\NormalTok{(linear\_rms, age, }\AttributeTok{conf.int =} \ConstantTok{FALSE}\NormalTok{)}

\NormalTok{calibration\_of\_training\_testing\_models\_df }\OtherTok{\textless{}{-}} \FunctionTok{data.frame}\NormalTok{(}
  \AttributeTok{original\_observed\_values =}\NormalTok{ Wage}\SpecialCharTok{$}\NormalTok{wages,}
  \AttributeTok{prediction\_linear\_values =}\NormalTok{ linear\_prediction}\SpecialCharTok{$}\NormalTok{yhat,}
  \AttributeTok{predictions\_RCS\_spline =}\NormalTok{ rcs\_spline\_prediction}\SpecialCharTok{$}\NormalTok{yhat}
\NormalTok{)}

\NormalTok{make\_calib\_plot }\OtherTok{\textless{}{-}} \ControlFlowTok{function}\NormalTok{(df, pred\_var, title) \{}
  \FunctionTok{ggplot}\NormalTok{(df, }\FunctionTok{aes\_string}\NormalTok{(}\AttributeTok{x =}\NormalTok{ pred\_var, }\AttributeTok{y =} \StringTok{"original\_observed\_values"}\NormalTok{)) }\SpecialCharTok{+}
    \FunctionTok{geom\_point}\NormalTok{(}\AttributeTok{alpha =} \FloatTok{0.3}\NormalTok{, }\AttributeTok{color =} \StringTok{"gray25"}\NormalTok{, }\AttributeTok{size =} \FloatTok{0.8}\NormalTok{) }\SpecialCharTok{+}
    \FunctionTok{geom\_smooth}\NormalTok{(}\AttributeTok{method =} \StringTok{"loess"}\NormalTok{, }\AttributeTok{se =} \ConstantTok{TRUE}\NormalTok{, }\AttributeTok{color =} \StringTok{"steelblue"}\NormalTok{, }\AttributeTok{linewidth =} \FloatTok{1.2}\NormalTok{) }\SpecialCharTok{+}
    \FunctionTok{geom\_abline}\NormalTok{(}\AttributeTok{slope =} \DecValTok{1}\NormalTok{, }\AttributeTok{intercept =} \DecValTok{0}\NormalTok{, }\AttributeTok{color =} \StringTok{"black"}\NormalTok{, }\AttributeTok{linetype =} \StringTok{"dashed"}\NormalTok{, }\AttributeTok{linewidth =} \DecValTok{1}\NormalTok{) }\SpecialCharTok{+}
    \FunctionTok{labs}\NormalTok{(}
      \AttributeTok{title =}\NormalTok{ title,}
      \AttributeTok{subtitle =} \StringTok{"Loess smooth (blue) vs. perfect calibration (dashed line)"}\NormalTok{,}
      \AttributeTok{x =} \StringTok{"Predicted Wages ($)"}\NormalTok{,}
      \AttributeTok{y =} \StringTok{"Observed Wages ($)"}
\NormalTok{    ) }\SpecialCharTok{+}
    \FunctionTok{theme\_bw}\NormalTok{(}\AttributeTok{base\_size =} \DecValTok{11}\NormalTok{) }\SpecialCharTok{+}
    \FunctionTok{theme}\NormalTok{(}
      \AttributeTok{panel.grid.major =} \FunctionTok{element\_line}\NormalTok{(}\AttributeTok{color =} \StringTok{"grey80"}\NormalTok{, }\AttributeTok{linewidth =} \FloatTok{0.35}\NormalTok{),}
      \AttributeTok{panel.grid.minor =} \FunctionTok{element\_line}\NormalTok{(}\AttributeTok{color =} \StringTok{"grey90"}\NormalTok{, }\AttributeTok{linewidth =} \FloatTok{0.2}\NormalTok{),}
      \AttributeTok{plot.title =} \FunctionTok{element\_text}\NormalTok{(}\AttributeTok{face =} \StringTok{"bold"}\NormalTok{, }\AttributeTok{size =} \DecValTok{13}\NormalTok{),}
      \AttributeTok{legend.position =} \StringTok{"none"}
\NormalTok{    )}
\NormalTok{\}}



\NormalTok{linear\_model\_calibration }\OtherTok{\textless{}{-}} \FunctionTok{make\_calib\_plot}\NormalTok{(calibration\_of\_training\_testing\_models\_df, }\StringTok{"prediction\_linear\_values"}\NormalTok{, }\StringTok{"Calibration: Linear Model (Training Data Model)"}\NormalTok{)}
\end{Highlighting}
\end{Shaded}

\begin{verbatim}
Warning: `aes_string()` was deprecated in ggplot2 3.0.0.
i Please use tidy evaluation idioms with `aes()`.
i See also `vignette("ggplot2-in-packages")` for more information.
\end{verbatim}

\begin{Shaded}
\begin{Highlighting}[]
\NormalTok{rcs\_spline\_model\_calibration }\OtherTok{\textless{}{-}} \FunctionTok{make\_calib\_plot}\NormalTok{(calibration\_of\_training\_testing\_models\_df, }\StringTok{"predictions\_RCS\_spline"}\NormalTok{, }\StringTok{"Calibration: 4{-}Knot RCS Spline Model (Testing Data Model)"}\NormalTok{)}

\FunctionTok{print}\NormalTok{(linear\_model\_calibration)}
\end{Highlighting}
\end{Shaded}

\begin{verbatim}
`geom_smooth()` using formula = 'y ~ x'
\end{verbatim}

\pandocbounded{\includegraphics[keepaspectratio]{HW7_files/figure-pdf/unnamed-chunk-1-3.pdf}}

\begin{Shaded}
\begin{Highlighting}[]
\FunctionTok{print}\NormalTok{(rcs\_spline\_model\_calibration)}
\end{Highlighting}
\end{Shaded}

\begin{verbatim}
`geom_smooth()` using formula = 'y ~ x'
\end{verbatim}

\pandocbounded{\includegraphics[keepaspectratio]{HW7_files/figure-pdf/unnamed-chunk-1-4.pdf}}

\begin{Shaded}
\begin{Highlighting}[]
\CommentTok{\# miscalibration across the predicted range?}
\CommentTok{\# After the creation of the calibration plots for both models and the data model that houses all of the data from the models in the code above,}
\CommentTok{\# it would become clear to anyone that there is a huge systematic miscalibration across the predicted range for the linear model in comparison to the 4{-}knot restricted cubic spline model.}
\CommentTok{\# This systematic miscalibration is expressed in the how the 4{-}knot restricted cubic spline model tends to have a line that bends and curves across all of the scattered data values/points of}
\CommentTok{\# the entire distribution observed vs predicted wages data while the pure linear model only shows a pure linear relationship line between the predicted and observed wages data values. From these details,}
\CommentTok{\# it would be clear to say that there is a miscalibration or misalignment in the conclusions of these two lines and both of the lines are attempting to tell a different story about the predicted vs observed}
\CommentTok{\# values of wages. Upon some further evaluations of this data model and its other models, it then become clear to me that the closest accurate calibration of this data model would be the}
\CommentTok{\# 4{-}knot restricted cubic spline model considering that the non{-}linear and non{-}normality relationship between the predicted and observed wages data shows a more nuanced and somewhat accurate}
\CommentTok{\# representation of how the observed data is very close to the precticted wages data since the data across the whole data model is very scattered, shapeless and non{-}linear.}
\CommentTok{\# Altogether though, this data model reveals a huge disconnect and an ability to show a good and accurate calibration plot between the predicted and observed values of wages}
\CommentTok{\# regardless of the linear or 4{-}knot restricted cubic spline model used. In conclusion, there does not really seem to be huge systematic miscalibration in the data lines of these models at all, but instead }
\CommentTok{\# there does not seem to be any good calibration plot (regardless of the model) to show how the predicted values of wages are close to the observed values of wages and thus there is not good relationship between the predicted and observed values of wages.}


\CommentTok{\#Task 5. State a final conclusion: should the nonlinear spline term be retained, }
\CommentTok{\# or is the linear model adequate? Address practical predictive improvement }
\CommentTok{\#(from MSE and calibration).}

\CommentTok{\# In conclusion, the nonlinear 4{-}knot restricted cubic spline model/term should be retained and should be considered to be a somewhat good model to interpret}
\CommentTok{\# the non{-}linear relationship between the predicted and observed values of wages along with interpreting that the wages and ages variables are connected in some way. }
\CommentTok{\# The linear model in this context would not be considered to be a good model to interpret any of the data within this context considering that the linear model is}
\CommentTok{\# only seeking to seek out a pure linear relationship between the ages and wages along with the other covariates in the models and only seeks a linear}
\CommentTok{\# relationship between the predicted and observed values of wages. In other words, the linear model is almost being mutually exclusive in seeking its conclusions and interpreations}
\CommentTok{\# of the relationships within the data in a linear context and structure when it was clear that all of the data in the model was complex, nuanced and scattered from each other which signals}
\CommentTok{\# the need for a non{-}linear relationship measure and formula to find the relationships in the data. The only thing of noteworthyness in terms of the practical improvements that}
\CommentTok{\# could have been made to enforce this need of keeping the non{-}linear 4{-}knot restricted cubic spline model/term would be to improve the MSE score difference to be closer or greater in value}
\CommentTok{\# than the linear model\textquotesingle{}s MSE score in order to make the non{-}linear nature of the 4{-}knot restricted cubic spline model could be approximately normal or greater in value than the linear model.}
\CommentTok{\# Then, the non{-}linear nature of the 4{-}knot restricted cubic spline model would be able to offer a more predictive accuracy improvement in the prediction of the observed values of wages to the}
\CommentTok{\# predicted values of wages along with showing the positive relationships between the wages and the other covariates (such as ages, education, etc...).}
\end{Highlighting}
\end{Shaded}

\bookmarksetup{startatroot}

\chapter{Name: Maxwell Braseth}\label{name-maxwell-braseth-4}

\bookmarksetup{startatroot}

\chapter{Date: 07/12/2026}\label{date-07122026}

\bookmarksetup{startatroot}

\chapter{Course : DS 6021 - Machine Learning I - Introduction to
Predictive
Modeling}\label{course-ds-6021---machine-learning-i---introduction-to-predictive-modeling-1}

\bookmarksetup{startatroot}

\chapter{Module 8 Homework}\label{module-8-homework}

\begin{Shaded}
\begin{Highlighting}[]
\CommentTok{\# installing the packages}
\CommentTok{\#install.packages("glmnet")}
\CommentTok{\#install.packages("pls")}
\CommentTok{\#install.packages("hdi")}
\CommentTok{\#install.packages("lmtest")}


\CommentTok{\# loading and using the packages in the program}
\FunctionTok{library}\NormalTok{(glmnet)   }\CommentTok{\# contains the methods and functions for ridge, lasso, elastic net + cv.glmnet}
\end{Highlighting}
\end{Shaded}

\begin{verbatim}
Loading required package: Matrix
\end{verbatim}

\begin{verbatim}
Loaded glmnet 5.0
\end{verbatim}

\begin{Shaded}
\begin{Highlighting}[]
\FunctionTok{library}\NormalTok{(pls)      }\CommentTok{\# contains the methods and functions for PCR (principal component regression)}
\end{Highlighting}
\end{Shaded}

\begin{verbatim}

Attaching package: 'pls'
\end{verbatim}

\begin{verbatim}
The following object is masked from 'package:stats':

    loadings
\end{verbatim}

\begin{Shaded}
\begin{Highlighting}[]
\FunctionTok{library}\NormalTok{(lmtest)}
\end{Highlighting}
\end{Shaded}

\begin{verbatim}
Loading required package: zoo
\end{verbatim}

\begin{verbatim}

Attaching package: 'zoo'
\end{verbatim}

\begin{verbatim}
The following objects are masked from 'package:base':

    as.Date, as.Date.numeric
\end{verbatim}

\begin{Shaded}
\begin{Highlighting}[]
\CommentTok{\#library(hdi)}

\CommentTok{\# loading the riboflavin dataset}
\CommentTok{\# d1 \textless{}{-} readRDS("data/riboflavin.rds")}
\NormalTok{d1 }\OtherTok{\textless{}{-}} \FunctionTok{readRDS}\NormalTok{(}\StringTok{"/Users/Mzmz12/Documents/UVA Data Science Online Masters/DS6021\_Predictability\_Modeling/Module7\_Linearmodel\_≠\_linear\_relationship/riboflavin.rds"}\NormalTok{)}
\CommentTok{\# extracting the x and y values}
\NormalTok{x }\OtherTok{\textless{}{-}}\NormalTok{ d1}\SpecialCharTok{$}\NormalTok{x   }\CommentTok{\# 71 × 4088 matrix. n = 71, p = 4088}
\NormalTok{y }\OtherTok{\textless{}{-}}\NormalTok{ d1}\SpecialCharTok{$}\NormalTok{y   }\CommentTok{\# length{-}71 vector. n = 71}
\CommentTok{\# extracting the column names}
\NormalTok{x\_column\_names }\OtherTok{\textless{}{-}} \FunctionTok{colnames}\NormalTok{(x)}
\CommentTok{\# setting the seed for reproducibility}
\FunctionTok{set.seed}\NormalTok{(}\DecValTok{123}\NormalTok{)}

\CommentTok{\#Problem 1. Confirm that OLS cannot create a unique solution for this dataset.}
\CommentTok{\# Attempt to fit a linear model using lm() and report what happens.}
\CommentTok{\# How many coefficients does R return as NA, and why?}

\CommentTok{\# fitting the OLS model with y values to the x values will not work}
\CommentTok{\#  due to how x \textless{}{-} d1$x has the dimensions of 71 × 4088 matrix. n = 71, p = 4088}
\CommentTok{\# and p \textgreater{} n = 4088 \textgreater{} 71 violates the OLS requirements for a unique solution }
\CommentTok{\# ols\_fitted\_model \textless{}{-} ols(y \textasciitilde{} x)}

\CommentTok{\# fitting the y values to the x values in a linear model}
\CommentTok{\# this will work regardless of 𝑝 \textgreater{} 𝑛 {-}{-}\textgreater{} rank(X) \textless{} p {-}{-}\textgreater{} due to how d1$x has the dimensions of 71 × 4088 matrix. n = 71, p = 4088}
\CommentTok{\# which is 4088 \textgreater{} 71 that agrees with the lm() requirements for a non{-}unique solution and}
\CommentTok{\# due to the rank(x) \textless{} p which is 71 \textless{} 4088 which is true then lm() assures that there is a non{-}unique solution and that the output non{-}matching }
\CommentTok{\# coefficients from y to x will be NA}

\CommentTok{\# fitting the y values to the x values in a linear model}
\NormalTok{linear\_fitted\_model }\OtherTok{\textless{}{-}} \FunctionTok{lm}\NormalTok{(y }\SpecialCharTok{\textasciitilde{}}\NormalTok{ x)}
\CommentTok{\# calculating the linear algebra rank}
\NormalTok{linear\_rank }\OtherTok{\textless{}{-}} \FunctionTok{qr}\NormalTok{(x)}\SpecialCharTok{$}\NormalTok{rank}


\CommentTok{\# Count how many coefficients are NA}
\NormalTok{number\_of\_na\_coefficients\_in\_linear\_model }\OtherTok{\textless{}{-}} \FunctionTok{sum}\NormalTok{(}\FunctionTok{is.na}\NormalTok{(linear\_fitted\_model}\SpecialCharTok{$}\NormalTok{coefficients))}
\FunctionTok{cat}\NormalTok{(}\StringTok{"Number of coefficients within the Linear Model combination of y to x that returns as NA:"}\NormalTok{, number\_of\_na\_coefficients\_in\_linear\_model, }\StringTok{"}\SpecialCharTok{\textbackslash{}n}\StringTok{"}\NormalTok{)}
\end{Highlighting}
\end{Shaded}

\begin{verbatim}
Number of coefficients within the Linear Model combination of y to x that returns as NA: 4018 
\end{verbatim}

\begin{Shaded}
\begin{Highlighting}[]
\CommentTok{\# The reason why we are getting NA is because the lm() requirements to have a unique solution are not met.}
\CommentTok{\# This is evidenced through how this dataset has 71 observations  {-}{-}\textgreater{} n = 71 and 4088 predictors {-}{-}\textgreater{} p = 4088.}
\CommentTok{\# Due to these n and p values, the lm() requirements to have a unique solution are not met and the full column rank}
\CommentTok{\# of the design matrix X is not correct due to p \textgreater{} n = 4088 \textgreater{} 71. It is also important to note that the requirement}
\CommentTok{\# of rank(x) \textless{} p being met due to 71 \textless{} 4088 that makes X\^{}T * X singular and prevents OLS from having a unique solution.}
\CommentTok{\# It is even important to realize that when the design matrix X is fitted to Y via lm(y \textasciitilde{} x), R uses QR decomposition and}
\CommentTok{\# any predictor that is linearly dependent on others is dropped. The coefficients for these dropped predictors appear as NA.}
\CommentTok{\# So, the NA values calculated here is only a result of the many dependent or non{-}ranked columns in this QR composition linear model}
\CommentTok{\# fitting of y to x will automatically assign NA values to those columns values.}


\CommentTok{\# Problem 2. Fit ridge regression using glmnet (alpha = 0). Select lambda λ}
\CommentTok{\# via 10{-}fold cross{-}validation. Report λmin and λlse (LSE). Plot the Cross Valdidation (CV)}
\CommentTok{\# curve with the two vertical lines. How does CV{-}MSE change as lambda λ increases?}

\CommentTok{\# Ridge Regression: β(ridge) = arg β min ||Y {-} Xβ||\^{}(2, 2) + λ ||β||\^{}(2, 2)}
\CommentTok{\# where λ \textgreater{} 0 is a tuning parameter (penalty) that controls the amount of shrinkage.}
\CommentTok{\# The first term is the OLS loss where λ {-}{-}\textgreater{} 0, then the ridge regression reduces the OLS loss.}
\CommentTok{\# the second term penalizes the large coefficients of this regression line and as λ {-}{-}\textgreater{} infinity,}
\CommentTok{\# then this ridge regression gets closer to 0 (penalty).}

\CommentTok{\# L2 penalty (ridge regression): ||x|| = x\^{}t * x = √(x1 \^{} 2 + x2 \^{} 2 + ... + xd \^{} 2)}

\CommentTok{\# β\^{}(LSE) ​= (X\^{}T * X)\^{}−1 * X\^{}T * y}
\CommentTok{\# X = design matrix of predictors,}
\CommentTok{\# y = vector of observed responses,}
\CommentTok{\# β\^{}(LSE) = estimated regression coefficients.}

\CommentTok{\# Fit ridge regression using glmnet (alpha = 0, fold cross{-}validation = 10).}
\NormalTok{cv\_10\_ridge }\OtherTok{\textless{}{-}} \FunctionTok{cv.glmnet}\NormalTok{(x, y, }\AttributeTok{alpha =} \DecValTok{0}\NormalTok{, }\AttributeTok{nfolds =} \DecValTok{10}\NormalTok{, }\AttributeTok{type.measure =} \StringTok{"mse"}\NormalTok{)}

\CommentTok{\# Plot the Cross Valdidation (CV) curve with the two vertical lines.}
\FunctionTok{plot}\NormalTok{(cv\_10\_ridge, }\AttributeTok{main =} \StringTok{"Ridge Regression – 10{-}fold CV Error Curve"}\NormalTok{)}
\end{Highlighting}
\end{Shaded}

\pandocbounded{\includegraphics[keepaspectratio]{HW8_files/figure-pdf/unnamed-chunk-1-1.pdf}}

\begin{Shaded}
\begin{Highlighting}[]
\CommentTok{\# Getting the LSE and min lambda values}
\NormalTok{lambda\_min\_ridge }\OtherTok{\textless{}{-}}\NormalTok{ cv\_10\_ridge}\SpecialCharTok{$}\NormalTok{lambda.min}
\NormalTok{lambda\_1se\_ridge }\OtherTok{\textless{}{-}}\NormalTok{ cv\_10\_ridge}\SpecialCharTok{$}\NormalTok{lambda}\FloatTok{.1}\NormalTok{se}

\FunctionTok{cat}\NormalTok{(}\StringTok{"Selected lambda (λ) (min CV{-}MSE):"}\NormalTok{, }\FunctionTok{round}\NormalTok{(lambda\_min\_ridge, }\DecValTok{6}\NormalTok{), }\StringTok{"}\SpecialCharTok{\textbackslash{}n}\StringTok{"}\NormalTok{)}
\end{Highlighting}
\end{Shaded}

\begin{verbatim}
Selected lambda (λ) (min CV-MSE): 5.934162 
\end{verbatim}

\begin{Shaded}
\begin{Highlighting}[]
\FunctionTok{cat}\NormalTok{(}\StringTok{"lambda.1se:"}\NormalTok{, }\FunctionTok{round}\NormalTok{(lambda\_1se\_ridge, }\DecValTok{6}\NormalTok{), }\StringTok{"}\SpecialCharTok{\textbackslash{}n\textbackslash{}n}\StringTok{"}\NormalTok{)}
\end{Highlighting}
\end{Shaded}

\begin{verbatim}
lambda.1se: 36.41147 
\end{verbatim}

\begin{Shaded}
\begin{Highlighting}[]
\CommentTok{\# From the ridge regression plot that is produced from this code, we can witness how the Cross Valdidation (CV) Mean Squared Error (MSE)}
\CommentTok{\# curve that describes the diference between this model\textquotesingle{}s y values and this Riboflavin\textquotesingle{}s y production values as a result of their relationship to their x values}
\CommentTok{\# as the lambda λ increases or decreases along the {-}log(λ) x axis. In this case, the Cross Valdidation (CV) Mean Squared Error (MSE) curve is a downward sloping curve}
\CommentTok{\# that is a result of the ridge regression model\textquotesingle{}s ability to reduce the bias in the y values as the lambda λ increases. This decrease of bias is reflecting in how}
\CommentTok{\# when the lambda λ is lower in value (such as {-}(log(λ=6))) then the bias of y is high due to the large difference in the MSE value at this point and how this MSE value}
\CommentTok{\# difference seems to be getting smaller as the lambda λ gets higher in value (such as {-}(log(λ=1 or λ=0))) and thus the bias in the y values decreases. It is also important}
\CommentTok{\# to note that this MSE value difference also reflects the changes in the variance of the y values as the lambda λ goes along the x axis. So, the variance of the MSE value difference}
\CommentTok{\# is at its highest with the y value is biased at {-}(log(λ=6))) and tends to get smaller as the lambda λ decreases and moves along with x axis to get to the lowest variances at}
\CommentTok{\# {-}(log(λ=1 or λ=0)). Finally, the penalty or shrinkage of the ridge regression model is reflected in how the MSE value difference is getting smaller as the lambda λ goes along the x axis.}


\CommentTok{\# Problem 3. Plot the ridge coefficient path (}
\CommentTok{\# on the x{-}axis, coefficient values on the y{-}axis). }
\CommentTok{\# Describe qualitatively what happens to the coefficients as }
\CommentTok{\# increases. Does any coefficient reach exactly zero?}

\FunctionTok{plot}\NormalTok{(cv\_10\_ridge}\SpecialCharTok{$}\NormalTok{glmnet.fit, }
     \AttributeTok{xvar =} \StringTok{"lambda"}\NormalTok{, }
     \AttributeTok{label =} \ConstantTok{FALSE}\NormalTok{,}
     \AttributeTok{main =} \StringTok{"Ridge Coefficient Path (all coefficients shrink toward 0 to show when predicator variables don\textquotesingle{}t effect the 𝑦 outcome)"}\NormalTok{)}
\end{Highlighting}
\end{Shaded}

\pandocbounded{\includegraphics[keepaspectratio]{HW8_files/figure-pdf/unnamed-chunk-1-2.pdf}}

\begin{Shaded}
\begin{Highlighting}[]
\CommentTok{\# L2 penalty (ridge regression): ||x|| = x\^{}t * x = √(x1 \^{} 2 + x2 \^{} 2 + ... + xd \^{} 2)}

\CommentTok{\# From this ridge regression plot and its coefficient pathway on the y axis,}
\CommentTok{\# we can see how the 𝛽 coefficients values are going from smaller (𝛽 = 0.0 {-}{-}\textgreater{} no effect on the predicted y outcomes) to the higher}
\CommentTok{\# (𝛽 = 0.02 {-}{-}\textgreater{} some to a lot of effect on the predicted y outcomes) as the lambda λ goes along the x axis. In this case, the coefficient values are getting}
\CommentTok{\# bigger and you can tell from the lines in the plot that the predict y value outcomes are heavily influenced by 𝛽 coefficient as we move along the lambda λ values along the x axis.}
\CommentTok{\# Thus, the coefficient values are getting quantitatively bigger and the variances of these predicted y values that derives from the linear combinations of these increasing coefficients}
\CommentTok{\# to their x predictor variables are along getting quantitatively bigger and more diverse as we move along the lambda λ x axis values.}
\CommentTok{\# The only two things to note in this case, is that the coefficient of these predicted y values and all of these data lines actually starts at 0.0 when the lambda λ value is at its highest}
\CommentTok{\# {-}log(λ \textgreater{} {-}6). The second thing is that the coefficient values of this plot seem to be kept within a small range of values due to }
\CommentTok{\# the penalties and loss functions that are being used to keep the coefficients from getting too large and imposes the L2 penalty on the ridge regression model to keep the coefficent pathways smooth as it goes from zero to larger values.}
\CommentTok{\# As reflected by this equation, Loss = min𝛽 * ||y {-} x𝛽||\^{}2 + 𝜆 ∑ 𝛽\^{}2j}


\CommentTok{\# Problem 4. Fit lasso regression (alpha = 1) and select λ}
\CommentTok{\# via 10{-}fold CV. At λlse,}
\CommentTok{\# how many gene expression features are retained }
\CommentTok{\# (non{-}zero coefficients)? Plot the coefficient path and the }
\CommentTok{\# CV error curve.}

\CommentTok{\# Fit lasso regression}
\NormalTok{cv\_lasso }\OtherTok{\textless{}{-}} \FunctionTok{cv.glmnet}\NormalTok{(x, y, }\AttributeTok{alpha =} \DecValTok{1}\NormalTok{, }\AttributeTok{nfolds =} \DecValTok{10}\NormalTok{, }\AttributeTok{type.measure =} \StringTok{"mse"}\NormalTok{)}

\CommentTok{\# Getting the LSE and min lambda values}
\NormalTok{lambda\_min\_lasso }\OtherTok{\textless{}{-}}\NormalTok{ cv\_lasso}\SpecialCharTok{$}\NormalTok{lambda.min}
\NormalTok{lambda\_1se\_lasso }\OtherTok{\textless{}{-}}\NormalTok{ cv\_lasso}\SpecialCharTok{$}\NormalTok{lambda}\FloatTok{.1}\NormalTok{se}

\CommentTok{\# Number of non{-}zero gene expression features at lambda.lse}
\NormalTok{coef\_lasso }\OtherTok{\textless{}{-}} \FunctionTok{coef}\NormalTok{(cv\_lasso, }\AttributeTok{s =} \StringTok{"lambda.1se"}\NormalTok{)}
\NormalTok{num\_retained\_lasso }\OtherTok{\textless{}{-}} \FunctionTok{sum}\NormalTok{(coef\_lasso[}\SpecialCharTok{{-}}\DecValTok{1}\NormalTok{, ] }\SpecialCharTok{!=} \DecValTok{0}\NormalTok{)}

\CommentTok{\# Print the Lambda }
\FunctionTok{cat}\NormalTok{(}\StringTok{"lambda.1se (lasso):"}\NormalTok{, }\FunctionTok{round}\NormalTok{(lambda\_1se\_lasso, }\DecValTok{6}\NormalTok{), }\StringTok{"}\SpecialCharTok{\textbackslash{}n}\StringTok{"}\NormalTok{)}
\end{Highlighting}
\end{Shaded}

\begin{verbatim}
lambda.1se (lasso): 0.092316 
\end{verbatim}

\begin{Shaded}
\begin{Highlighting}[]
\FunctionTok{cat}\NormalTok{(}\StringTok{"Number of retained (non{-}zero) gene expression features at lambda.lse:"}\NormalTok{, }
\NormalTok{    num\_retained\_lasso, }\StringTok{"}\SpecialCharTok{\textbackslash{}n\textbackslash{}n}\StringTok{"}\NormalTok{)}
\end{Highlighting}
\end{Shaded}

\begin{verbatim}
Number of retained (non-zero) gene expression features at lambda.lse: 25 
\end{verbatim}

\begin{Shaded}
\begin{Highlighting}[]
\CommentTok{\#Plot CV error curve with the two vertical lines}
\FunctionTok{plot}\NormalTok{(cv\_lasso, }\AttributeTok{main =} \StringTok{"Lasso Regression – 10{-}fold CV Error Curve"}\NormalTok{)}
\end{Highlighting}
\end{Shaded}

\pandocbounded{\includegraphics[keepaspectratio]{HW8_files/figure-pdf/unnamed-chunk-1-3.pdf}}

\begin{Shaded}
\begin{Highlighting}[]
\CommentTok{\# Coefficient path plot}
\FunctionTok{plot}\NormalTok{(cv\_lasso}\SpecialCharTok{$}\NormalTok{glmnet.fit, }
     \AttributeTok{xvar =} \StringTok{"lambda"}\NormalTok{, }
     \AttributeTok{main =} \StringTok{"Lasso Coefficient Path (many coefficients hit exactly zero)"}\NormalTok{)}
\end{Highlighting}
\end{Shaded}

\pandocbounded{\includegraphics[keepaspectratio]{HW8_files/figure-pdf/unnamed-chunk-1-4.pdf}}

\begin{Shaded}
\begin{Highlighting}[]
\CommentTok{\# Problem 5. Fit an elastic net with α ϵ \{0.25, 0.5, 0.75\}.}
\CommentTok{\# For each α, record the minimum CV{-}MSE. }
\CommentTok{\# Create a small table showing (α, λmin, min CV{-}MSE).}
\CommentTok{\# Which α yields the lowest CV{-}MSE? }
\CommentTok{\# How does the number of selected features compare to lasso?}

\CommentTok{\# Setting the alpha, min lambda, min CV{-}MSE, and number of selected elastic net features}
\NormalTok{alphas }\OtherTok{\textless{}{-}} \FunctionTok{seq}\NormalTok{(}\FloatTok{0.25}\NormalTok{, }\FloatTok{0.75}\NormalTok{, }\AttributeTok{by =} \FloatTok{0.25}\NormalTok{)}
\NormalTok{min\_cv\_mse\_elastic\_net }\OtherTok{\textless{}{-}} \FunctionTok{numeric}\NormalTok{(}\FunctionTok{length}\NormalTok{(alphas))}
\NormalTok{lambda\_min\_elastic\_net }\OtherTok{\textless{}{-}} \FunctionTok{numeric}\NormalTok{(}\FunctionTok{length}\NormalTok{(alphas))}
\NormalTok{num\_selected\_elastic\_net }\OtherTok{\textless{}{-}} \FunctionTok{numeric}\NormalTok{(}\FunctionTok{length}\NormalTok{(alphas))}

\CommentTok{\# Fitting elastic net through looping and appending the values to the vectors}
\ControlFlowTok{for}\NormalTok{ (i }\ControlFlowTok{in} \FunctionTok{seq\_along}\NormalTok{(alphas)) \{}
\NormalTok{  cv\_elastic\_net }\OtherTok{\textless{}{-}} \FunctionTok{cv.glmnet}\NormalTok{(x, y, }\AttributeTok{alpha =}\NormalTok{ alphas[i], }\AttributeTok{nfolds =} \DecValTok{10}\NormalTok{, }\AttributeTok{type.measure =} \StringTok{"mse"}\NormalTok{)}
\NormalTok{  min\_cv\_mse\_elastic\_net[i] }\OtherTok{\textless{}{-}} \FunctionTok{min}\NormalTok{(cv\_elastic\_net}\SpecialCharTok{$}\NormalTok{cvm)}
\NormalTok{  lambda\_min\_elastic\_net[i] }\OtherTok{\textless{}{-}}\NormalTok{ cv\_elastic\_net}\SpecialCharTok{$}\NormalTok{lambda.min}
\NormalTok{  coefs }\OtherTok{\textless{}{-}} \FunctionTok{coef}\NormalTok{(cv\_elastic\_net, }\AttributeTok{s =}\NormalTok{ cv\_elastic\_net}\SpecialCharTok{$}\NormalTok{lambda.min)}
\NormalTok{  num\_selected\_elastic\_net[i] }\OtherTok{\textless{}{-}} \FunctionTok{sum}\NormalTok{(coefs[}\SpecialCharTok{{-}}\DecValTok{1}\NormalTok{, ] }\SpecialCharTok{!=} \DecValTok{0}\NormalTok{)}
\NormalTok{\}}

\CommentTok{\# creating a final data frame table}
\NormalTok{results\_elastic\_net }\OtherTok{\textless{}{-}} \FunctionTok{data.frame}\NormalTok{(}
  \AttributeTok{alpha =}\NormalTok{ alphas,}
  \AttributeTok{min\_lambda =} \FunctionTok{round}\NormalTok{(lambda\_min\_elastic\_net, }\DecValTok{6}\NormalTok{),}
  \AttributeTok{min\_CV\_MSE =} \FunctionTok{round}\NormalTok{(min\_cv\_mse\_elastic\_net, }\DecValTok{6}\NormalTok{),}
  \AttributeTok{num\_nonzero\_features =}\NormalTok{ num\_selected\_elastic\_net}
\NormalTok{)}

\CommentTok{\# printing the results and comparing the number of selected features}
\FunctionTok{print}\NormalTok{(results\_elastic\_net)}
\end{Highlighting}
\end{Shaded}

\begin{verbatim}
  alpha min_lambda min_CV_MSE num_nonzero_features
1  0.25   0.167457   0.203687                   69
2  0.50   0.079923   0.222400                   55
3  0.75   0.040306   0.214340                   48
\end{verbatim}

\begin{Shaded}
\begin{Highlighting}[]
\FunctionTok{cat}\NormalTok{(}\StringTok{"The number of selected features from Elastic Net in comparison with Lasso:}\SpecialCharTok{\textbackslash{}n}\StringTok{"}\NormalTok{)}
\end{Highlighting}
\end{Shaded}

\begin{verbatim}
The number of selected features from Elastic Net in comparison with Lasso:
\end{verbatim}

\begin{Shaded}
\begin{Highlighting}[]
\ControlFlowTok{for}\NormalTok{(i }\ControlFlowTok{in} \FunctionTok{seq\_along}\NormalTok{(num\_selected\_elastic\_net))\{}
  \FunctionTok{cat}\NormalTok{(}\StringTok{" The Elastic Net (alpha="}\NormalTok{, alphas[i], }\StringTok{") selected "}\NormalTok{, num\_selected\_elastic\_net[i], }\StringTok{" features.}\SpecialCharTok{\textbackslash{}n}\StringTok{"}\NormalTok{)}
  \FunctionTok{cat}\NormalTok{(}\StringTok{" The Lasso (alpha=1) selected "}\NormalTok{, num\_retained\_lasso, }\StringTok{" features.}\SpecialCharTok{\textbackslash{}n}\StringTok{"}\NormalTok{)}
  \FunctionTok{cat}\NormalTok{(}\StringTok{" The Selected Features of Elastic Net ("}\NormalTok{, num\_selected\_elastic\_net[i], }\StringTok{") compared to Lasso ("}\NormalTok{, num\_retained\_lasso, }\StringTok{") — greater? "}\NormalTok{, num\_selected\_elastic\_net[i] }\SpecialCharTok{\textgreater{}}\NormalTok{ num\_retained\_lasso, }\StringTok{"}\SpecialCharTok{\textbackslash{}n\textbackslash{}n}\StringTok{"}\NormalTok{)}
\NormalTok{\}}
\end{Highlighting}
\end{Shaded}

\begin{verbatim}
 The Elastic Net (alpha= 0.25 ) selected  69  features.
 The Lasso (alpha=1) selected  25  features.
 The Selected Features of Elastic Net ( 69 ) compared to Lasso ( 25 ) — greater?  TRUE 

 The Elastic Net (alpha= 0.5 ) selected  55  features.
 The Lasso (alpha=1) selected  25  features.
 The Selected Features of Elastic Net ( 55 ) compared to Lasso ( 25 ) — greater?  TRUE 

 The Elastic Net (alpha= 0.75 ) selected  48  features.
 The Lasso (alpha=1) selected  25  features.
 The Selected Features of Elastic Net ( 48 ) compared to Lasso ( 25 ) — greater?  TRUE 
\end{verbatim}

\begin{Shaded}
\begin{Highlighting}[]
\CommentTok{\# From this data frame of all (α = alphas, λmin = lambda mins, min CV{-}MSE = minimum cross validation mean squared error, and num\_nonzero\_features = number of nonzero features extracted from the lambda min values),}
\CommentTok{\# it would be easy to witness that the alpha that yields the lowest CV{-}MSE is the alpha of 0.25 where the lowest CV{-}MSE value of 0.203687.}
\CommentTok{\# Also, it is easy to interpret that the number of selected features from the elastic net is always greater than the number of selected features from the lasso.}


\CommentTok{\# Problem 6. Fit PCR using the pls package (or manual SVD as shown in the chapter). Select the number of }
\CommentTok{\# components k{-}fold value(s) via cross{-}validation (use validationplot). Report the optimal k.}

\CommentTok{\# fitting PCR model with the y and x variable values and the k{-}fold value is 10}
\NormalTok{pcr\_fit }\OtherTok{\textless{}{-}} \FunctionTok{pcr}\NormalTok{(y }\SpecialCharTok{\textasciitilde{}}\NormalTok{ x, }\AttributeTok{scale =} \ConstantTok{TRUE}\NormalTok{, }\AttributeTok{validation =} \StringTok{"CV"}\NormalTok{, }\AttributeTok{segments =} \DecValTok{10}\NormalTok{)}

\CommentTok{\# Validation plot – choose optimal number of components}
\FunctionTok{validationplot}\NormalTok{(pcr\_fit, }\AttributeTok{val.type =} \StringTok{"MSE"}\NormalTok{, }
               \AttributeTok{main =} \StringTok{"PCR: k=10 fold {-} Cross{-}Validation Error vs Number of PC Components"}\NormalTok{,}
               \AttributeTok{xlab =} \StringTok{"Number of Principal Components (PC)"}\NormalTok{, }\AttributeTok{ylab =} \StringTok{"Cross{-}Validation Mean{-}Squared Error (MSE)"}\NormalTok{,}
               \AttributeTok{legendpos =} \StringTok{"topright"}\NormalTok{)}
\end{Highlighting}
\end{Shaded}

\pandocbounded{\includegraphics[keepaspectratio]{HW8_files/figure-pdf/unnamed-chunk-1-5.pdf}}

\begin{Shaded}
\begin{Highlighting}[]
\CommentTok{\# Identify optimal k}
\NormalTok{optimal\_ncomp }\OtherTok{\textless{}{-}} \FunctionTok{which.min}\NormalTok{(pcr\_fit}\SpecialCharTok{$}\NormalTok{validation}\SpecialCharTok{$}\NormalTok{PRESS)}
\FunctionTok{cat}\NormalTok{(}\StringTok{"The Optimal k (number of PC components) in the calculated Partial Least Square Regression Value (PLS):"}\NormalTok{, optimal\_ncomp, }\StringTok{"}\SpecialCharTok{\textbackslash{}n\textbackslash{}n}\StringTok{"}\NormalTok{)}
\end{Highlighting}
\end{Shaded}

\begin{verbatim}
The Optimal k (number of PC components) in the calculated Partial Least Square Regression Value (PLS): 22 
\end{verbatim}

\begin{Shaded}
\begin{Highlighting}[]
\CommentTok{\#  a. Back{-}transform the PCR fit to obtain the full 4,088{-}dimensional coefficient vector 𝛽pcr.}
\CommentTok{\# transforming }

\CommentTok{\# this performs a backdoor linear transformation on all x values with the coefficients altered by the most optimal k number of principal components value for the PCR fitted object}
\NormalTok{backdoor\_coefficient\_transformation\_on\_k\_optimal\_pcr }\OtherTok{\textless{}{-}} \FunctionTok{coef}\NormalTok{(pcr\_fit, }\AttributeTok{ncomp =}\NormalTok{ optimal\_ncomp)}
\NormalTok{dimensions }\OtherTok{\textless{}{-}} \FunctionTok{dim}\NormalTok{(backdoor\_coefficient\_transformation\_on\_k\_optimal\_pcr)}
\FunctionTok{cat}\NormalTok{(}\StringTok{"Checking all of the Dimensions in the PCR object to check if the transformations of the coefficients in the coefficient vector is correct for 4088: "}\NormalTok{, }\FunctionTok{paste}\NormalTok{(dimensions, }\AttributeTok{collapse =} \StringTok{" x "}\NormalTok{), }\StringTok{"}\SpecialCharTok{\textbackslash{}n\textbackslash{}n}\StringTok{"}\NormalTok{)}
\end{Highlighting}
\end{Shaded}

\begin{verbatim}
Checking all of the Dimensions in the PCR object to check if the transformations of the coefficients in the coefficient vector is correct for 4088:  4088 x 1 x 1 
\end{verbatim}

\begin{Shaded}
\begin{Highlighting}[]
\CommentTok{\#  b. Use 𝛽 pcr to generate a prediction for the first observation in the dataset (X[1, ]). Show your R code and the predicted value.}
\NormalTok{pred\_first\_observation\_of\_the\_pcr }\OtherTok{\textless{}{-}} \FunctionTok{predict}\NormalTok{(pcr\_fit, }
                      \AttributeTok{newdata =}\NormalTok{ x[}\DecValTok{1}\NormalTok{, , }\AttributeTok{drop =} \ConstantTok{FALSE}\NormalTok{], }
                      \AttributeTok{ncomp =}\NormalTok{ optimal\_ncomp)}

\CommentTok{\# predicted value for the first observation}
\FunctionTok{cat}\NormalTok{(}\StringTok{"Predicted value for the first observation (X[1,]):"}\NormalTok{, }
    \FunctionTok{round}\NormalTok{(}\FunctionTok{as.numeric}\NormalTok{(pred\_first\_observation\_of\_the\_pcr), }\DecValTok{6}\NormalTok{), }\StringTok{"}\SpecialCharTok{\textbackslash{}n\textbackslash{}n}\StringTok{"}\NormalTok{)}
\end{Highlighting}
\end{Shaded}

\begin{verbatim}
Predicted value for the first observation (X[1,]): -6.562936 
\end{verbatim}

\begin{Shaded}
\begin{Highlighting}[]
\CommentTok{\#  c. In one sentence, explain why you needed all 4,088 gene expression values to generate}
\CommentTok{\# this prediction even though you only fit k principal components.}

\CommentTok{\# For any prection to be accurate, we need to use all 4088 genes in linear combination with }
\CommentTok{\# the new principal components (uncorrelated representations of the original predictor variables x) }
\CommentTok{\# along with the coefficients in the coefficent vector to generate the accurate full{-}wieghted predictions of all }
\CommentTok{\# the 4088 genes in this data set.}

\CommentTok{\# Problem 7. Comparing across methods:}

  \CommentTok{\# a. Report the CV{-}MSE at the selected tuning parameter for each of the four methods (ridge, lasso, elastic net, PCR).}
\NormalTok{min\_cv\_mse\_ridge }\OtherTok{\textless{}{-}} \FunctionTok{min}\NormalTok{(cv\_10\_ridge}\SpecialCharTok{$}\NormalTok{cvm)}
\NormalTok{min\_cv\_mse\_lasso }\OtherTok{\textless{}{-}} \FunctionTok{min}\NormalTok{(cv\_lasso}\SpecialCharTok{$}\NormalTok{cvm)}
\NormalTok{min\_cv\_mse\_elastic\_net\_min }\OtherTok{\textless{}{-}} \FunctionTok{min}\NormalTok{(min\_cv\_mse\_elastic\_net)}
\NormalTok{cv\_mse\_pcr   }\OtherTok{\textless{}{-}} \FunctionTok{min}\NormalTok{(pcr\_fit}\SpecialCharTok{$}\NormalTok{validation}\SpecialCharTok{$}\NormalTok{PRESS) }\SpecialCharTok{/} \FunctionTok{length}\NormalTok{(y)}

  \CommentTok{\# b. Which method achieves the lowest CV{-}MSE on this dataset?}

\CommentTok{\# Compare the lowest CV{-}MSE values for all methods}
\NormalTok{lowest\_cv\_mse\_methods }\OtherTok{\textless{}{-}} \FunctionTok{c}\NormalTok{(}\StringTok{"Ridge Regression"}\NormalTok{, }\StringTok{"Lasso Regression"}\NormalTok{, }\StringTok{"Elastic Net Regression"}\NormalTok{, }\StringTok{"Partial Least Square Regression (PCR)"}\NormalTok{)}
\NormalTok{lowest\_cv\_mse\_values }\OtherTok{\textless{}{-}} \FunctionTok{c}\NormalTok{(min\_cv\_mse\_ridge, min\_cv\_mse\_lasso, min\_cv\_mse\_elastic\_net\_min, cv\_mse\_pcr)}

\CommentTok{\# Find the method with the lowest CV{-}MSE}
\NormalTok{lowest\_cv\_mse\_index }\OtherTok{\textless{}{-}} \FunctionTok{which.min}\NormalTok{(lowest\_cv\_mse\_values)}
\NormalTok{lowest\_cv\_mse\_method\_val }\OtherTok{\textless{}{-}}\NormalTok{ lowest\_cv\_mse\_methods[lowest\_cv\_mse\_index]}

\CommentTok{\# Print the results}
\FunctionTok{cat}\NormalTok{(}\StringTok{"The lowest CV{-}MSE achieved by the methods are:}\SpecialCharTok{\textbackslash{}n}\StringTok{"}\NormalTok{)}
\end{Highlighting}
\end{Shaded}

\begin{verbatim}
The lowest CV-MSE achieved by the methods are:
\end{verbatim}

\begin{Shaded}
\begin{Highlighting}[]
\ControlFlowTok{for}\NormalTok{ (i }\ControlFlowTok{in} \DecValTok{1}\SpecialCharTok{:}\FunctionTok{length}\NormalTok{(lowest\_cv\_mse\_methods)) \{}
  \FunctionTok{cat}\NormalTok{(lowest\_cv\_mse\_methods[i], }\StringTok{": "}\NormalTok{, }\FunctionTok{round}\NormalTok{(lowest\_cv\_mse\_values[i], }\DecValTok{6}\NormalTok{), }\StringTok{"}\SpecialCharTok{\textbackslash{}n}\StringTok{"}\NormalTok{)}
\NormalTok{\}}
\end{Highlighting}
\end{Shaded}

\begin{verbatim}
Ridge Regression :  0.269497 
Lasso Regression :  0.209623 
Elastic Net Regression :  0.203687 
Partial Least Square Regression (PCR) :  0.248206 
\end{verbatim}

\begin{Shaded}
\begin{Highlighting}[]
\CommentTok{\# Print the method with the lowest CV{-}MSE}
\FunctionTok{cat}\NormalTok{(}\StringTok{"}\SpecialCharTok{\textbackslash{}n}\StringTok{The method with the lowest CV{-}MSE is:"}\NormalTok{, lowest\_cv\_mse\_method\_val, }\StringTok{"}\SpecialCharTok{\textbackslash{}n\textbackslash{}n}\StringTok{"}\NormalTok{)}
\end{Highlighting}
\end{Shaded}

\begin{verbatim}

The method with the lowest CV-MSE is: Elastic Net Regression 
\end{verbatim}

\begin{Shaded}
\begin{Highlighting}[]
  \CommentTok{\# c. Based on the guidance in the chapter, would you expect ridge or lasso to work better here?}
  \CommentTok{\# Does the data agree with your expectation? Why or why not?}

  \CommentTok{\# Due to all of the details and guidance within this chapter 7., I would say that since this data set is }
  \CommentTok{\# a high{-}dimensional genetics dataset that would need to have some kind of correctlation, association and}
  \CommentTok{\# even collinearity within the features/variables in the dataset since the calculation of the CV{-}MSE values}
  \CommentTok{\# needs to be accurate considering that this measure is attempto to determine which method(s) actually predicts acccurately}
  \CommentTok{\# to the already observed values/data in the genetics dataset along with being able to predict any changes in the prediction model }
  \CommentTok{\# that would occur such as the inputting of new data and how these predictions reacts to changing values across the x axis such as }
  \CommentTok{\# how the lambda λ values in the ridge regression and lasso regression models. It is also important to note that these}
  \CommentTok{\# MSE values are also dependent on the lambda λ values that are used in the ridge regression and lasso regression models }
  \CommentTok{\# along with helping determine the bias in the y prediction values and how the variances that these MSE help define behave and/or are structured}
  \CommentTok{\# as the lambda λ values increase or decrease. In this case though, the data here really does meet my expectation in accurately predicting}
  \CommentTok{\# the predicted y values from how these MSE values from each model reveals in the most non{-}linear way and the non{-}consideration of intercept and zero values}
  \CommentTok{\# which model is more accurate or not based on its lower CV{-}MSE value. In conclusion after some comparisons between the ridge and lassoo prediction models}
  \CommentTok{\# along with the the MSE curve values calculated in this chapter (Ridge Regression :  0.2694971, Lasso Regression :  0.2096233, Elastic Net Regression :  0.2036869 )}
  \CommentTok{\# it becomes easy to conclude that the Lasso Regression (and possibly the Elastic Net Regression) is more accurate in predicting the y values in this genetics dataset.}
  \CommentTok{\# This conclusion is backed up by the MSE values listed above.}
\end{Highlighting}
\end{Shaded}

\bookmarksetup{startatroot}

\chapter{Summary}\label{summary}

In summary, this book has no content whatsoever.

\bookmarksetup{startatroot}

\chapter*{References}\label{references}
\addcontentsline{toc}{chapter}{References}

\markboth{References}{References}

\protect\phantomsection\label{refs}
\begin{CSLReferences}{1}{1}
\bibitem[\citeproctext]{ref-knuth84}
Knuth, Donald E. 1984. {``Literate Programming.''} \emph{Comput. J.}
(USA) 27 (2): 97--111. \url{https://doi.org/10.1093/comjnl/27.2.97}.

\end{CSLReferences}




\end{document}
